\pdfoutput=1
\documentclass[11pt]{amsart}
\usepackage[utf8]{inputenc}
\usepackage[T1]{fontenc}
\usepackage{lmodern}
\ifdefined\pdfgentounicode\input{glyphtounicode}\pdfgentounicode=1\fi
\usepackage{amsmath,amssymb,amsthm}
\usepackage[margin=2.8cm]{geometry}
\usepackage{enumitem}
\usepackage{booktabs}
\usepackage{needspace}
\usepackage{xcolor}
\usepackage[colorlinks=true,linkcolor=blue!50!black,citecolor=blue!50!black,urlcolor=blue!50!black]{hyperref}

\newtheorem{theorem}{Theorem}[section]
\newtheorem{proposition}[theorem]{Proposition}
\newtheorem{lemma}[theorem]{Lemma}
\newtheorem{corollary}[theorem]{Corollary}
\newtheorem*{theoremA}{Theorem A}
\newtheorem*{theoremB}{Theorem B}
\newtheorem*{theoremC}{Theorem C}
\newtheorem*{theoremD}{Theorem D}
\newtheorem*{theoremE}{Theorem E}
\newtheorem*{theoremF}{Theorem F}
\newtheorem*{theoremG}{Theorem G}
\theoremstyle{definition}
\newtheorem{definition}[theorem]{Definition}
\newtheorem{example}[theorem]{Example}
\theoremstyle{remark}
\newtheorem{remark}[theorem]{Remark}

\setlength{\emergencystretch}{2em}
\newcommand{\PP}{\mathbb{P}}
\newcommand{\QQ}{\mathbb{Q}}
\newcommand{\RR}{\mathbb{R}}
\newcommand{\ZZ}{\mathbb{Z}}
\newcommand{\OO}{\mathcal{O}}
\newcommand{\cE}{\mathcal{E}}
\newcommand{\cF}{\mathcal{F}}
\newcommand{\cI}{\mathcal{I}}
\newcommand{\cQ}{\mathcal{Q}}
\newcommand{\cU}{\mathcal{U}}
\newcommand{\cT}{\mathcal{T}}
\newcommand{\Tt}{\mathsf{T}}
\newcommand{\CC}{\mathcal{C}}
\newcommand{\BS}{\mathrm{BS}}
\DeclareMathOperator{\rk}{rk}
\DeclareMathOperator{\gr}{gr}
\DeclareMathOperator{\Ext}{Ext}
\DeclareMathOperator{\coker}{coker}

\title[Boij--Söderberg and Harder--Narasimhan polygons]{Boij--Söderberg and Harder--Narasimhan polygons of vector bundles on projective space}
\author{Cristian Anghel}
\address{``Simion Stoilow'' Institute of Mathematics of the Romanian Academy, Bucharest, Romania}
\email{cristian.anghel@imar.ro}
\author{Filip Chindea}
\address{``Simion Stoilow'' Institute of Mathematics of the Romanian Academy, Bucharest, Romania}
\email{fchindea@imar.ro}
\subjclass[2020]{Primary 14J60; Secondary 13D02, 14F06, 14F08}
\keywords{Boij--Söderberg theory, cohomology tables, supernatural bundles, Harder--Narasimhan filtration, Bridgeland stability, Ulrich bundles}

\begin{document}

\begin{abstract}
Eisenbud and Schreyer proved that the cohomology table of a vector bundle $F$ on $\PP^n$ is a unique positive combination of supernatural tables along a chain of root sequences. The slopes of the pieces, weighted by their ranks, form a concave polygon $\Psi_F$, the Boij--Söderberg polygon, with the same endpoints as the Harder--Narasimhan polygon $P_F$. We prove that $P_F\le\Psi_F$, in any characteristic, by testing subsheaves against sheaves with pure resolutions on thickened lines; the bound is optimal among bounds depending only on the ray of the table. The chain majorizes every other supernatural decomposition, and $\Psi$ is monotone under cancellations, hence upper semicontinuous in flat families. The comparison extends to coherent sheaves and, by linear projection, to sheaves of positive dimension. On $\PP^2$ the defect is a sum of four cohomological terms; we characterize equality and show that the inequality survives the large volume limit of Bridgeland stability. For a smooth projective variety $X$ of dimension $n\ge2$ in characteristic zero, with $H$ ample and base-point-free, the mean normalized squared spread $G(\cE)$ of the roots in the cohomology table of every vector bundle $\cE$ satisfies $G(\cE)\ge(K_X^2-2c_2(X))\cdot H^{n-2}/((n+1)H^n)$, a bound independent of $\cE$. Since $G(\cE)=1$ exactly for arithmetically Cohen--Macaulay bundles, when this bound exceeds $1$ it is a uniform form of their absence.
\end{abstract}

\hypersetup{pdftitle={Boij--Soderberg and Harder--Narasimhan polygons of vector bundles on projective space},pdfauthor={Cristian Anghel and Filip Chindea}}
\maketitle

\enlargethispage{2pt}
\setcounter{tocdepth}{1}
\tableofcontents

\section{Introduction}\label{sec:intro}

Let $K$ be an algebraically closed field of characteristic zero. The characteristic enters through Bogomolov's inequality in Sections~\ref{sec:plane} and~\ref{sec:polarized}, through Theorem~\ref{thm:lvl} and Remark~\ref{rem:lvl}, where we work over $\mathbb C$ as in the references on Bridgeland stability, and through results we quote: the comparison with \cite{ErSa16} in Remark~\ref{rem:main}(3) and the stability statements in Appendix~\ref{app:P3}. Theorem~A, Sections~\ref{sec:tables}--\ref{sec:monotone}, Theorems~\ref{thm:three} and~\ref{thm:HNmono}, Proposition~\ref{prop:fourdefect}, the results of Section~\ref{sec:equality} and Theorem~G hold over any algebraically closed field (see the table below), since the pure resolutions and the supernatural bundles of \cite[Thms.~5.1, 6.1]{ES09} exist over any field. For a coherent sheaf $\cF$ on $\PP^n$ let $\gamma(\cF)=(h^j(\cF(t)))_{0\le j\le n,\,t\in\ZZ}$ be its cohomology table. Eisenbud and Schreyer \cite{ES09} proved that the cohomology tables of vector bundles span, up to positive scalars, the cone generated by the tables $\gamma^z$ of supernatural bundles, and that every such table has a unique decomposition
\[
\gamma(F)=\sum_kq_k\,\gamma^{z^k},\qquad q_k>0,
\]
along a chain $z^1>z^2>\dots$ of root sequences; see \cite{ES10,ESi,Fl12} for the theory and its extension to coherent sheaves. A root sequence is a strictly decreasing sequence of integers $z=(z_1>\dots>z_n)$, and we normalize $\gamma^z$ to have Hilbert polynomial $\prod_j(t-z_j)$, i.e.\ rank $n!$. We attach to $z$ its \emph{center} $c(z)=\frac1n\sum_jz_j$ and its \emph{slope} $\mu(z)=-c(z)-\frac{n+1}2$; a supernatural bundle with root sequence $z$ has slope $\mu(z)$. Decreasing slopes correspond to increasing centers. To link the Eisenbud--Schreyer chain decomposition to slope stability, we introduce the following polygon.

\begin{definition}
The \emph{Boij--Söderberg polygon} $\Psi_F\colon[0,\rk F]\to\RR$ of a vector bundle $F$ on $\PP^n$ is the concave polygon whose edges have slopes $\mu(z^k)$ and horizontal lengths $n!\,q_k$, arranged in decreasing order of slope.
\end{definition}

Its endpoints are $(0,0)$ and $(\rk F,\deg F)$, the same as those of the Harder--Narasimhan polygon $P_F$, whose edges have the slopes of the Harder--Narasimhan factors and lengths their ranks. The two polygons record different kinds of information: $P_F$ is defined by subsheaves, $\Psi_F$ by cohomology. Our main result compares them.

\begin{theoremA}[Theorem~\ref{thm:main}]
Let $F$ be a vector bundle on $\PP^n$. Then $\deg F'\le\Psi_F(\rk F')$ for every subsheaf $F'\subset F$; equivalently, $P_F\le\Psi_F$. In particular, a vector bundle whose cohomology table is a multiple of a single supernatural table is semistable.
\end{theoremA}

Theorem~A gives an upper bound for the slopes of all subsheaves in terms of the table alone, and it yields $V_{HN}(F)\le V_c(F)$ for the variances of the Harder--Narasimhan slopes and of the Boij--Söderberg slopes (Corollary~\ref{cor:supernatural}). The semistability of supernatural bundles is not the main point: it also follows from known results, in any characteristic from the semistability of Ulrich bundles \cite[Thm.~2.9]{CH12}, applied to an Ulrich bundle on $\PP^n\times(\PP^1)^W$ built from $F$, and in characteristic zero from the work of Erman and Sam \cite{ErSa16} (Remark~\ref{rem:main}(3)). The content of Theorem~A is the comparison for arbitrary chains, and it is optimal among bounds that depend only on the ray of the table: every table has a multiple that is the table of a direct sum of supernatural bundles with $P=\Psi$. For a fixed table, integrality gives a sharper bound (Remark~\ref{rem:main}(1)). Moreover a saturated subsheaf of a supernatural bundle with the same slope, and its quotient, are again supernatural with the same roots (Proposition~\ref{prop:sameslope}), so supernatural bundles are Gieseker semistable.

The proof rests on a formula for the concave conjugate
\[
\Lambda_F(\lambda)=\sum_kn!\,q_k\big(\mu(z^k)-\lambda\big)_+ ,\qquad \Psi_F(m)=\min_\lambda\big(m\lambda+\Lambda_F(\lambda)\big).
\]
For a root sequence $z$ comparable with every element of the chain, $\Lambda_F(\mu(z))$ is the divided difference of truncated Euler characteristics
\[
\phi_z(\gamma(F))=(n-1)!\sum_{r=1}^n\frac{\sum_{j<r}(-1)^jh^j(F(z_r))}{\prod_{i\ne r}(z_r-z_i)} .
\]
By the theorem of Eisenbud, Fløystad, Weyman and Eisenbud, Schreyer on pure resolutions \cite{EFW11,ES09}, there is a sheaf $\Tt$, set-theoretically supported on a general line $L$ (in general on a thickening of $L$), whose minimal free resolution is pure, with twists $z_1,\dots,z_n$. For a saturated subsheaf $F'\subset F$ we may choose $L$ disjoint from the loci where $F'$ and $F/F'$ are not locally free; then $F'\otimes \Tt\subset F\otimes \Tt$, and
\[
\chi(F'\otimes \Tt)\ \le\ h^0(F'\otimes \Tt)\ \le\ h^0(F\otimes \Tt).
\]
Riemann--Roch turns the left-hand side into a positive multiple of $\deg F'-\mu(z)\rk F'$, and the hypercohomology spectral sequence of $F$ tensored with the resolution bounds the right-hand side by the same multiple of $\phi_z(\gamma(F))$. This spectral sequence argument is a variant of the one by which Eisenbud and Schreyer prove the nonnegativity of their pairing between Betti tables and cohomology tables \cite[\S3]{ES09}; here the sheaf $\Tt$ has one-dimensional support, and $h^0(F\otimes \Tt)$ replaces the vanishing of $F\otimes \Tt$. On $\PP^2$ the test sheaf can be taken to be $\OO_{(a-b)L}(a)$, with resolution $\OO(b)\to\OO(a)$, for a line $L$ and $a>b$, and the argument is elementary (Section~\ref{sec:plane}).

\medskip\noindent\textbf{Structural results on \texorpdfstring{$\PP^n$}{Pn}.}
Let $\CC$ be the set of nonnegative rational combinations of supernatural tables on $\PP^n$. Every $\gamma\in\CC$ has a unique chain decomposition, and $\Psi_\gamma$, $\Lambda_\gamma$ are defined as above. Applied to supernatural bundles, the inequality behind Theorem~A says that $\phi_z(\gamma^y)\ge n!\,(c(z)-c(y))_+$ for all root sequences $y,z$, with equality if $y$ and $z$ are comparable. So, for $\lambda\in\frac1n\ZZ$, $\Lambda_\gamma(\lambda)$ is the minimum of the linear functionals $\phi_z$ with $\mu(z)=\lambda$, attained at the root sequences comparable with the chain (Proposition~\ref{prop:minformula}). Two consequences follow.

\begin{theoremB}[Theorem~\ref{thm:concave}]
Let $\gamma\in\CC$, and let $\gamma=\sum_jq'_j\gamma^{y_j}$ be any decomposition with $q'_j\ge0$ and arbitrary root sequences $y_j$. Then the distribution of centers of the chain decomposition, with masses $q_k$, majorizes the distribution of the centers $c(y_j)$, with masses $q'_j$. Equivalently, $\Psi_\gamma$ is the largest slope polygon of any decomposition of $\gamma$; and for every convex function $f$ the map $\gamma\mapsto\sum_kq_kf(c(z^k))$ is concave on $\CC$.
\end{theoremB}

Let $e_{j,t}$ be the table with a single entry $1$ in row $j$ and column $t$. We call the tables $e_{j,t}+e_{j+1,t}$, $0\le j<n$, \emph{cancellation vectors}. If $0\to F'\to F\to F''\to0$ is an exact sequence, the long exact cohomology sequences give
\[
\gamma(F')+\gamma(F'')-\gamma(F)=\sum_{t}\sum_{j=0}^{n-1}\rk\big(\partial^j_t\big)\,\big(e_{j,t}+e_{j+1,t}\big),
\]
where $\partial^j_t\colon H^j(F''(t))\to H^{j+1}(F'(t))$ are the connecting homomorphisms. The value of $\phi_z$ on $e_{r-1,t}+e_{r,t}$ is $(n-1)!/\prod_{j\ne r}|t-z_j|$ if $t=z_r$, and $0$ otherwise; in particular it is nonnegative.

\begin{theoremC}[Theorem~\ref{thm:monotone}, Corollary~\ref{cor:monoext}]
Let $\gamma,\gamma'\in\CC$, and suppose that $\gamma'-\gamma$ is a nonnegative combination of cancellation vectors. Then $\Psi_\gamma\le\Psi_{\gamma'}$. Consequently:
\begin{enumerate}[label=\textup{(\roman*)}]
\item $\Psi_F\le\Psi_{F'\oplus F''}$ for every exact sequence $0\to F'\to F\to F''\to0$ of vector bundles on $\PP^n$; in particular $\Psi_F\le\Psi_{\gr F}$ for every filtration with locally free subquotients;
\item if $(F_s)_{s\in S}$ is a flat family of vector bundles on $\PP^n$ over a variety $S$, then every $s_0\in S$ has a neighborhood $U$ with $\Psi_{F_s}\le\Psi_{F_{s_0}}$ for all $s\in U$; if $S$ is irreducible, $\Psi_{F_s}\le\Psi_{F_{s_0}}$ for general $s$ and every $s_0$.
\end{enumerate}
\end{theoremC}

On $\PP^2$ the hypothesis only says that $\gamma\le\gamma'$ entrywise and that both have the same Hilbert polynomial. For every $n\ge3$ this weaker hypothesis does not suffice; on $\PP^3$ the tables of $\OO(-4)^{6}\oplus\OO(-2)^{6}$ and $\OO(-5)\oplus\OO(-4)^{2}\oplus\OO(-3)^{6}\oplus\OO(-2)^{2}\oplus\OO(-1)$ are entrywise comparable, with the same Hilbert polynomial, but their polygons are not (Example~\ref{ex:weak}). Part (ii) is the analogue for $\Psi$ of Shatz's semicontinuity of the Harder--Narasimhan polygon \cite{Sh77}. Theorem~C gives a structural constraint on two problems of Erman and Sam \cite[Problem~4.3 and \S5]{ErS}, on the cohomology tables of the extensions of two given bundles and on the jumps of cohomology tables in flat families: from an extension to the direct sum, and from a general to a special fiber, the center distribution of the chain only moves by mean-preserving spreads. The same computation bounds how much cancellation can lower $\Lambda$ (Proposition~\ref{prop:defectineq}). With $\Lambda$ and $\Psi$ defined by the infimum formula, Theorem~A and Theorem~C(i) hold for all coherent sheaves on $\PP^n$, and Theorem~D for all coherent sheaves on $\PP^2$ (Section~\ref{sec:coherent}); through a general linear projection they extend to sheaves of lower dimension (Theorem~G below). On $\PP^2$ it follows that a coherent sheaf whose decomposition involves only tables of vector bundles is torsion-free, which answers in this case a question of Eisenbud and Schreyer \cite{ES10} (Corollary~\ref{cor:ESquestion}).

\medskip\noindent\textbf{The plane.}
On $\PP^2$ a root sequence is a pair $z=(z_1>z_2)$, and $\phi_{(a,b)}(\gamma)=\big(\gamma_{0,a}-\gamma_{0,b}+\gamma_{1,b}\big)/(a-b)$ involves three entries of the table.

\begin{theoremD}[Theorem~\ref{thm:three}, Proposition~\ref{prop:fourdefect}]
Let $F$ be a vector bundle on $\PP^2$, $\lambda\in\frac12\ZZ$, and $a>b$ integers with $a+b=-2\lambda-3$. Then
\[
(a-b)\,\Lambda_F(\lambda)\ \le\ h^0(F(a))-h^0(F(b))+h^1(F(b)),
\]
with equality whenever $(a,b)$ is comparable with every element of the chain of $F$; such pairs exist for every $\lambda$. For such a pair and a subsheaf $F'\subset F$ of rank $n'$ and slope $\mu'$ with torsion-free quotient $F''$,
\begin{multline*}
(a-b)\big(\Lambda_F(\lambda)-n'(\mu'-\lambda)\big)\\=\big(h^2_b(F')-h^2_a(F')\big)+\big(h^0_a(F'')-h^0_b(F'')\big)+\big(h^1_a(F')-\rk\partial^0_a\big)+\big(h^1_b(F'')-\rk\partial^1_b\big),
\end{multline*}
where $h^j_t(G)=h^j(G(t))$, and all four terms are nonnegative.
\end{theoremD}

The nonnegativity of the first two terms is the monotonicity of $h^0$ and $h^2$ under multiplication by a linear form, so Theorem~D gives an elementary proof of Theorem~A on $\PP^2$, together with a description of its equality cases, which we use in Section~\ref{sec:lvl}. With the gap moment $G(F)=\sum_kp_k(z^k_1-z^k_2)^2$, $p_k=2q_k/\rk F$, Theorem~A gives $G(F)\ge1+4\sum_iw_i\,\delta(\cQ_i)$, where $\cQ_i$ are the Harder--Narasimhan factors, with weights $w_i=\rk\cQ_i/\rk F$ and normalized discriminants $\delta(\cQ_i)\ge0$ (Corollary~\ref{cor:Q}).

On $\PP^2$ the argument uses only the tables of $F'$ and $F''$ and the ranks of the connecting homomorphisms, so it gives necessary conditions on the tables of extensions of two given bundles (Corollary~\ref{cor:conicQ}). On $\PP^3$ no argument of this kind can prove Theorem~A: there are tables $\gamma_1,\gamma_2$ of direct sums of supernatural bundles and cancellations compatible with the source--target bounds for the ranks of the connecting homomorphisms, for which the analogous conclusion fails (Example~\ref{ex:conicP3}). The test sheaves of Theorem~A see the injection $F'\otimes \Tt\to F\otimes \Tt$, which the tables do not record, and they exclude these tables. On $\PP^2$ a stronger comparison holds: the Boij--Söderberg polygons of the steps of the Harder--Narasimhan filtration increase, $\Psi_{F_j}\le\Psi_{F_k}$ on $[0,\rk F_j]$ for $j\le k$ (Theorem~\ref{thm:HNmono}); and $P_F=\Psi_F$ holds exactly when the Harder--Narasimhan factors are supernatural with pairwise comparable roots (Corollary~\ref{cor:equalityP2}).

\medskip\noindent\textbf{Bridgeland stability.}
On $\PP^2$ the Harder--Narasimhan polygon is the large volume limit of the Harder--Narasimhan polygons of the Bridgeland stability conditions $\sigma_{\alpha,\beta}$ \cite{Br08,ABCH13}. The chain decomposition defines a polygon $\Psi^\sigma_F$ in the same plane: its edges are the central charges $q_kZ_{\alpha,\beta}(E_{z^k})$ of the pieces, in order of decreasing phase.

\begin{theoremE}[Theorem~\ref{thm:lvl}]
Let $F$ be a vector bundle on $\PP^2$ over $\mathbb C$ and $\beta<\mu_{\min}(F)$ rational. For $\alpha\gg0$ the region bounded by the $\sigma_{\alpha,\beta}$-Harder--Narasimhan polygon of $F$ is contained in the region bounded by $\Psi^\sigma_F$.
\end{theoremE}

The leading order is Theorem~A. The next order is controlled by the equality cases of Theorem~D: at a tight vertex, the pieces above the vertex have the Chern character of the Harder--Narasimhan step (Lemma~\ref{lem:Ktight}); and on a contact edge the Harder--Narasimhan factor is a supernatural bundle (Proposition~\ref{prop:contactedge}), hence Gieseker semistable (Proposition~\ref{prop:sameslope}), so only the vertices of the Gieseker polygon need to be compared. For small $\alpha$ the supernatural pieces cross walls of their own, and the containment has to be formulated with the Harder--Narasimhan polygon of a supernatural model $\bigoplus_kE_{z^k}^{\oplus mq_k}$; whether it then holds for all $\sigma_{\alpha,\beta}$ with $\beta$ below the Boij--Söderberg slopes of $F$ is open. It does for $F=\OO\oplus E$ with $E$ a stable bundle with $c_1=0$, $c_2=2$ (Remark~\ref{rem:lvl}).

\medskip\noindent\textbf{Application to polarized varieties.}
Our motivation comes from \cite{An1,An2}. Let $X$ be a smooth projective variety of dimension $n\ge2$, $H$ an ample and base-point-free divisor, $d=H^n$, and $\sigma=(K_X^2-2c_2(X))\cdot H^{n-2}/(n+1)$, a number unrelated to the stability conditions $\sigma_{\alpha,\beta}$ above; for a surface, $\sigma=K_X^2-8\chi(\OO_X)$, the topological signature over $\mathbb C$. A general projection $\pi\colon X\to\PP^n$ is finite and flat, and $H^i(X,\cE(tH))=H^i(\PP^n,\pi_*\cE(t))$. In \cite{An2} it is shown that if $d<\sigma$ then $X$ carries no arithmetically Cohen--Macaulay bundle, i.e.\ no $\cE$ with $\pi_*\cE$ split. The \emph{gap moment} $G(\cE)=G(\pi_*\cE)$, the mean of a normalized spread of the root sequences of the chain of $\pi_*\cE$, satisfies $G\ge1$, with equality iff $\cE$ is aCM. Like the variance $V_c(\cE)$ of the centers, it depends only on the table $(h^i(\cE(tH)))_{i,t}=\gamma(\pi_*\cE)$, not on $\pi$.

\begin{theoremF}[Proposition~\ref{prop:defect}, Corollary~\ref{cor:cases}]
For every vector bundle $\cE$ of rank $r$ on $X$, with $\mu_H$-Harder--Narasimhan factors $\cQ_i$,
\[
\frac{n+1}{12}\big(d\,G(\cE)-\sigma\big)=d\big(V_c(\cE)-V_\beta(\cE)\big)+\sum_iw_i\big(\delta(\cQ_i)-\alpha(\cQ_i)^2\big),
\]
where $V_c$ is the variance of the centers of the chain, $V_\beta$ that of the normalized Harder--Narasimhan slopes, $\delta\ge0$ is a normalized Bogomolov discriminant and $\alpha^2\le0$ by the Hodge index theorem. The terms $d(V_c-V_\beta)$, $\sum_iw_i\delta(\cQ_i)$ and $-\sum_iw_i\alpha(\cQ_i)^2$ are nonnegative, the first by Theorem~A. Consequently
\[
G(\cE)\ \ge\ \frac\sigma d\qquad\text{and}\qquad\sum_{0<i<n}\sum_th^i(\cE(tH))\ \ge\ C_n\,r\,(\sigma-d)_+ .
\]
\end{theoremF}

So the absence of aCM bundles in \cite{An2}, which is the case $G=1$, comes with a uniform lower bound for the gap moment of every cohomology table of a vector bundle on $(X,H)$. The underlying Riemann--Roch identity on $X$ is \cite[Prop.~10.6, (10.3)]{An3} (for aCM bundles on surfaces, \cite[Prop.~2.6]{An2}); what is new here is its expression through the table, by $G$ and $V_c$, and the comparison $V_\beta\le V_c$ for every bundle, which is Theorem~A (for split push-forwards it is the majorization of \cite[Lemma~10.8]{An3}). Here $C_n=n\,(\lfloor n/2\rfloor-1)!\,(\lceil n/2\rceil-1)!/(12\,(n-2)!)$, which is optimal in the intermediate-cohomology estimate of Lemma~\ref{lem:gapPn}(iv), and $C_2=\frac16$. For the Hesse surfaces $Y_\nu$, $\nu\ge3$, of \cite{An2}, $\sigma/d\ge16/7$. For surfaces the integrality of the gaps gives the stronger bound $\sum_th^1(\cE(tH))\ge rd\,L(\kappa)$, $\kappa=\max(1,\sigma/d)$, where $L$ interpolates $k^2\mapsto\frac{k^3-k}{12}$ linearly.

The structure of $\Psi$ also bears on a question of Eisenbud and Schreyer (Corollary~\ref{cor:ES}): if all cohomology of $\cE$ vanishes at $n$ twists $a_1<\dots<a_n$, then $\pi_*\cE$ has a supernatural table with roots $a_j$, $\cE$ is $\mu_H$-semistable, and $\frac{n+1}{12}(d\,\rho-\sigma)=\delta(\cE)-\alpha(\cE)^2\ge0$, where $\rho\ge1$ is the normalized spread of the $a_j$, equal to $e^2$ for an arithmetic progression of difference $e$. In that case a twist of $\cE$ satisfies the cohomological condition for Ulrich bundles on $(X,eH)$, so when $eH$ is very ample the semistability and the inequality also follow from \cite{CH12} and from the inequality of Lopez for Ulrich bundles \cite[Lem.~3.1]{Lo26}; what is new is the case of root sequences not in arithmetic progression, which occurs only for $n\ge3$ and does not reduce by a twist to the Ulrich condition for a multiple of $H$, and the interpretation as the condition for a supernatural ray to be realized. On the Hesse surfaces $Y_\nu$ of \cite{An2} with $\nu\ge4$ no vector bundle $\cE$ has $H^i(\cE(-H))=H^i(\cE(H))=0$ for all $i$.

\medskip\noindent\textbf{Sheaves of lower dimension.}
For a coherent sheaf $\cF$ of dimension $d$ on $\PP^n$ write $\chi(\cF(t))=\sum_i(-1)^ie_i(\cF)\binom{t+d-i}{d-i}$. The multiplicity $e_0$ plays the role of the rank and $-e_1$ that of the degree; for pure sheaves we use the normalized slope $\mu_d=-e_1/e_0=\hat\mu-\frac{d+1}2$, where $\hat\mu$ is the slope of \cite[\S1.6]{HL10}. The coefficients of every subsheaf are computed in the same binomial basis of dimension $d$. A general linear projection to $\PP^d$ is finite on the support of $\cF$, and it preserves the cohomology table, the numbers $e_0$ and $e_1$, and inclusions of subsheaves (Lemma~\ref{lem:projection}).

\begin{theoremG}[Theorem~\ref{thm:lowerdim}]
Let $\cF$ be a coherent sheaf of dimension $d\ge1$ on $\PP^n$, and let $\Psi_\cF$ be defined from its table by the infimum formula, with root sequences of length $d$. Then $-e_1(\cF')\le\Psi_\cF(e_0(\cF'))$ for every subsheaf $\cF'\subset\cF$; for $\cF$ pure, the Harder--Narasimhan polygon of $\cF$ with respect to $\mu_d=-e_1/e_0$ lies below $\Psi_\cF$. If $\cF$ is pure of dimension $d$ and locally Cohen--Macaulay, for instance a vector bundle on a Cohen--Macaulay subvariety, then $\Psi_\cF$ is the Boij--Söderberg polygon of a vector bundle on $\PP^d$.
\end{theoremG}

For a graded module $M$ of dimension $d+1$ over $K[x_0,\dots,x_n]$, Theorem~G bounds $e_0$ and $e_1$ of every submodule of the same dimension in terms of the Hilbert functions of $M$ and of its local cohomology modules (Remark~\ref{rem:lowerdim}(2)). The multiplicity bounds of Herzog, Huneke and Srinivasan \cite{HS98}, proved in \cite{ES09,BS12}, concern $e_0(M)$ alone; their naive analogue for submodules, on the side of Betti tables, fails (Remark~\ref{rem:lowerdim}(3)).

\medskip\noindent\textbf{Organization.}
Section~\ref{sec:tables} collects notation on $\PP^n$. Section~\ref{sec:main} proves Theorem~A. Section~\ref{sec:monotone} proves Theorems~B and~C and the decomposition criterion, and extends Theorems~A, C(i) and~D to coherent sheaves (Section~\ref{sec:coherent}). Section~\ref{sec:plane} specializes to $\PP^2$: gap moment, extreme slopes, Theorem~D, the tables of extensions, monotonicity along the Harder--Narasimhan filtration, and the example on $\PP^3$. Section~\ref{sec:lvl} characterizes the equality $P_F=\Psi_F$ on $\PP^2$ and proves Theorem~E, Section~\ref{sec:polarized} contains the application to polarized varieties, and Section~\ref{sec:lowerdim} proves Theorem~G. Appendix~\ref{app:P3} describes, by Beilinson's theorem, the vector bundles on $\PP^3$ whose table is a multiple of $\gamma^z$ with $z_1-z_3\le4$; by Theorem~A they are all semistable, and for $z_1-z_3=4$ they are, up to a twist, instanton bundles, Steiner bundles or their duals (Corollary~\ref{cor:appP3}).

\begin{table}[ht]\centering\small
\textsc{Ground fields of the main results}\par\smallskip
\begin{tabular}{@{}lll@{}}
\toprule
Result & Section & Ground field\\
\midrule
Theorem~A, Propositions~\ref{prop:sameslope}, \ref{prop:marginPn} & \ref{sec:main} & any\\
Theorems~B, C & \ref{sec:monotone} & any\\
Propositions~\ref{prop:cohA}--\ref{prop:cohD}, Corollary~\ref{cor:ESquestion} & \ref{sec:coherent} & any\\
Theorem~D, Theorem~\ref{thm:HNmono} & \ref{sec:plane} & any\\
Corollary~\ref{cor:Q} (Bogomolov) & \ref{sec:plane} & characteristic $0$\\
Lemmas~\ref{lem:contact}--\ref{lem:gapdisc}, Proposition~\ref{prop:contactedge}, Corollary~\ref{cor:equalityP2} & \ref{sec:equality} & any\\
Lemma~\ref{lem:twist} & \ref{sec:lvl} & any\\
Theorem~E (rational $\beta$), Lemma~\ref{lem:support}, Remark~\ref{rem:lvl} & \ref{sec:lvl} & $\mathbb C$\\
Lemmas~\ref{lem:gapPn}, \ref{lem:push}, Proposition~\ref{prop:defect} & \ref{sec:polarized} & any\\
Theorem~F, Corollary~\ref{cor:ES} (Bogomolov) & \ref{sec:polarized} & characteristic $0$\\
Theorem~G & \ref{sec:lowerdim} & any\\
Stability of general bundles in Appendix~\ref{app:P3} & \ref{app:P3} & characteristic $0$\\
\bottomrule
\end{tabular}
\end{table}

\section{Cohomology tables on \texorpdfstring{$\PP^n$}{Pn}}\label{sec:tables}

\subsection{Supernatural tables}
A \emph{root sequence} is a strictly decreasing sequence of integers $z=(z_1>\dots>z_n)$, with center $c(z)=\frac1n\sum_jz_j\in\frac1n\ZZ$ and slope $\mu(z)=-c(z)-\frac{n+1}2\in\frac1n\ZZ$; note that $\frac{n+1}2=\frac1n\cdot\frac{n(n+1)}2\in\frac1n\ZZ$. Root sequences are ordered termwise, and a \emph{chain} is a totally ordered set. Two comparable root sequences with the same center are equal. The normalized supernatural table $\gamma^z$ is the cohomology table of a supernatural bundle with root sequence $z$ \cite[Thm.~0.4, \S6]{ES09}, scaled so that its Hilbert polynomial is $\prod_j(t-z_j)$:
\[
\gamma^z_{j,t}=\begin{cases}\big|\prod_{i=1}^n(t-z_i)\big| & \text{if } j=\#\{i:z_i>t\},\\ 0 & \text{otherwise.}\end{cases}
\]
So each column has at most one nonzero entry, the column $t$ vanishes iff $t$ is a root, and $\sum_j(-1)^j\gamma^z_{j,t}=\prod_i(t-z_i)$.

\begin{theorem}[Eisenbud--Schreyer {\cite[Thm.~0.5]{ES09}}]\label{thm:ES}
For every vector bundle $F$ on $\PP^n$ there are a unique chain $z^1>\dots>z^m$ and unique $q_k\in\QQ_{>0}$ with $\gamma(F)=\sum_kq_k\gamma^{z^k}$.
\end{theorem}

Since supernatural bundles exist for every root sequence \cite[Thm.~0.4]{ES09}, a positive integer multiple of every nonnegative rational combination of supernatural tables is the table of a direct sum of supernatural bundles. So Theorem~\ref{thm:ES} applies to all such combinations. We write $\CC$ for the set of nonnegative rational combinations of the $\gamma^z$; every $\gamma\in\CC$ has a unique chain decomposition with positive rational coefficients. We compare tables entrywise. For $\gamma=0$ the chain is empty, $\Psi_0$ is the point $(0,0)$ and $\Lambda_0=0$; statements about the first and last elements of a chain assume $\gamma\ne0$. Throughout, (semi)stable means slope (Mumford--Takemoto) (semi)stable with respect to $\OO(1)$; $P_F$, $\mu_{\max}(F)$ and $\mu_{\min}(F)$ refer to the corresponding Harder--Narasimhan filtration.

\subsection{Rank, degree and the polygon}
For $\gamma=\sum_kq_k\gamma^{z^k}\in\CC$ put
\[
\rk\gamma=n!\sum_kq_k,\qquad \deg\gamma=\sum_kn!\,q_k\,\mu(z^k).
\]
Both are linear in $\gamma$, because they are read off from the two leading coefficients of the Hilbert polynomial $\sum_kq_k\prod_j(t-z^k_j)$:
\[
\sum_kq_k\prod_j(t-z^k_j)=\frac{\rk\gamma}{n!}\,t^n+\Big(\deg\gamma+\frac{n+1}2\rk\gamma\Big)\frac{t^{n-1}}{(n-1)!}+\dots
\]
By Riemann--Roch, $\chi(\cF(t))=\frac{\rk\cF}{n!}t^n+\big(\deg\cF+\frac{n+1}2\rk\cF\big)\frac{t^{n-1}}{(n-1)!}+\dots$ for every coherent sheaf $\cF$, where $\deg\cF=c_1(\cF)\cdot h^{n-1}$ and $\rk\cF$ may be zero, so $\rk\gamma(F)=\rk F$ and $\deg\gamma(F)=\deg F$ for a vector bundle $F$. The \emph{Boij--Söderberg slopes} of $\gamma$ are the $\mu(z^k)$, with ranks $n!\,q_k$, and
\[
\Psi_\gamma(m)=\max\Big\{\sum_kx_k\,\mu(z^k)\ :\ 0\le x_k\le n!\,q_k,\ \sum_kx_k=m\Big\},\qquad 0\le m\le\rk\gamma .
\]
We also use the concave conjugate in the form
\[
\Lambda_\gamma(\lambda)=\sum_kn!\,q_k\,\big(\mu(z^k)-\lambda\big)_+=\sum_kn!\,q_k\,\big(c_\lambda-c(z^k)\big)_+,\qquad c_\lambda=-\lambda-\tfrac{n+1}2 ,
\]
so that $c(z)=c_\lambda$ iff $\mu(z)=\lambda$. It is convex and piecewise linear in $\lambda$, with breakpoints at the Boij--Söderberg slopes, which lie in $\frac1n\ZZ$. For a polygon $P$ with slopes $\mu_i$ and lengths $n_i$ put $\Lambda_P(\lambda)=\sum_in_i(\mu_i-\lambda)_+$.

\begin{lemma}\label{lem:legendre}
\begin{enumerate}[label=\textup{(\roman*)}]
\item Let $\gamma\ne0$. Then $\Psi_\gamma(m)=\min_\lambda\big(m\lambda+\Lambda_\gamma(\lambda)\big)$ for $0\le m\le\rk\gamma$, and the minimum is attained at a Boij--Söderberg slope of $\gamma$; in particular at a point of $\frac1n\ZZ$.
\item For two concave polygons $P,\Psi$ on $[0,N]$ with the same endpoints, $P\le\Psi$ iff $\Lambda_P\le\Lambda_\Psi$ on $\RR$, iff the distribution of slopes of $P$ is majorized by that of $\Psi$. Since $\Psi$ is concave and $P$ is linear between its vertices, $P\le\Psi$ iff $P(m)\le\Psi(m)$ at the vertices $m$ of $P$.
\item Let $\gamma=\gamma(F)$ and let $F_1\subset F$ have rank $n_1>0$ and slope $\mu_1$. Then $\Psi_F(n_1)\ge n_1\mu_1$ iff $\Lambda_F(\lambda)\ge n_1(\mu_1-\lambda)$ for all $\lambda\in\frac1n\ZZ$. If $F/F_1$ has rank $n_2>0$ and slope $\mu_2<\mu_1$, it suffices to check $\lambda\in(\mu_2,\mu_1)\cap\frac1n\ZZ$.
\end{enumerate}
\end{lemma}

\begin{proof}
(i) is the duality of the fractional knapsack problem: $m\lambda+\Lambda_\gamma(\lambda)$ is convex and piecewise linear in $\lambda$, with slope $m-\sum_{\mu(z^k)>\lambda}n!\,q_k$, which is negative for $\lambda$ below the smallest and positive above the largest Boij--Söderberg slope when $0<m<\rk\gamma$; so its minimum is at a breakpoint, and the endpoints $m=0,\rk\gamma$ are checked directly. Since $\Lambda_P(\lambda)=\max_m(P(m)-\lambda m)$ for a concave polygon with $P(0)=0$, the minimum equals $\Psi_\gamma(m)$. (ii) The first equivalence is the comparison of concave conjugates; the second is the theorem of Hardy, Littlewood and P\'olya \cite[Thm.~108]{HLP}, since both distributions have the same total mass and mean (for rational masses, repeat the points; the general case follows by continuity). The last assertion holds because $\Psi$ is concave. (iii) The first assertion follows from (i). For $\lambda\ge\mu_1$ the right-hand side is $\le0$. For $\lambda\le\mu_2$, convexity of $x\mapsto(x-\lambda)_+$ gives $\Lambda_F(\lambda)\ge\deg F-\rk F\cdot\lambda=n_1(\mu_1-\lambda)+n_2(\mu_2-\lambda)\ge n_1(\mu_1-\lambda)$.
\end{proof}

\subsection{Duality}
By Serre duality $h^j(F^\vee(t))=h^{n-j}(F(-t-n-1))$. On tables this is the involution $\gamma\mapsto\gamma^\vee$, $\gamma^\vee_{j,t}=\gamma_{n-j,-t-n-1}$, and $(\gamma^z)^\vee=\gamma^{z^\vee}$ with $z^\vee=(-n-1-z_n,\dots,-n-1-z_1)$. The map $z\mapsto z^\vee$ reverses the order and $\mu(z^\vee)=-\mu(z)$. So the chain of $F^\vee$ is $\{(z^k)^\vee\}$ with the same coefficients, its Boij--Söderberg slopes are the negatives of those of $F$, and
\begin{equation}\label{eq:dualpsi}
\Psi_{F^\vee}(N-m)=\Psi_F(m)-\deg F,\qquad N=\rk F .
\end{equation}
The Harder--Narasimhan polygon satisfies the same relation.

\subsection{The ends of the chain}
The first and last elements of the chain are read off from the supports of the cohomology rows.

\begin{lemma}[ends of the chain]\label{lem:ends}
Let $\gamma=\sum_lq'_l\gamma^{y_l}$ with all $q'_l>0$, for arbitrary root sequences $y_l$. Then the termwise maximum $y^{\max}$ and the termwise minimum $y^{\min}$ of the $y_l$ are given by
\[
y^{\max}_i=1-i+\max\{t+j:\ j\ge i,\ \gamma_{j,t}\ne0\},\qquad y^{\min}_i=\min\{t+j:\ j<i,\ \gamma_{j,t}\ne0\}-i .
\]
In particular they depend only on $\gamma$, and they are the first and the last element of the chain of $\gamma$.
\end{lemma}

\begin{proof}
All entries are nonnegative, so $\gamma_{j,t}\ne0$ iff $\gamma^{y_l}_{j,t}\ne0$ for some $l$; it suffices to prove the formulas for a single root sequence $y$. If $\gamma^y_{j,t}\ne0$ with $j\ge i$, then exactly $j$ roots exceed $t$, so $y_j>t$ and $y_i\ge y_j+(j-i)\ge t+1+j-i$. Conversely, let $j\ge i$ be maximal with $y_j=y_i-(j-i)$, and $t=y_j-1$. Then $t$ is not a root, because $y_{j+1}<y_j-1$ by maximality (or $j=n$), and exactly $y_1,\dots,y_j$ exceed $t$; so $\gamma^y_{j,t}\ne0$, and $1-i+t+j=y_i$. Symmetrically, if $\gamma^y_{j,t}\ne0$ with $j<i$, then $y_{j+1}<t$ and $y_i\le y_{j+1}-(i-j-1)\le t+j-i$; and if $j+1\le i$ is minimal with $y_{j+1}=y_i+(i-j-1)$, then $t=y_{j+1}+1$ is not a root, exactly $j$ roots exceed it, and $t+j-i=y_i$. Applied to the chain, whose elements are totally ordered, this identifies $y^{\max}$ and $y^{\min}$ with its first and last elements.
\end{proof}

\section{The comparison theorem}\label{sec:main}

\subsection{Divided differences}
For a root sequence $z$ put $D_r(z)=\prod_{i\ne r}(z_r-z_i)$ and $P_r(z)=|D_r(z)|$. Exactly $r-1$ of the $z_i$ exceed $z_r$, so $D_r(z)=(-1)^{r-1}P_r(z)$. For a table $\gamma$ put
\begin{equation}\label{eq:phi}
\phi_z(\gamma)=(n-1)!\sum_{r=1}^n\frac{1}{D_r(z)}\sum_{j<r}(-1)^j\gamma_{j,z_r} .
\end{equation}
This is a linear functional, which involves the entries of the rows $0,\dots,r-1$ in the column $z_r$. For a polynomial $f$, $\sum_rf(z_r)/D_r(z)$ is the divided difference $f[z_1,\dots,z_n]$; if $f$ has degree at most $n$, $f(t)=a_nt^n+a_{n-1}t^{n-1}+\dots$, then
\begin{equation}\label{eq:divdiff}
\sum_{r=1}^n\frac{f(z_r)}{D_r(z)}=a_n\sum_iz_i+a_{n-1}.
\end{equation}
For $n=2$ and $z=(a,b)$, $\phi_z(\gamma)=\big(\gamma_{0,a}-\gamma_{0,b}+\gamma_{1,b}\big)/(a-b)$. For $n=1$, $\phi_z(\gamma)=\gamma_{0,z_1}$.

\begin{lemma}\label{lem:comparable}
If $y$ and $z$ are comparable root sequences, then $\phi_z(\gamma^y)=n!\,\big(c(z)-c(y)\big)_+$.
\end{lemma}

\begin{proof}
If $y\le z$, then for $i\ge r$ we have $y_i\le z_i\le z_r$, so at most $r-1$ roots of $y$ exceed $z_r$, and the nonzero entry of the column $z_r$ of $\gamma^y$, if any, lies in a row $j<r$. Hence $\sum_{j<r}(-1)^j\gamma^y_{j,z_r}=\prod_i(z_r-y_i)$, also when $z_r$ is a root of $y$. By \eqref{eq:divdiff} applied to the monic polynomial $\prod_i(t-y_i)$, $\phi_z(\gamma^y)=(n-1)!\big(\sum_iz_i-\sum_iy_i\big)=n!\,(c(z)-c(y))\ge0$. If $y\ge z$, then $y_i\ge z_i\ge z_r$ for $i\le r$, so either $z_r$ is a root of $y$ or at least $r$ roots of $y$ exceed it; in both cases $\gamma^y_{j,z_r}=0$ for $j<r$, and $\phi_z(\gamma^y)=0=n!\,(c(z)-c(y))_+$.
\end{proof}

For root sequences $x\le y$ let $[x,y]=\{z:x\le z\le y\}$, and write $\epsilon_1,\dots,\epsilon_n$ for the unit vectors of $\ZZ^n$.

\begin{lemma}\label{lem:chains}
Let $x\le y$ be root sequences and $N=\sum_i(y_i-x_i)$. Lowering first the coordinate $n$ from $y_n$ to $x_n$, then the coordinate $n-1$, and so on, one unit at a time, passes only through root sequences. So there is a chain $y=z^0>z^1>\dots>z^N=x$ in $[x,y]$ with $z^k-z^{k+1}$ a unit vector and $c(z^k)=c(y)-\frac kn$.
\end{lemma}

\begin{proof}
Suppose the coordinates $n,\dots,l+1$ have been lowered to $x$ and the coordinate $l$ is being lowered from a value $z_l>x_l$. Then $z_l-1\ge x_l>x_{l+1}=z_{l+1}$ and $z_{l-1}=y_{l-1}>y_l\ge z_l$, so $z-\epsilon_l$ is a root sequence.
\end{proof}

\begin{lemma}[adapted root sequences]\label{lem:adapted}
Let $\gamma\in\CC$, $\gamma\ne0$, with chain $z^1>\dots>z^m$, and let $\lambda\in\frac1n\ZZ$. There is a root sequence $z$ with $\mu(z)=\lambda$ that is comparable with every $z^k$; we call it \emph{adapted} to $\gamma$ at $\lambda$. For every such $z$, $\phi_z(\gamma)=\Lambda_\gamma(\lambda)$. In particular $\phi_{z^k}(\gamma)=\Lambda_\gamma(\mu(z^k))$ for every $k$.
\end{lemma}

\begin{proof}
Put $c=c_\lambda\in\frac1n\ZZ$. If $c\ge c(z^1)$, take $z=z^1+n(c-c(z^1))\,\epsilon_1$, which lies above $z^1$, hence above every $z^k$. If $c\le c(z^m)$, take $z=z^m-n(c(z^m)-c)\,\epsilon_n$. Otherwise $c(z^{k+1})\le c\le c(z^k)$ for some $k$, and the chain of Lemma~\ref{lem:chains} for $[z^{k+1},z^k]$ passes through a root sequence $z$ with $c(z)=c$; it satisfies $z^{k+1}\le z\le z^k$, so it is comparable with every element of the chain. In all cases Lemma~\ref{lem:comparable} and linearity give $\phi_z(\gamma)=\sum_kq_k\,n!\,(c-c(z^k))_+=\Lambda_\gamma(\lambda)$.
\end{proof}

\subsection{Test sheaves}
A sheaf set-theoretically supported on a line, with a pure resolution of prescribed type, exists by the theorem on pure resolutions.

\begin{lemma}[test sheaves]\label{lem:testsheaf}
Let $z$ be a root sequence and $L\subset\PP^n$ a line. There are a nonzero coherent sheaf $\Tt$ with support $L$, an integer $B>0$ and an exact sequence
\[
0\to\OO(z_n)^{b_{n-1}}\to\dots\to\OO(z_2)^{b_1}\to\OO(z_1)^{b_0}\to \Tt\to0,\qquad b_i=\frac B{P_{i+1}(z)} .
\]
\end{lemma}

\begin{proof}
For $n=1$ take $L=\PP^1$ and $\Tt=\OO(z_1)$. Let $n\ge2$, $R=K[x_1,\dots,x_{n-1}]$ and $d_i=-z_{i+1}$, so that $d_0<\dots<d_{n-1}$. By the existence of pure resolutions, due to Eisenbud, Fløystad and Weyman in characteristic zero \cite{EFW11} and to Eisenbud and Schreyer in general \cite[Thm.~5.1]{ES09}, there is a graded $R$-module $M\ne0$ of finite length with a pure resolution $0\to R(-d_{n-1})^{b_{n-1}}\to\dots\to R(-d_0)^{b_0}\to M\to0$. By Herzog and Kühl \cite{HK84}, the Betti numbers of a pure resolution are proportional to $1/\prod_{j\ne i}|d_j-d_i|=1/P_{i+1}(z)$. Tensoring with the flat extension $R\to S=R[x_0,x_n]$ and passing to sheaves on $\PP^n=\operatorname{Proj}S$ gives a resolution of the sheaf associated with $M\otimes_RS$, which is supported on the line $x_1=\dots=x_{n-1}=0$. A linear change of coordinates moves this line to $L$.
\end{proof}

\begin{lemma}[truncation]\label{lem:truncation}
Let $F$ be a vector bundle on $\PP^n$, $z$ a root sequence and $\Tt$, $B$ as in Lemma~\textup{\ref{lem:testsheaf}}. Then
\[
h^0(F\otimes \Tt)\ \le\ \frac B{(n-1)!}\,\phi_z\big(\gamma(F)\big).
\]
\end{lemma}

\begin{proof}
Since $F$ is locally free, $F$ tensored with the resolution of $\Tt$ is a resolution $K^\bullet$ of $F\otimes \Tt$, with $K^{-i}=F(z_{i+1})^{b_i}$, $0\le i\le n-1$. The hypercohomology spectral sequence of $K^\bullet$ has $E_1^{-i,j}=H^j(F(z_{i+1}))^{b_i}$ and converges to $H^{j-i}(F\otimes \Tt)$. On each page put $S_r=\sum_{k\le0}(-1)^k\sum_{j-i=k}\dim E_r^{-i,j}$. The differentials raise the total degree $k$ by one; those between total degrees $k-1$ and $k$ with $k\le0$ do not change $S_r$, and those from total degree $0$ to total degree $1$ lower it by their rank. So $S_{r+1}\le S_r$, and at the limit $S_\infty=\sum_{k\le0}(-1)^kh^k(F\otimes \Tt)=h^0(F\otimes \Tt)$, because the hypercohomology of a sheaf vanishes in negative degrees. Hence
\[
h^0(F\otimes \Tt)\le S_1=\sum_{i=0}^{n-1}b_i\sum_{j\le i}(-1)^{j-i}h^j\big(F(z_{i+1})\big)=B\sum_{r=1}^n\frac{(-1)^{r-1}}{P_r(z)}\sum_{j<r}(-1)^jh^j\big(F(z_r)\big),
\]
with $r=i+1$, and $(-1)^{r-1}/P_r(z)=1/D_r(z)$.
\end{proof}

\subsection{The theorem}
\begin{theorem}\label{thm:main}
Let $F$ be a vector bundle on $\PP^n$ and $F'\subset F$ a subsheaf of rank $n'$.
\begin{enumerate}[label=\textup{(\alph*)}]
\item For every root sequence $z$, $\deg F'-\mu(z)\,n'\le\phi_z(\gamma(F))$.
\item $\deg F'\le\Psi_F(n')$. Equivalently, $P_F\le\Psi_F$.
\end{enumerate}
\end{theorem}

\begin{proof}
(a) If $F'=0$, the assertion is $\phi_z(\gamma(F))\ge0$, which holds by Lemma~\ref{lem:truncation}. Otherwise, the saturation of $F'$ has the same rank and at least the same degree, so we may assume that $F/F'$ is torsion-free. Then $F'$ and $F/F'$ are locally free outside a closed subset of codimension at least two, and a general line $L$ is disjoint from it. Having chosen $L$, let $\Tt$ be the sheaf of Lemma~\ref{lem:testsheaf} supported on $L$. Near $L$ the sequence $0\to F'\to F\to F/F'\to0$ is a locally split exact sequence of vector bundles, and away from $L$ the resolution of $\Tt$ is a locally split exact complex; so $F'\otimes \Tt\to F\otimes \Tt$ is injective, and $F'$ tensored with the resolution of $\Tt$ is a resolution of $F'\otimes \Tt$. Since $F'\otimes \Tt$ has one-dimensional support,
\[
\chi(F'\otimes \Tt)\ \le\ h^0(F'\otimes \Tt)\ \le\ h^0(F\otimes \Tt)\ \le\ \frac B{(n-1)!}\,\phi_z(\gamma(F))
\]
by Lemma~\ref{lem:truncation}. On the other hand $\chi(F'\otimes \Tt)=\sum_i(-1)^ib_i\chi(F'(z_{i+1}))=B\sum_r\chi(F'(z_r))/D_r(z)$, and by Riemann--Roch and \eqref{eq:divdiff},
\[
(n-1)!\sum_r\frac{\chi(F'(z_r))}{D_r(z)}=(n-1)!\Big(\frac{n'}{n!}\sum_iz_i+\frac{\deg F'+\frac{n+1}2n'}{(n-1)!}\Big)=\deg F'-\mu(z)\,n'.
\]
(b) If $F'=0$ there is nothing to prove. Otherwise apply (a) to the elements $z^k$ of the chain of $F$: by Lemma~\ref{lem:adapted}, $\deg F'\le n'\mu(z^k)+\Lambda_F(\mu(z^k))$ for every $k$, and by Lemma~\ref{lem:legendre}(i) the minimum over $k$ of the right-hand side is $\Psi_F(n')$. Since $P_F$ is the upper envelope of the points $(\rk F',\deg F')$ for subsheaves $F'\subset F$ (its vertices are the steps of the Harder--Narasimhan filtration, and every subsheaf lies below it, see \cite[\S1.3]{HL10} and \cite{Sh77}), this is $P_F\le\Psi_F$.
\end{proof}

\begin{corollary}\label{cor:supernatural}
Let $F$ be a vector bundle of rank $N$ on $\PP^n$.
\begin{enumerate}[label=\textup{(\alph*)}]
\item If $\gamma(F)$ is a multiple of a single supernatural table, then $F$ is semistable.
\item $\mu_{\max}(F)\le\mu^\BS_{\max}(F)=\mu(z^m)$ and $\mu_{\min}(F)\ge\mu^\BS_{\min}(F)=\mu(z^1)$, where $z^1$ and $z^m$ are the first and the last element of the chain, given by Lemma~\textup{\ref{lem:ends}} in terms of the supports of the rows of $\gamma(F)$.
\item $V_{HN}(F)\le V_c(F)$, where $V_{HN}(F)=\sum_iw_i(\mu_i-\mu(F))^2$ is the variance of the slopes $\mu_i$ of the Harder--Narasimhan factors $\cQ_i$ of $F$, with weights $w_i=\rk\cQ_i/N$, and $V_c(F)$ is the variance of the Boij--Söderberg slopes, with weights $p_k=n!\,q_k/N$. Equality holds if and only if $P_F=\Psi_F$.
\end{enumerate}
\end{corollary}

\begin{proof}
(a) Then $\Psi_F$ is linear of slope $\mu(F)$, and Theorem~\ref{thm:main} says $\mu(F')\le\mu(F)$ for every subsheaf. (b) For small $x>0$ we have $P_F(x)=x\,\mu_{\max}(F)$ and $\Psi_F(x)=x\,\mu^\BS_{\max}(F)$, so $P_F\le\Psi_F$ gives the first inequality; the second follows at the other end, or by \eqref{eq:dualpsi}. (c) Both polygons are concave with the same endpoints, so $P_F\le\Psi_F$ says that the distribution of the Harder--Narasimhan slopes is majorized by that of the Boij--Söderberg slopes, both with mean $\mu(F)$; apply Karamata's inequality \cite{Ka32} to $x\mapsto x^2$. After clearing denominators both distributions have the same finite total mass; for the strictly convex function $x^2$, equality in the inequality of Hardy, Littlewood and P\'olya \cite{HLP} forces the two distributions to coincide, i.e.\ $P_F=\Psi_F$. The converse is clear.
\end{proof}

\begin{proposition}[subsheaves of the same slope]\label{prop:sameslope}
Let $F$ be a vector bundle on $\PP^n$ whose table is a multiple of $\gamma^z$, and let $G\subset F$ be a saturated subsheaf with $0\ne G\ne F$ and $\mu(G)=\mu(F)$. Then $G$ and $F/G$ are vector bundles whose tables are multiples of $\gamma^z$. Consequently $F$ is Gieseker semistable, and $F$ is Gieseker stable if and only if it is $\mu$-stable.
\end{proposition}

\begin{proof}
Let $D=z_1-z_n+1$, let $W$ be the set of integers in $[z_n,z_1]$ that are not roots, so $|W|=D-n$, and let $X=\PP^n\times(\PP^1)^W$ with the very ample $A=\OO_X(1,\dots,1)$; $\dim X=D$. For a coherent sheaf $\cT$ on $\PP^n$ put
\[
\cU_\cT=\Big(\cT\boxtimes\mathop{\boxtimes}_{w\in W}\OO_{\PP^1}(-w-1)\Big)\big((z_1+1)A\big).
\]
For $1\le j\le D$ put $t=z_1+1-j\in[z_n,z_1]$; by the Künneth formula $H^*(\cU_\cT(-jA))=H^*(\cT(t))\otimes\bigotimes_{w}H^*(\OO_{\PP^1}(t-w-1))$. If $t$ is a root, $H^*(F(t))=0$; if $t\in W$, the factor $\OO_{\PP^1}(-1)$ is acyclic. So $\cU_F$ has no cohomology at the twists $-A,\dots,-DA$, i.e.\ it is an Ulrich bundle on $(X,A)$ \cite[Prop.~2.1]{ESW03}. (This goes in the opposite direction to the construction of supernatural bundles in \cite[\S6]{ES09}, as push-forwards of line bundles from products of projective spaces under finite linear projections.) The functor $\cT\mapsto\cU_\cT$ is exact and preserves torsion-freeness, so $0\to\cU_G\to\cU_F\to\cU_{F/G}\to0$ is exact with torsion-free quotient. For $\cT$ of rank $s>0$, $c_1(\cU_\cT)\cdot A^{D-1}=\frac{(D-1)!}{(n-1)!}\deg\cT+s\kappa$ with $\kappa$ independent of $\cT$, so $\mu_A(\cU_G)=\mu_A(\cU_F)$. By \cite[Thm.~2.9(b)]{CH12} (Casanellas and Hartshorne work over an algebraically closed field of arbitrary characteristic, and the proof uses Bertini's theorem for general hyperplane sections, Serre vanishing and Serre duality), $\cU_G$ and $\cU_{F/G}$ are Ulrich bundles; restricting to $\PP^n\times\{p\}$ shows that $G$ and $F/G$ are locally free. For $t=z_i$, every factor $\OO_{\PP^1}(z_i-w-1)$, $w\in W$, has a nonzero cohomology group, while $\cU_G(-(z_1+1-z_i)A)$ is acyclic; so $H^*(G(z_i))=0$, and likewise $H^*((F/G)(z_i))=0$. Hence the columns $z_1,\dots,z_n$ of the tables of $G$ and $F/G$ vanish, and since these tables are positive combinations of supernatural tables, every element of their chains is $z$.

The reduced Hilbert polynomial of a bundle with table $q\gamma^z$ is $\prod_i(t-z_i)/n!$. If $G'\subset F$ is a subsheaf with $\mu(G')<\mu(F)$, its reduced Hilbert polynomial is smaller. If $\mu(G')=\mu(F)$, its saturation $G''$ has $\mu(G'')\ge\mu(G')=\mu(F)$, hence $\mu(G'')=\mu(F)$ by Corollary~\ref{cor:supernatural}(a); if $G''\ne F$ it has the reduced Hilbert polynomial of $F$ by the first part, and if $G''=F$ this is trivial. The reduced Hilbert polynomial of $G'$ is at most that of $G''$, since $G''/G'$ is a torsion sheaf. So $F$ is Gieseker semistable; and if it is not $\mu$-stable, a saturated $G\subset F$ with $\mu(G)=\mu(F)$ has the same reduced Hilbert polynomial as $F$, so $F$ is not Gieseker stable. A $\mu$-stable sheaf is Gieseker stable.
\end{proof}

\begin{remark}\label{rem:main}
\begin{enumerate}[label=(\arabic*)]
\item Equality $P_F=\Psi_F$ holds whenever the Harder--Narasimhan factors $\cQ_j=F_j/F_{j-1}$, $1\le j\le s$, of $F$ are vector bundles with tables $q_j\gamma^{z(j)}$ and pairwise comparable $z(j)$, for instance for a direct sum of supernatural bundles with comparable root sequences. Indeed the slopes decrease, so $z(1)<\dots<z(s)$; for $z'<z$ the row of $\gamma^{z'}$ in each column is at most the row of $\gamma^z$, so, by induction on $j$, at each twist the cohomology of $F_{j-1}$ lies in rows at most the row of $\cQ_j$, every connecting homomorphism $H^i(\cQ_j(t))\to H^{i+1}(F_{j-1}(t))$ vanishes, and $\gamma(F_j)=\gamma(F_{j-1})+\gamma(\cQ_j)$. Then $\gamma(F)=\sum_j\gamma(\cQ_j)$ is a positive combination along a chain, hence the chain decomposition, and its slopes and lengths are those of $P_F$. Consequently Theorem~\ref{thm:main} is optimal among bounds that depend only on the ray of the table: by \cite[Thm.~6.1]{ES09} every ray $\RR_{>0}\gamma^z$ contains the table of a supernatural bundle, so every $\gamma\in\CC$ has a positive integer multiple $M\gamma$ that is the table of a direct sum $F$ of supernatural bundles along the chain of $\gamma$, and then $P_F(Mm)=M\Psi_\gamma(m)$. So $\Psi_\gamma(m)$ is the maximum of $P_F(Mm)/M$ over all $M\ge1$ and all vector bundles $F$ with $\gamma(F)=M\gamma$. For a fixed table integrality gives more: the vertices of $P_F$ are lattice points, so $P_F\le\Psi^{\ZZ}_F$, the upper concave hull of the points $(m,\lfloor\Psi_F(m)\rfloor)$, $m=0,\dots,\rk F$, which also depends only on $\gamma(F)$. In particular $F$ is semistable if $\Psi^{\ZZ}_F$ is the segment from $(0,0)$ to $(\rk F,\deg F)$. This happens for the table of Example~\ref{ex:sum} below: there $\Psi_F$ takes the values $0,-\frac52,-\frac{16}3,-\frac{17}2,-12$ at $m=0,\dots,4$, so every bundle with this table is semistable, although $\Psi_F$ is not linear.
\item The root sequences of the chain lie within the cohomology range of $F$ (proof of \cite[Cor.~0.6]{ES09}), so $\mu^\BS_{\max}(F)$ and $\mu^\BS_{\min}(F)$ lie between the two bounds that \cite[Cor.~0.6]{ES09} gives for $\mu(F)$ in terms of the cohomology range. By Corollary~\ref{cor:supernatural}(b), the upper bound holds for $\mu_{\max}(F)$, i.e.\ for the slope of every subsheaf of $F$, and the lower bound for $\mu_{\min}(F)$, i.e.\ for the slope of every torsion-free quotient. The reverse bounds fail: on $\PP^2$, $0\to\OO(-d)\to\OO^2\to\cI_Z(d)\to0$, with $Z$ a complete intersection of two curves of degree $d$, has a quotient of slope $d$ and a subsheaf of slope $-d$.
\item Two comparable root sequences with the same slope are equal, so $\Psi_F$ is linear iff $\gamma(F)$ is a multiple of a single $\gamma^z$; Corollary~\ref{cor:supernatural}(a) is the case of a linear polygon. If the roots form an arithmetic progression, $z_i=z_1-(i-1)k$, then the bundle $F(z_1+k)$ has no cohomology at the twists $-k,-2k,\dots,-nk$, so it is an Ulrich bundle on $(\PP^n,\OO(k))$, and its semistability also follows from \cite[Thm.~2.9]{CH12}. For arbitrary $z$ the Segre product of Proposition~\ref{prop:sameslope} gives the same reduction: a destabilizing subsheaf $G\subset F$ would give a destabilizing subsheaf $\cU_G\subset\cU_F$ of the Ulrich bundle $\cU_F$, contradicting \cite[Thm.~2.9(a)]{CH12}; this is a second proof of Corollary~\ref{cor:supernatural}(a), in any characteristic. In characteristic zero, Corollary~\ref{cor:supernatural}(a) can also be deduced from \cite{ErSa16}: for $n+1$ integers containing the roots $z$, the transform $\Phi^W_2$ of Erman and Sam is exact on coherent sheaves \cite[Prop.~4.4]{ErSa16}, multiplies classes in $K_0$ by a positive integer $N_W$ \cite[Cor.~1.4]{ErSa16}, and takes a bundle $F$ with table $q\gamma^z$ to a direct sum of copies of the homogeneous bundle $\mathbf S_\lambda(\mathcal Q^*)(w)$ of \cite[(3.2), Cor.~1.9]{ErSa16}, which is stable, being irreducible homogeneous \cite{Ra66,Um78}; a destabilizing subsheaf of $F$ would give a destabilizing subsheaf of $\Phi^W_2(F)$. Theorem~\ref{thm:main} needs neither the characteristic nor a single ray. Appendix~\ref{app:P3} describes the bundles with supernatural tables on $\PP^3$ with $z_1-z_3\le4$.
\item The proof gives the equality cases of Theorem~\ref{thm:main}(a): for $F/F'$ torsion-free and $L$ general, $\deg F'-\mu(z)n'=\phi_z(\gamma(F))$ iff $h^1(F'\otimes \Tt)=0$, $H^0(F'\otimes \Tt)=H^0(F\otimes \Tt)$, and every differential of the spectral sequence of Lemma~\ref{lem:truncation} from total degree $0$ to total degree $1$ vanishes. Proposition~\ref{prop:marginPn} below expresses the defect in terms of the cohomology of $F'$ and $F/F'$.
\end{enumerate}
\end{remark}

The defect of (a) has an exact expression in terms of the cohomology of $F'$ and $F/F'$ at the $n$ twists $z_1,\dots,z_n$.

\begin{proposition}[the defect]\label{prop:marginPn}
Let $0\to F'\to F\to F''\to0$ be an exact sequence of coherent sheaves on $\PP^n$ with $F$ a vector bundle and $F'\ne0$, $n'=\rk F'$, $\mu'=\mu(F')$, let $\lambda\in\frac1n\ZZ$, and let $z$ be adapted to $\gamma(F)$ at $\lambda$. Write $h^j_t(G)=h^j(G(t))$, and let $\partial^j_t\colon H^j(F''(t))\to H^{j+1}(F'(t))$ be the connecting homomorphisms. Then
\begin{equation}\label{eq:marginPn}
\Lambda_F(\lambda)-n'(\mu'-\lambda)\ =\ (n-1)!\sum_{r=1}^n\frac{B_r}{P_r(z)},
\end{equation}
where
\[
B_r=\sum_{j=r}^{n}(-1)^{j-r}h^j_{z_r}(F')+\sum_{j=0}^{r-1}(-1)^{r-1-j}h^j_{z_r}(F'')-\rk\partial^{r-1}_{z_r}.
\]
\end{proposition}

\begin{proof}
For a sheaf $G$ put $A_r(G)=\sum_{j<r}(-1)^jh^j_{z_r}(G)$. By Lemma~\ref{lem:adapted},
\[
\Lambda_F(\lambda)=(n-1)!\sum_r\frac{(-1)^{r-1}A_r(F)}{P_r(z)} .
\]
Fix $r$ and $t=z_r$. In the long exact cohomology sequence, $h^j_t(F)=h^j_t(F')+h^j_t(F'')-\rk\partial^j_t-\rk\partial^{j-1}_t$, and the alternating sum over $j<r$ telescopes:
\[
A_r(F)=A_r(F')+A_r(F'')-(-1)^{r-1}\rk\partial^{r-1}_t .
\]
Moreover $A_r(F')=\chi(F'(t))-\sum_{j\ge r}(-1)^jh^j_t(F')$. Hence $(-1)^{r-1}A_r(F)=(-1)^{r-1}\chi(F'(z_r))+B_r$, and $(n-1)!\sum_r(-1)^{r-1}\chi(F'(z_r))/P_r(z)=n'(\mu'-\lambda)$, as in the proof of Theorem~\ref{thm:main}.
\end{proof}

For $n\ge3$ the intermediate cohomology enters $B_r$ with both signs, and the nonnegativity of the right-hand side of \eqref{eq:marginPn}, which Theorem~\ref{thm:main} implies, is not visible term by term. For $n=2$ the $B_r$ can still be negative, but after regrouping the summands the defect becomes a sum of four nonnegative terms: see Proposition~\ref{prop:fourdefect} and the paragraph after it.

\section{Majorization and monotonicity}\label{sec:monotone}

\subsection{The functionals as upper bounds}
Applied to supernatural bundles, Theorem~\ref{thm:main}(a) and Lemma~\ref{lem:truncation} give lower bounds for $\phi_z$ on all supernatural tables.

\begin{lemma}\label{lem:positivity}
For all root sequences $y$ and $z$,
\[
\phi_z(\gamma^y)\ \ge\ n!\,\big(c(z)-c(y)\big)_+ ,
\]
with equality if $y$ and $z$ are comparable.
\end{lemma}

\begin{proof}
Let $E$ be a supernatural bundle with root sequence $y$ \cite[Thm.~0.4]{ES09}, so that $\gamma(E)=q\gamma^y$ with $q>0$, $\rk E=n!\,q$ and $\mu(E)=\mu(y)$. Lemma~\ref{lem:truncation} gives $\phi_z(\gamma(E))\ge0$, and Theorem~\ref{thm:main}(a) with $F'=F=E$ gives $\phi_z(\gamma(E))\ge\rk E\,(\mu(y)-\mu(z))=n!\,q\,(c(z)-c(y))$. Divide by $q$. Equality for comparable $y,z$ is Lemma~\ref{lem:comparable}.
\end{proof}

\begin{proposition}\label{prop:minformula}
Let $\gamma\in\CC$ and $\lambda\in\frac1n\ZZ$. For every decomposition $\gamma=\sum_yq'_y\gamma^y$ with $q'_y\ge0$ and every root sequence $z$ with $\mu(z)=\lambda$,
\[
\sum_yq'_y\,n!\,\big(\mu(y)-\lambda\big)_+\ \le\ \phi_z(\gamma).
\]
For the chain decomposition the left-hand side is $\Lambda_\gamma(\lambda)$, and equality holds when $z$ is adapted to $\gamma$ at $\lambda$. Hence
\[
\Lambda_\gamma(\lambda)=\min\{\phi_z(\gamma):\ \mu(z)=\lambda\}\ \ge\ \sum_yq'_y\,n!\,\big(\mu(y)-\lambda\big)_+ .
\]
\end{proposition}

\begin{proof}
The first inequality is Lemma~\ref{lem:positivity} and linearity, since $c(z)-c(y)=\mu(y)-\lambda$. The equality is Lemma~\ref{lem:adapted}.
\end{proof}

So each $\phi_z$ is an upper bound for $\Lambda_\lambda=\Lambda_{(\cdot)}(\lambda)$ on all of $\CC$, and $\Lambda_\lambda$ is their minimum, attained at the root sequences comparable with the chain.

\begin{theorem}\label{thm:concave}
Let $\gamma\in\CC$, and let $\gamma=\sum_jq'_j\gamma^{y_j}$ be any decomposition with $q'_j\ge0$.
\begin{enumerate}[label=\textup{(\alph*)}]
\item The distribution of centers of the chain, with masses $q_k$, majorizes the distribution of the $c(y_j)$, with masses $q'_j$. Equivalently, the slope polygon of the decomposition lies below $\Psi_\gamma$, and $\Lambda_\gamma(\lambda)\ge\sum_jn!\,q'_j(\mu(y_j)-\lambda)_+$ for all $\lambda$.
\item For each $\lambda$ the function $\gamma\mapsto\Lambda_\gamma(\lambda)$ is concave, positively homogeneous and superadditive on $\CC$. For every convex function $f$ on $\frac1n\ZZ$ the function $\gamma\mapsto\sum_kq_kf(c(z^k))$, computed on the chain, is concave; in particular so is $\gamma\mapsto\rk\gamma\cdot V_c(\gamma)$. The function $(\gamma,m)\mapsto\Psi_\gamma(m)$, on $\{(\gamma,m):0\le m\le\rk\gamma\}$, is jointly concave and positively homogeneous; in particular $\Psi_{\gamma+\gamma'}(m+m')\ge\Psi_\gamma(m)+\Psi_{\gamma'}(m')$.
\end{enumerate}
\end{theorem}

\begin{proof}
(a) Both sides of the last inequality are convex and piecewise linear in $\lambda$, with breakpoints in $\frac1n\ZZ$, and they agree for $\lambda$ large; by Proposition~\ref{prop:minformula} the inequality holds on $\frac1n\ZZ$, hence everywhere. Both distributions have the same total mass and mean, because $\rk$ and $\deg$ are linear, so the equivalences are Lemma~\ref{lem:legendre}(ii).

(b) By Proposition~\ref{prop:minformula}, $\Lambda_\lambda$ is a minimum of linear functionals for $\lambda\in\frac1n\ZZ$, hence concave and positively homogeneous, so superadditive; for other $\lambda$ it is a convex combination of its values at the neighboring points of $\frac1n\ZZ$. Let $\gamma,\gamma'\in\CC$ be nonzero (if one of them is zero there is nothing to prove). The chains of all $s\gamma+t\gamma'$ with $s,t\in\QQ_{\ge0}$, not both zero, lie between the termwise minimum of the last elements and the termwise maximum of the first elements of the chains of $\gamma$ and $\gamma'$ (Lemma~\ref{lem:ends}), so only finitely many centers occur on the segment. On a finite set of centers, a convex function $f$ on $\frac1n\ZZ$ is an affine function plus a nonnegative combination of the functions $x\mapsto(a-x)_+$, $a\in\frac1n\ZZ$, and $\sum_kq_k(a-c(z^k))_+=\Lambda_\gamma(\lambda_a)/n!$ with $c_{\lambda_a}=a$, while $\sum_kq_k$ and $\sum_kq_kc(z^k)$ are linear in $\gamma$. Further, with $\rk\gamma\cdot V_c(\gamma)=0$ for $\gamma=0$, $\rk\gamma\cdot V_c(\gamma)=n!\sum_kq_kc(z^k)^2-\big(n!\sum_kq_kc(z^k)\big)^2/\rk\gamma$, and the second term is convex in $\gamma$, being a square of a linear function divided by a linear function. Finally, by Lemma~\ref{lem:legendre}(i), $\Psi_\gamma(m)$ is the minimum over $\lambda$ of $m\lambda+\Lambda_\gamma(\lambda)$, which is jointly concave and positively homogeneous in $(\gamma,m)$.
\end{proof}

For fixed $m>0$, $\gamma\mapsto\Psi_\gamma(m)$ is concave but neither homogeneous nor superadditive: $\Psi_\gamma(1)=\Psi_{2\gamma}(1)=1$ for $\gamma=\gamma(\OO(1)\oplus\OO(-1))$.

\begin{example}\label{ex:sum}
On $\PP^2$, let $F=\pi_*\OO_{\PP^1\times\PP^1}(-4,-1)\oplus\OO(-3)^{\oplus2}$, where $\pi$ is a finite linear projection of the quadric \cite[Thm.~6.1]{ES09}. Then $\gamma(F)=\gamma^{(3,0)}+\gamma^{(2,1)}$, with incomparable root sequences, and the chain decomposition is $\gamma(F)=\frac23\big(\gamma^{(3,1)}+\gamma^{(3,0)}+\gamma^{(2,0)}\big)$. The given decomposition has centers $\frac32,\frac32$; the chain has centers $2,\frac32,1$, each with mass $\frac23$, a mean-preserving spread.
\end{example}

\begin{proposition}[decomposition criterion]\label{prop:splitting}
Let $\gamma\in\CC$, $m\in\QQ$ with $0\le m\le\rk\gamma$, and $d\in\QQ$. Then $\Psi_\gamma(m)\ge d$ if and only if
\[
\gamma=x+y\qquad\text{with }x,y\in\CC,\quad \rk x=m,\quad \deg x\ge d .
\]
\end{proposition}

\begin{proof}
Let $\gamma=\sum_kq_k\gamma^{z^k}$ be the chain decomposition. If $\Psi_\gamma(m)\ge d$, choose rational $0\le x_k\le q_k$ with $\sum_kn!\,x_k=m$ and $\sum_kn!\,x_k\mu(z^k)=\Psi_\gamma(m)$, and put $x=\sum_kx_k\gamma^{z^k}$, $y=\gamma-x$. Conversely, decompositions of $x$ and of $y$ together form a decomposition of $\gamma$; by Theorem~\ref{thm:concave}(a) its slope polygon lies below $\Psi_\gamma$, and taking all the pieces of $x$ shows that this polygon is at least $\deg x\ge d$ at $m$.
\end{proof}

\begin{corollary}[additive extensions]\label{cor:additive}
Let $0\to F'\to F\to F''\to0$ be an exact sequence of vector bundles on $\PP^n$ with $\gamma(F)=\gamma(F')+\gamma(F'')$, i.e.\ with all connecting homomorphisms zero. Then $\Psi_F(m'+m'')\ge\Psi_{F'}(m')+\Psi_{F''}(m'')$ for $0\le m'\le\rk F'$, $0\le m''\le\rk F''$.
\end{corollary}

\begin{proof}
This is the superadditivity of Theorem~\ref{thm:concave}(b).
\end{proof}

\subsection{Monotonicity}
A table $\Delta$ is \emph{$\chi$-neutral} if $\sum_j(-1)^j\Delta_{j,t}=0$ for all $t$. Two tables in $\CC$ have the same Hilbert polynomial iff their difference is $\chi$-neutral. The cancellation vectors $e_{j,t}+e_{j+1,t}$ are $\chi$-neutral, and by the long exact sequences, for every exact sequence $0\to F'\to F\to F''\to0$,
\begin{equation}\label{eq:extdiff}
\gamma(F')+\gamma(F'')-\gamma(F)=\sum_t\sum_{j=0}^{n-1}\rk\big(\partial^j_t\big)\big(e_{j,t}+e_{j+1,t}\big),\qquad \partial^j_t\colon H^j(F''(t))\to H^{j+1}(F'(t)).
\end{equation}
On $\PP^2$ every nonnegative $\chi$-neutral table is a nonnegative combination of cancellation vectors, since a column $(a,a+b,b)$ is $a(1,1,0)+b(0,1,1)$. For $n\ge3$ this fails: $e_{0,t}+e_{3,t}$ is nonnegative and $\chi$-neutral, but not a nonnegative combination of cancellation vectors.

\begin{lemma}[local values]\label{lem:localvalues}
Let $z$ be a root sequence, $1\le r\le n$ and $t\in\ZZ$. Then
\[
\phi_z\big(e_{r-1,t}+e_{r,t}\big)=v_{r,t}(z)=\begin{cases}\dfrac{(n-1)!}{P_r(z)} & \text{if } z_r=t,\\[2ex] 0 & \text{otherwise.}\end{cases}
\]
\end{lemma}

\begin{proof}
By \eqref{eq:phi}, $\phi_z(e_{j,t})=(n-1)!(-1)^j/D_s(z)$ if $t=z_s$ and $j<s$, and $\phi_z(e_{j,t})=0$ otherwise. If $t=z_s$ with $s>r$, the contributions of $j=r-1$ and $j=r$ cancel; if $s=r$, only $j=r-1$ contributes, with $(-1)^{r-1}/D_r(z)=1/P_r(z)$; if $s<r$, neither contributes.
\end{proof}

\begin{theorem}[monotonicity]\label{thm:monotone}
Let $\gamma,\gamma'\in\CC$, and suppose that $\gamma'-\gamma$ is a nonnegative combination of cancellation vectors. Then $\Lambda_\gamma(\lambda)\le\Lambda_{\gamma'}(\lambda)$ for every $\lambda$; equivalently $\Psi_\gamma\le\Psi_{\gamma'}$, i.e.\ the center distribution of the chain of $\gamma'$ majorizes that of the chain of $\gamma$. On $\PP^2$ the hypothesis says that $\gamma\le\gamma'$ entrywise and that $\gamma,\gamma'$ have the same Hilbert polynomial.
\end{theorem}

\begin{proof}
Let $\lambda\in\frac1n\ZZ$ and let $z$ be adapted to $\gamma'$ at $\lambda$. By Lemma~\ref{lem:adapted}, Lemma~\ref{lem:localvalues} and Proposition~\ref{prop:minformula},
\[
\Lambda_{\gamma'}(\lambda)=\phi_z(\gamma')=\phi_z(\gamma)+\phi_z(\gamma'-\gamma)\ \ge\ \phi_z(\gamma)\ \ge\ \Lambda_\gamma(\lambda).
\]
Both sides are piecewise linear with breakpoints in $\frac1n\ZZ$, so the inequality holds for all $\lambda$. Since $\gamma'-\gamma$ is $\chi$-neutral, $\gamma$ and $\gamma'$ have the same rank and degree, and the equivalence with $\Psi_\gamma\le\Psi_{\gamma'}$ is Lemma~\ref{lem:legendre}(ii). On $\PP^2$ the hypothesis is equivalent to the stated one by the remark before the lemma.
\end{proof}

\begin{corollary}\label{cor:monoext}
\begin{enumerate}[label=\textup{(\alph*)}]
\item For every exact sequence $0\to F'\to F\to F''\to0$ of vector bundles on $\PP^n$, $\Psi_F\le\Psi_{F'\oplus F''}$. Hence $\Psi_F\le\Psi_{\gr F}$ for every filtration of $F$ with locally free subquotients, and $V_c(F)\le V_c(\gr F)$.
\item Let $(F_s)_{s\in S}$ be a flat family of vector bundles on $\PP^n$ over a variety $S$. Every $s_0\in S$ has an open neighborhood $U$ with $\Psi_{F_s}\le\Psi_{F_{s_0}}$ for all $s\in U$. If $S$ is irreducible, then $\Psi_{F_s}\le\Psi_{F_{s_0}}$ for all $s$ in a dense open subset and every $s_0\in S$.
\item Let $F',F''$ be vector bundles and let $F$ be an extension of $F''$ by $F'$ with general class in $\Ext^1(F'',F')$. Then $\Psi_F\le\Psi_{F_1}$ for every extension $F_1$ of $F''$ by $F'$.
\end{enumerate}
\end{corollary}

\begin{proof}
(a) By \eqref{eq:extdiff}, $\gamma(F'\oplus F'')-\gamma(F)$ is a nonnegative combination of cancellation vectors; apply Theorem~\ref{thm:monotone}. This relation between tables is transitive and is preserved by adding the same table to both sides, so we may induct along the filtration. The center distributions of $F$ and $\gr F$ have the same mean, and the second majorizes the first; applying the convex function $x\mapsto x^2$ gives $V_c(F)\le V_c(\gr F)$.

(b) We may shrink $S$ to an affine neighborhood $\operatorname{Spec}A$ of $s_0$. Being $m$-regular is an open condition on $S$, since it is the vanishing of finitely many cohomology groups, which are upper semicontinuous in $s$; so after shrinking $S$ the fibers $F_s$ and $F_s^\vee$ are $m$-regular for a fixed $m$. This gives a bounded range of $t$ outside which the column $t$ of $\gamma(F_s)$ has a single nonzero row, equal to $\pm\chi(F_s(t))$, which does not depend on $s$. Fix $t$ in that range. There is a complex $K^0\to K^1\to\dots\to K^n$ of finitely generated projective, hence after shrinking $S$ free, $A$-modules computing $H^j(F_s(t))$ for every $s$, after base change to $s$ \cite[Ch.~II, \S5]{Mu70}. With $\rho_j(s)$ the rank of $K^j\otimes k(s)\to K^{j+1}\otimes k(s)$, we get $h^j(F_s(t))=\rk K^j-\rho_j(s)-\rho_{j-1}(s)$. Each $\rho_j$ is lower semicontinuous, so $\rho_j(s)\ge\rho_j(s_0)$ for all $j$ and all $s$ in an open neighborhood $U_t$ of $s_0$, and for these $s$
\[
\gamma(F_{s_0})-\gamma(F_s)=\sum_{j=0}^{n-1}\big(\rho_j(s)-\rho_j(s_0)\big)\big(e_{j,t}+e_{j+1,t}\big)\quad\text{in the column }t .
\]
On the intersection $U$ of the $U_t$, $t$ in the bounded range, $\gamma(F_{s_0})-\gamma(F_s)$ is a nonnegative combination of cancellation vectors, and Theorem~\ref{thm:monotone} applies. If $S$ is irreducible, then, $S$ being quasi-compact, finitely many open sets as above cover it, so there is one $m$ for which all $F_s$ and $F_s^\vee$, $s\in S$, are $m$-regular; the locus $V$ where the finitely many $h^j(F_s(t))$ with $t$ in the resulting bounded range take their minimal values is dense and open, $\gamma(F_s)$ is constant on $V$, and $V$ meets every neighborhood $U$ as above. (c) Apply (b) to the family of extensions over $\Ext^1(F'',F')$. The proof of (b) works for a vector bundle on $\PP^n_S$ over any noetherian scheme $S$, with the tables, and their chain decompositions, of the geometric fibers, over the algebraic closures of the residue fields; for instance, over the spectrum of a discrete valuation ring of mixed characteristic it gives $\Psi_{F_\eta}\le\Psi_{F_s}$ for the generic and the special fiber.
\end{proof}

The same computation bounds how much cancellation can lower $\Lambda$.

\begin{proposition}[defect inequality]\label{prop:defectineq}
Let $\gamma,\gamma'\in\CC$ with
\[
\gamma'-\gamma=\sum_{r=1}^n\sum_td^r_t\,\big(e_{r-1,t}+e_{r,t}\big),\qquad d^r_t\ge0 ,
\]
let $\lambda\in\frac1n\ZZ$, and let $z$ be adapted to $\gamma$ at $\lambda$. Then
\begin{equation}\label{eq:defect}
\Lambda_\gamma(\lambda)\ \ge\ \Lambda_{\gamma'}(\lambda)-\sum_{r=1}^nd^r_{z_r}\,\frac{(n-1)!}{P_r(z)}.
\end{equation}
If $z'$ is adapted to $\gamma'$ at $\lambda$, then also $\Lambda_{\gamma'}(\lambda)-\Lambda_\gamma(\lambda)\ge\sum_{r=1}^nd^r_{z'_r}\,(n-1)!/P_r(z')$.
\end{proposition}

\begin{proof}
By Lemma~\ref{lem:adapted}, Proposition~\ref{prop:minformula} and Lemma~\ref{lem:localvalues}, $\Lambda_\gamma(\lambda)=\phi_z(\gamma)=\phi_z(\gamma')-\phi_z(\gamma'-\gamma)\ge\Lambda_{\gamma'}(\lambda)-\sum_{r,t}d^r_t\,v_{r,t}(z)$. Similarly $\Lambda_{\gamma'}(\lambda)=\phi_{z'}(\gamma)+\phi_{z'}(\gamma'-\gamma)\ge\Lambda_\gamma(\lambda)+\sum_{r,t}d^r_t\,v_{r,t}(z')$.
\end{proof}

For an exact sequence $0\to F'\to F\to F''\to0$ of vector bundles, \eqref{eq:extdiff} gives $d^r_t=\rk\partial^{r-1}_t$. So Theorem~\ref{thm:monotone} and \eqref{eq:defect} bound $\Lambda_F$ from both sides: cancellation lowers $\Lambda_F(\lambda)$ below $\Lambda_{F'\oplus F''}(\lambda)$, but at most by the ranks of the connecting homomorphisms at the $n$ twists $z_1,\dots,z_n$, weighted by $(n-1)!/P_r(z)$.

For $n\ge3$ the hypothesis of Theorem~\ref{thm:monotone} cannot be weakened to entrywise comparison with the same Hilbert polynomial. The local value of $e_{0,t}+e_{3,t}$ at a root sequence $z$ with $z_2=t$ is $\phi_z(e_{0,t})=(n-1)!/D_2(z)<0$, and this already happens for split bundles.

\begin{example}[entrywise comparison is not enough]\label{ex:weak}
On $\PP^3$ let $F=\OO(-4)^{6}\oplus\OO(-2)^{6}$ and $G=\OO(-5)\oplus\OO(-4)^{2}\oplus\OO(-3)^{6}\oplus\OO(-2)^{2}\oplus\OO(-1)$. In the Grothendieck group $G-F=\OO(-5)-4\,\OO(-4)+6\,\OO(-3)-4\,\OO(-2)+\OO(-1)$, and the Hilbert polynomials agree because the fourth difference of the cubic $a\mapsto\chi(\OO(a+t))$ vanishes. Moreover $\gamma(G)-\gamma(F)=e_{0,1}+e_{3,1}$: in the column $1$ the table of $F$ vanishes, while $h^0(\OO(-1)(1))=h^3(\OO(-5)(1))=1$, and in all other columns the two tables agree. So $\gamma(F)\le\gamma(G)$ entrywise, with the same Hilbert polynomial. For split bundles $\Psi=P$, and $\Psi_F(6)=-12>-14=\Psi_G(6)$, while $\Psi_F(1)=-2<-1=\Psi_G(1)$. So the polygons are not comparable. In terms of chains, $\gamma(F)=\gamma^{(3,2,1)}+\gamma^{(1,0,-1)}$ and $\gamma(G)=\frac16\gamma^{(4,3,2)}+\frac13\gamma^{(3,2,1)}+\gamma^{(2,1,0)}+\frac13\gamma^{(1,0,-1)}+\frac16\gamma^{(0,-1,-2)}$.
\end{example}

The same happens on every $\PP^n$, $n\ge3$: appending the tail $(-3,-4,\dots,-(n-1))$ to each root sequence of Example~\ref{ex:weak}, with the same coefficients, gives tables $\gamma\le\gamma'$ in $\CC$ with the same Hilbert polynomial and $\gamma'-\gamma=\frac{n!}6(e_{0,1}+e_{3,1})$, whose polygons are not comparable: $\Psi_\gamma-\Psi_{\gamma'}$ is $-\frac{(n-1)!}2$ at $m=\frac{n!}6$ and $(n-1)!$ at $m=n!$.

Theorem~\ref{thm:monotone} has no analogue for the Harder--Narasimhan polygon as a function of the table, since $P_F$ is not determined by $\gamma(F)$. For instance, let $\pi\colon\PP^1\times\PP^1\to\PP^2$ be the projection of a smooth quadric surface from a point off it, so $\pi^*\OO(1)=\OO(1,1)$, $G=\pi_*\OO(-1,3)$ and $H=G^\vee$. Then $\gamma(G)=\gamma^{(0,-4)}$ by the Künneth formula, $G$ is stable (semistable by Corollary~\ref{cor:supernatural}(a), of rank $2$ and odd degree), $\gamma(H)=\gamma^{(1,-3)}$, and $\Ext^1(G,H)$ has dimension $-\chi(G,H)=15$. For every extension $0\to H\to E\to G\to0$ the connecting maps vanish, since $H^0(G(t))$ and $H^1(H(t))$, and likewise $H^1(G(t))$ and $H^2(H(t))$, are never both nonzero; so $\gamma(E)=\gamma(G\oplus H)$. But a nonsplit $E$ is semistable: $H^0(E)=H^0(E^\vee)=0$ excludes destabilizing subsheaves of rank $1$ and $3$, and by the stability of $G$ and $H$ a saturated subsheaf of rank $2$ with $c_1>0$ maps isomorphically onto $G$ and splits the extension. So $P_E$ is a segment, while $P_{G\oplus H}$ has slopes $\pm\frac12$. For extensions and specializations, however, $P_F\le P_{F'\oplus F''}$ and Shatz's semicontinuity give the analogous monotonicity \cite{Sh77}. Nor is the chain itself monotone: on $\PP^2$, $\gamma=\gamma^{(4,3)}+\gamma^{(4,2)}$ and $\gamma'=\gamma+e_{1,2}+e_{2,2}=\frac32\gamma^{(4,3)}+\frac12\gamma^{(4,1)}$, and neither chain contains the other.

\subsection{Coherent sheaves}\label{sec:coherent}
Let $\cF$ be a coherent sheaf on $\PP^n$, with $\rk\cF\ge0$, $\deg\cF=c_1(\cF)\cdot h^{n-1}$ and torsion subsheaf $\cF_{\rm tors}$; $\deg\cF_{\rm tors}\ge0$ is the degree of the $(n-1)$-dimensional part of $\cF_{\rm tors}$. Following Proposition~\ref{prop:minformula}, put
\[
\Lambda_\cF(\lambda)=\inf\big\{\phi_z(\gamma(\cF)):\ \mu(z)=\lambda\big\},\ \ \lambda\in\tfrac1n\ZZ,\qquad \Psi_\cF(m)=\inf_{\lambda\in\frac1n\ZZ}\big(m\lambda+\Lambda_\cF(\lambda)\big),\ \ 0\le m\le\rk\cF .
\]
For vector bundles these are the functions defined before, by Proposition~\ref{prop:minformula} and Lemma~\ref{lem:legendre}(i). Let $P_\cF$ be the upper envelope of the points $(\rk\cF',\deg\cF')$ for coherent subsheaves $\cF'\subset\cF$; it is the Harder--Narasimhan polygon of $\cF/\cF_{\rm tors}$ raised by $\deg\cF_{\rm tors}$.

\begin{proposition}[Theorem~A for coherent sheaves]\label{prop:cohA}
Let $\cF$ be a coherent sheaf on $\PP^n$ and $\cF'\subset\cF$ a coherent subsheaf.
\begin{enumerate}[label=\textup{(\alph*)}]
\item $\deg\cF'-\mu(z)\rk\cF'\le\phi_z(\gamma(\cF))$ for every root sequence $z$. Hence $\deg\cF'\le\Psi_\cF(\rk\cF')$, i.e.\ $P_\cF\le\Psi_\cF$.
\item $\Lambda_\cF(\lambda)\ge\max(\deg\cF_{\rm tors},\deg\cF-\lambda\rk\cF)$ for all $\lambda\in\frac1n\ZZ$, and $\Lambda_\cF(\lambda)=\deg\cF-\lambda\rk\cF$ for $\lambda\ll0$. So $\Psi_\cF$ is concave, $\Psi_\cF(0)\ge\deg\cF_{\rm tors}$ and $\Psi_\cF(\rk\cF)=\deg\cF$.
\end{enumerate}
\end{proposition}

\begin{proof}
(a) Let $\Tt$ be the sheaf of Lemma~\ref{lem:testsheaf}. By the Kleiman--Bertini theorem for tensor products of Miller and Speyer \cite{MiSp08}, which holds in any characteristic, $\operatorname{Tor}_i(\mathcal G,g_*\Tt)=0$ for $i>0$ and $\mathcal G=\cF',\cF,\cF/\cF'$, for a general $g\in\mathrm{PGL}_{n+1}(K)$; and $g_*\Tt$ has the same resolution as $\Tt$. So $\mathcal G$ tensored with this resolution is a resolution of $\mathcal G\otimes g_*\Tt$, and $\cF'\otimes g_*\Tt\to\cF\otimes g_*\Tt$ is injective. With $g_*\Tt$ in place of $\Tt$, and without saturating $\cF'$, the proofs of Lemma~\ref{lem:truncation} and Theorem~\ref{thm:main}(a) apply verbatim, with Riemann--Roch in the form of Section~\ref{sec:tables}, where the rank may be zero. This gives the first inequality, and the second follows by taking the infimum over $z$ and then over $\lambda$. For torsion-free $\cF$ the original proofs apply as they stand, with a general line $L$ that is also disjoint from the locus where $\cF$ is not locally free, and \cite{MiSp08} is not needed.

(b) The lower bound is (a) for $\cF'=\cF_{\rm tors}$ and for $\cF'=\cF$. If $z_n$ is at least the regularity of $\cF$, the higher cohomology of $\cF$ vanishes at the twists $z_r$, and the Riemann--Roch computation in the proof of Theorem~\ref{thm:main} gives $\phi_z(\gamma(\cF))=\deg\cF-\mu(z)\rk\cF$; this gives the equality for $\lambda\ll0$. The function $\Psi_\cF$ is concave as an infimum of affine functions, and the values at $0$ and $\rk\cF$ follow.
\end{proof}

\begin{example}\label{ex:Ip}
Let $\cI_p\subset\OO_{\PP^2}$ be the ideal sheaf of a point. In the rows $0$ and $1$ of its table, $h^0(\cI_p(t))=\binom{t+2}2-1$ for $t\ge0$ and $h^1(\cI_p(t))=1$ for $t<0$, and the other entries vanish. So for $a>b$ with $a+b=-2\lambda-3$, $\phi_{(a,b)}(\gamma(\cI_p))$ equals $-\lambda$ if $b\ge0$, $\frac{(a+1)(a+2)}{2(a-b)}$ if $b<0\le a$, and $\frac1{a-b}$ if $a<0$. Minimizing, with the lower bound $-\lambda$ of Proposition~\ref{prop:cohA}(b),
\[
\Lambda_{\cI_p}(\lambda)=\begin{cases}-\lambda, & \lambda\le-\frac12,\\[0.5ex] \dfrac1{2\lambda+3}, & \lambda\ge-\frac12,\end{cases}\qquad \Psi_{\cI_p}(m)=\inf_{x\in\ZZ,\,x\ge2}\Big(\frac{x-3}2\,m+\frac1x\Big).
\]
So $\Psi_{\cI_p}$ has infinitely many vertices $\big(\frac2{x(x+1)},\frac{2(x-1)}{x(x+1)}\big)$, $x\ge1$, accumulating at $(0,0)$, with an edge of slope $\frac{x-3}2$ between the vertices for $x$ and $x-1$, for $x\ge2$; $\Psi_{\cI_p}(0)=0$ and $\Psi_{\cI_p}(m)=\sqrt{2m}+O(m)$ as $m\to0$, while $P_{\cI_p}=0$.
\end{example}

By \cite[Thm.~0.2]{ES10}, $\gamma(\cF)=\sum_{y\in Y}q_y\gamma^y$, entrywise convergent, for a unique chain $Y$ and unique $q_y>0$. Here $y=(y_1>\dots>y_s)$, $0\le s\le n$, is completed by $y_i=-\infty$ for $i>s$, and $\gamma^y$ is the table of a supernatural sheaf of dimension $s$ with Hilbert polynomial $\prod_{i\le s}(t-y_i)$; such sheaves exist for every $y$ \cite{ES10}. We call $Y$ the chain of $\cF$, and we write $s=s(y)$.

\begin{proposition}[the series of Eisenbud and Schreyer]\label{prop:cohseries}
Let $\cF$ be a coherent sheaf on $\PP^n$ with chain $Y$, and put $\tau_\cF=(n-1)!\sum_{y:\,s=n-1}q_y$.
\begin{enumerate}[label=\textup{(\alph*)}]
\item For every $\lambda\in\frac1n\ZZ$ there is a root sequence $z$ with $\mu(z)=\lambda$ comparable with every element of $Y$, and for every such $z$
\[
\Lambda_\cF(\lambda)=\phi_z(\gamma(\cF))=\sum_{y:\,s=n}n!\,q_y\big(\mu(y)-\lambda\big)_++\tau_\cF .
\]
\item $\sum_{s=n}n!\,q_y=\rk\cF$ and $\tau_\cF+\sum_{s=n}n!\,q_y\mu(y)=\deg\cF$. The graph of $\Psi_\cF$ starts at $(0,\tau_\cF)$, where $\tau_\cF=\Psi_\cF(0)\ge\deg\cF_{\rm tors}$, and has edges of slopes $\mu(y)$ and horizontal lengths $n!\,q_y$, $s=n$, in decreasing order of slope, possibly with infinitely many edges accumulating at the left endpoint.
\item If $\cF$ is torsion-free, only $s=n$ occurs and $\tau_\cF=0$. The chain $Y$ is finite with $s=n$ throughout if and only if $\cF$ is locally free.
\end{enumerate}
\end{proposition}

\begin{proof}
(a) Since $\phi_z$ involves finitely many entries, $\phi_z(\gamma(\cF))=\sum_yq_y\phi_z(\gamma^y)$. Let $s<n$. If $y\le z$, the computation in the proof of Lemma~\ref{lem:comparable}, with \eqref{eq:divdiff} applied to the polynomial $\prod_{i\le s}(t-y_i)$ of degree $s$, gives $\phi_z(\gamma^y)=(n-1)!$ for $s=n-1$ and $\phi_z(\gamma^y)=0$ for $s\le n-2$. For arbitrary $z$, Proposition~\ref{prop:cohA}(a) applied to a supernatural sheaf $\mathcal S$ of dimension $s$, with $\cF'=\mathcal S$ if $s=n-1$ and $\cF'=0$ if $s\le n-2$, gives $\phi_z(\gamma^y)\ge(n-1)!$, resp.\ $\phi_z(\gamma^y)\ge0$. The elements of $Y$ with $s=n$ lie above the others, and $Y$ is bounded above, since the rows $j\ge1$ of $\gamma(\cF)$ vanish in large degrees; so only finitely many elements of $Y$ with $s=n$ have center $\ge c_\lambda$. Hence, as in Lemma~\ref{lem:adapted}, there is a root sequence $z$ with $\mu(z)=\lambda$ comparable with every element of $Y$: let $x$ be the smallest element of $Y$ with $s=n$ and $c(x)\ge c_\lambda$. If $x$ is followed in $Y$ by an element $x'$ with $s=n$, take $z$ in $[x',x]$ as in Lemma~\ref{lem:adapted}; if $x$ exists and is not followed by such an element, take $z=x-n(c(x)-c_\lambda)\epsilon_n$; and if there is no such $x$, take $z$ above the largest element of $Y$. For a comparable $z$ the displayed value of $\phi_z(\gamma(\cF))$ follows from Lemma~\ref{lem:comparable} and the values above, and by Lemma~\ref{lem:positivity} and the lower bounds above it is a lower bound for $\phi_{z'}(\gamma(\cF))$ for every $z'$ with $\mu(z')=\lambda$.

(b) The sum in (a) converges, since $\Lambda_\cF(\lambda)<\infty$. The slopes $\mu(y)$, $s=n$, are bounded below. Choose $\lambda_0\in\frac1n\ZZ$ below all of them and in the range where $\Lambda_\cF(\lambda)=\deg\cF-\lambda\rk\cF$ (Proposition~\ref{prop:cohA}(b)). Subtracting the formulas at $\lambda_0$ and $\lambda_0-1/n$ gives $\sum_{s=n}n!\,q_y=\rk\cF$, and substitution gives $\tau_\cF+\sum_{s=n}n!\,q_y\mu(y)=\deg\cF$. These series have finite absolute first moment. For $\lambda\ge\lambda_0$, the summands $n!\,q_y(\mu(y)-\lambda)_+$ are dominated by the summable terms $n!\,q_y(\mu(y)-\lambda_0)$; hence their sum tends to $0$ as $\lambda\to\infty$, and $\Psi_\cF(0)=\inf_\lambda\Lambda_\cF(\lambda)=\tau_\cF$, which is at least $\deg\cF_{\rm tors}$ by Proposition~\ref{prop:cohA}(b). The description of the graph follows from (a) by the duality of Lemma~\ref{lem:legendre}(i). One may adjoin the vertical segment from $(0,0)$ to $(0,\tau_\cF)$ to obtain an augmented polygon.

(c) If $\cF$ is torsion-free, only $s=n$ occurs and the $q_y$ are rational \cite[Thm.~0.2]{ES10}. If $Y$ is finite with $s=n$ throughout, then $H^j(\cF(t))=0$ for $j<n$ and $t\ll0$, and then $\mathcal{E}xt^i(\cF,\OO)=0$ for $i>0$: Serre duality gives $\operatorname{Ext}^i(\cF,\OO(k))=0$ for $k\gg0$, and the local-to-global Ext spectral sequence, together with Serre vanishing, identifies these groups with $H^0(\mathcal{E}xt^i(\cF,\OO)(k))$. Vanishing for all large $k$ forces the Ext sheaves to vanish, so $\cF$ is locally free. The converse is Theorem~\ref{thm:ES}.
\end{proof}

By Example~\ref{ex:Ip} and Proposition~\ref{prop:cohseries}(b), the chain of $\cI_p$ is infinite.

\begin{corollary}\label{cor:ESquestion}
Let $\cF$ be a coherent sheaf on $\PP^n$, $n\ge1$, such that all the tables in the decomposition $\gamma(\cF)=\sum_{y\in Y}q_y\gamma^y$ of \cite[Thm.~0.2]{ES10} are tables of vector bundles on $\PP^n$, i.e.\ $s=n$ for every $y\in Y$ (the criterion of \cite{ES10} for $\gamma^y$ to be the table of a vector bundle). Then $\cF$ has no associated subvarieties of dimension $n-1$ or $0$. In particular, for $n\le2$ the sheaf $\cF$ is torsion-free.
\end{corollary}

\begin{proof}
By Proposition~\ref{prop:cohseries}(b), $\tau_\cF=0$ and $\tau_\cF\ge\deg\cF_{\rm tors}\ge0$; so $\deg\cF_{\rm tors}=0$, i.e.\ $\cF_{\rm tors}$ has no component of dimension $n-1$. Let $\cF_0\subset\cF$ be the largest subsheaf of finite length, and $\ell$ its length. Then $h^0(\cF_0(t))=\ell$ and $H^i(\cF_0(t))=0$ for $i>0$ and all $t$, so $\gamma(\cF)=\gamma(\cF/\cF_0)+\ell\sum_te_{0,t}$, and $\sum_te_{0,t}$ is the supernatural table with $s=0$, the table of a point. Its sequence lies below every other one, so adding it to the chain of $\cF/\cF_0$ gives a chain, which for $\ell>0$ is the chain of $\cF$ by the uniqueness in \cite[Thm.~0.2]{ES10}. By hypothesis the chain of $\cF$ has no element with $s=0$, so $\ell=0$. Finally, the associated subvarieties of $\cF$ of dimension less than $n$ are those of $\cF_{\rm tors}$, and for $n\le2$ they have dimension $n-1$ or $0$.
\end{proof}

Eisenbud and Schreyer observe that the decomposition of a torsion-free sheaf involves only tables of vector bundles, and ask whether the converse is true \cite[after Thm.~0.2]{ES10}. Corollary~\ref{cor:ESquestion} shows that it is for $n\le2$; for $n\ge3$ it leaves open the associated subvarieties of dimension $1,\dots,n-2$.

\begin{proposition}[Theorem~C for coherent sheaves]\label{prop:cohC}
\begin{enumerate}[label=\textup{(\alph*)}]
\item For every exact sequence $0\to\cF'\to\cF\to\cF''\to0$ of coherent sheaves on $\PP^n$, $\Lambda_\cF\le\Lambda_{\cF'\oplus\cF''}$ and $\Psi_\cF\le\Psi_{\cF'\oplus\cF''}$. Hence $\Psi_\cF\le\Psi_{\gr\cF}$ for every filtration of $\cF$ by coherent subsheaves, for instance for the Harder--Narasimhan filtration of a vector bundle, whose factors need not be locally free.
\item If $n\ge2$ and $\cF\subset\cE$ are coherent sheaves with $\cE/\cF$ of finite length, then $\Psi_\cE\le\Psi_\cF$.
\item Let $(\cF_s)_{s\in S}$ be a flat family of coherent sheaves on $\PP^n$ over a variety $S$, $s_0\in S$ and $m_0>0$. There is a neighborhood $U$ of $s_0$ with $\Psi_{\cF_s}\le\Psi_{\cF_{s_0}}$ on $[m_0,\rk\cF_{s_0}]$ for all $s\in U$.
\end{enumerate}
\end{proposition}

\begin{proof}
(a) The identity \eqref{eq:extdiff} holds for every exact sequence of coherent sheaves, and $\phi_z$ is nonnegative on cancellation vectors (Lemma~\ref{lem:localvalues}). So $\phi_z(\gamma(\cF'\oplus\cF''))\ge\phi_z(\gamma(\cF))$ for every $z$, and the inequalities follow by taking infima. The same holds along a filtration, as in the proof of Corollary~\ref{cor:monoext}(a).

(b) Let $\ell$ be the length of $\cE/\cF$. The table of $\cE/\cF$ is $\ell\sum_te_{0,t}$, on which every $\phi_z$ vanishes by \eqref{eq:divdiff}, so $\phi_z(\gamma(\cF))\ge\phi_z(\gamma(\cE))$ by \eqref{eq:extdiff}; and $\cE$, $\cF$ have the same rank.

(c) By Proposition~\ref{prop:cohseries}, for $m\ge m_0$ the infimum defining $\Psi_{\cF_{s_0}}(m)$ is attained at one of the finitely many slopes of the edges of $\Psi_{\cF_{s_0}}$ that meet $[m_0,\rk\cF_{s_0}]$. The argument of Corollary~\ref{cor:monoext}(b), applied to the columns $z_1,\dots,z_n$ of root sequences $z$ comparable with the chain of $\cF_{s_0}$ at these slopes, gives a neighborhood $U$ of $s_0$ on which $\phi_z(\gamma(\cF_s))\le\phi_z(\gamma(\cF_{s_0}))$ for these $z$; so $\Psi_{\cF_s}(m)\le m\lambda+\phi_z(\gamma(\cF_s))\le m\lambda+\Lambda_{\cF_{s_0}}(\lambda)=\Psi_{\cF_{s_0}}(m)$ for $m\ge m_0$ and the corresponding slope $\lambda$.
\end{proof}

On $\PP^2$, (b) applies to a torsion-free sheaf $\cF$ and $\cE=\cF^{\vee\vee}$, a vector bundle with the same Harder--Narasimhan polygon; so there Theorem~\ref{thm:main} for $\cF^{\vee\vee}$ bounds the subsheaves of $\cF$ at least as well as $\Psi_\cF$ does. In (c) we do not know whether $U$ can be chosen independently of $m_0$; the tables may change in infinitely many columns, as for the extensions $0\to\OO(-1)\to\cF_s\to\cI_p\to0$ on $\PP^2$, $s\in\Ext^1(\cI_p,\OO(-1))\cong K$, where $\cF_s\cong T_{\PP^2}(-2)$ and $h^1(\cF_s(t))=0$ for $s\ne0$ and $t\le-2$, while $h^1(\cF_0(t))=1$ for all $t<0$.

\begin{proposition}[Theorem~D for coherent sheaves]\label{prop:cohD}
Let $\cF$ be a coherent sheaf on $\PP^2$, $\lambda\in\frac12\ZZ$, and $a>b$ integers with $a+b=-2\lambda-3$. Then $\Lambda_\cF(\lambda)\le\phi_{a,b}(\gamma(\cF))$, with equality if $(a,b)$ is comparable with every element of the chain of $\cF$. For every coherent subsheaf $\cF'\subset\cF$, with $\cF''=\cF/\cF'$ and the notation of Proposition~\ref{prop:fourdefect},
\begin{multline*}
(a-b)\big(\phi_{a,b}(\gamma(\cF))-\deg\cF'+\lambda\rk\cF'\big)\\=\big(h^2_b(\cF')-h^2_a(\cF')\big)+\big(h^0_a(\cF'')-h^0_b(\cF'')\big)+\big(h^1_a(\cF')-\rk\partial^0_a\big)+\big(h^1_b(\cF'')-\rk\partial^1_b\big),
\end{multline*}
and all four terms are nonnegative.
\end{proposition}

\begin{proof}
The inequality is the definition of $\Lambda_\cF$, and the equality is Proposition~\ref{prop:cohseries}(a). Apart from its first sentence, the proof of Proposition~\ref{prop:fourdefect} below uses only the definition of $\phi_{a,b}$, and it applies to $\cF'\subset\cF$, with $\deg\cF'-\lambda\rk\cF'$ in place of $n'(\mu'-\lambda)$: Riemann--Roch is used in the form of Section~\ref{sec:tables}; $H^2(\cF'(t))\to H^2(\cF'(t+1))$ is surjective for every nonzero linear form, since $H^2$ is right exact on $\PP^2$ and the cokernel of $\cF'(t)\to\cF'(t+1)$ is supported on a line; and a general linear form is a nonzerodivisor on $\cF''$.
\end{proof}

In particular Proposition~\ref{prop:cohD} gives Proposition~\ref{prop:cohA}(a) on $\PP^2$ without \cite{MiSp08}.

\section{The plane}\label{sec:plane}

In this section and the next, $n=2$. A root sequence is a pair $z=(z_1>z_2)$, with center $c(z)=\frac12(z_1+z_2)$, gap $g(z)=z_1-z_2\ge1$ and slope $\mu(z)=-c(z)-\frac32$, and
\[
\gamma^z_{j,t}=\begin{cases}(t-z_1)(t-z_2) & j=0,\ t>z_1,\\ (z_1-t)(t-z_2) & j=1,\ z_2<t<z_1,\\ (t-z_1)(t-z_2) & j=2,\ t<z_2,\\ 0 & \text{otherwise.}\end{cases}
\]
Its Hilbert polynomial is $(t-z_1)(t-z_2)=(t-c)^2-\frac{g^2}4$, and on the $h^1$ row $\gamma^z_{1,t}=\big(\frac{g^2}4-(t-c)^2\big)_+$ for all $t\in\ZZ$. Ranks of pieces are $2q_k$, and Boij--Söderberg slopes lie in $\frac12\ZZ$.

\subsection{Gap moment and centers}
For a vector bundle $F$ of rank $N$ on $\PP^2$ with chain decomposition $\sum_kq_k\gamma^{z^k}$ put $p_k=2q_k/N$, $c_k=c(z^k)$, $g_k=g(z^k)$, and
\[
\Gamma(F)=\sum_kq_k\,g_k^2,\qquad G(F)=\sum_kp_k\,g_k^2=\tfrac2N\Gamma(F),\qquad V_c(F)=\sum_kp_k(c_k-\bar c)^2,\quad\bar c=\sum_kp_kc_k .
\]
For a torsion-free sheaf $\cF$ of rank $N$ on $\PP^2$ write $\mu(\cF)=c_1/N$ and $\delta(\cF)=\Delta(\cF)/N^2$ with $\Delta=2Nc_2-(N-1)c_1^2$. Riemann--Roch reads
\begin{equation}\label{eq:RRP2}
\frac{\chi(\cF(t))}{N}=\frac12\Big(t+\mu+\frac32\Big)^2-\frac18-\frac{\delta}{2}.
\end{equation}

\begin{lemma}\label{lem:basic}
Let $F$ be a vector bundle of rank $N$ on $\PP^2$.
\begin{enumerate}[label=\textup{(\roman*)}]
\item $\chi(F(t))/N=\frac12\sum_kp_k(t-c_k)^2-G(F)/8$.
\item $\bar c=-\mu(F)-\frac32$, i.e.\ $\deg\gamma(F)=\deg F$, and $G(F)-1=4V_c(F)+4\delta(F)$.
\item $G(F)\ge1$, with equality iff $F$ splits.
\item $\sum_th^1(F(t))=\sum_kq_k\frac{g_k^3-g_k}6=\frac N{12}\sum_kp_k(g_k^3-g_k)$, and this is at least $\frac N6\big(G(F)-1\big)$ and at least $\frac N{12}\sqrt{G(F)}\,\big(G(F)-1\big)$.
\item If $0=F_0\subset\dots\subset F_s=F$ is the Harder--Narasimhan filtration, with factors $\cQ_i$ of ranks $n_i$, weights $w_i=n_i/N$ and slopes $\mu_i$, then
\[
G(F)-1=4\big(V_c(F)-V_{HN}(F)\big)+4\sum_iw_i\,\delta(\cQ_i),\qquad V_{HN}(F)=\sum_iw_i(\mu_i-\mu(F))^2 .
\]
\end{enumerate}
\end{lemma}

\begin{proof}
(i) The Hilbert polynomial is linear in the table, and $N=2\sum_kq_k$.
(ii) Compare (i) with \eqref{eq:RRP2}: the coefficients of $t$ give $\bar c=-\mu-\frac32$, and the constant terms give $\frac12V_c-\frac G8=-\frac18-\frac\delta2$.
(iii) The $h^1$ row of $\gamma^z$ vanishes iff $g(z)=1$, and a bundle on $\PP^2$ without intermediate cohomology splits by Horrocks' criterion \cite{Ho64}.
(iv) $\sum_{u=1}^{g-1}u(g-u)=\frac{g^3-g}6$, and $\frac{g^3-g}6\ge\frac{g^2-1}3$ is $(g^2-1)(g-2)\ge0$ for integers $g\ge1$. For the second bound, $g^3-g=h(g^2)$ with $h(x)=x^{3/2}-x^{1/2}$, which is convex on $[0,\infty)$; Jensen's inequality gives $\sum_kp_kh(g_k^2)\ge h(G(F))$.
(v) Sum \eqref{eq:RRP2} over the $\cQ_i$ and compare constant terms with (i), using (ii).
\end{proof}

By (v), the inequality $V_{HN}(F)\le V_c(F)$ of Corollary~\ref{cor:supernatural}(c) is equivalent to $G(F)\ge1+4\sum_iw_i\delta(\cQ_i)$. By (ii) and Corollary~\ref{cor:monoext}(a), $G(F)\le G(\gr F)$ for every filtration with locally free subquotients, since $\delta$ depends only on the Hilbert polynomial. The constant term of the Hilbert polynomial, $\sum_kq_k(c_k^2-g_k^2/4)$, is linear in the table, so $\Gamma=4\sum_kq_kc_k^2$ minus a linear function; by Theorem~\ref{thm:concave}, $\Gamma$ is concave on $\CC$, and $\sum_jq'_jg(y_j)^2\le\Gamma(\gamma)$ for every decomposition $\gamma=\sum_jq'_j\gamma^{y_j}$.

\begin{lemma}[two twists]\label{lem:twotwist}
For every vector bundle $F$ of rank $N$ on $\PP^2$ and every $t\in\ZZ$,
\[
G(F)-1\ \ge\ \frac4N\Big(h^1(F(t))+h^1(F(t+1))\Big).
\]
\end{lemma}

\begin{proof}
Since $G-1=\frac2N\sum_kq_k(g_k^2-1)$, it suffices to show $g^2-1\ge2(\gamma^z_{1,t}+\gamma^z_{1,t+1})$ for every $z$. If $t,t+1\in(z_2,z_1)$, put $a=z_1-t$, $b=t-z_2$; then $g^2-1-2[ab+(a-1)(b+1)]=(a-b-1)^2$. If exactly one of $t,t+1$ lies in $(z_2,z_1)$, it is $z_2+1$ or $z_1-1$, the right-hand side is $2(g-1)$, and $g^2-1-2(g-1)=(g-1)^2$.
\end{proof}

\subsection{The extreme slopes}
For a nonzero table $\gamma\in\CC$ on $\PP^2$ put
\[
M_j=\max\{t:\gamma_{j,t}\ne0\}\ \ (j=1,2),\qquad m_j=\min\{t:\gamma_{j,t}\ne0\}\ \ (j=0,1),
\]
with $M_1=-\infty$ and $m_1=+\infty$ if the $h^1$ row vanishes.

\begin{corollary}\label{cor:endsP2}
For $\gamma\in\CC$ on $\PP^2$, $\gamma\ne0$, the first and the last element of the chain are
\[
y^{\max}=\big(\max(M_1+1,\,M_2+2),\ M_2+1\big),\qquad y^{\min}=\big(m_0-1,\ \min(m_1-1,\,m_0-2)\big),
\]
and the largest and smallest Boij--Söderberg slopes of $\gamma$ are
\[
\mu^\BS_{\max}(\gamma)=\max\Big(-m_0,\ -\frac{m_0+m_1+1}2\Big),\qquad
\mu^\BS_{\min}(\gamma)=\min\Big(-M_2-3,\ -\frac{M_1+M_2+5}2\Big).
\]
\end{corollary}

\begin{proof}
This is Lemma~\ref{lem:ends} for $n=2$; the slopes are $-c-\frac32$.
\end{proof}

By Corollary~\ref{cor:supernatural}(b), $\mu_{\max}(F)\le\mu^\BS_{\max}(F)$. At the first step of the Harder--Narasimhan filtration one can say more: the chain of $F$ reaches at least as far as that of the subsheaf.

\begin{proposition}\label{prop:firstedge}
Let $F$ be a vector bundle on $\PP^2$ and $F_1\subset F$ a saturated subsheaf with $\mu_{\max}(F/F_1)<\mu(F_1)$, for instance the maximal destabilizing subsheaf. Then
\[
\mu^\BS_{\max}(F)\ \ge\ \mu^\BS_{\max}(F_1)\ \ge\ \mu(F_1).
\]
\end{proposition}

\begin{proof}
$F_1$ is reflexive, hence locally free. Put $F_2=F/F_1$, $\mu_1=\mu(F_1)$, $\bar c_1=-\mu_1-\frac32$, and use the invariants $m_0,m_1$ of Corollary~\ref{cor:endsP2}. First, $m_0(F)\le m_0(F_1)$, because $H^0(F_1(t))\subset H^0(F(t))$. Second, for $t<-\mu_{\max}(F_2)$ we have $H^0(F_2(t))=0$, so $H^1(F_1(t))\to H^1(F(t))$ is injective.

If $m_1(F_1)-1\le m_0(F_1)-2$, the last chain element of $F_1$ is $(m_0(F_1)-1,m_1(F_1)-1)$, so $m_1(F_1)\le c(z^m(F_1))+\frac12$. The last center is the smallest, hence at most the mean center $\bar c_1$ (Lemma~\ref{lem:basic}(ii)); so $m_1(F_1)\le \bar c_1+\frac12=-\mu_1-1<-\mu_{\max}(F_2)$. Hence $m_1(F)\le m_1(F_1)$, and by Corollary~\ref{cor:endsP2}
\[
c(z^m(F))\le\tfrac12\big(m_0(F)+m_1(F)-2\big)\le\tfrac12\big(m_0(F_1)+m_1(F_1)-2\big)=c(z^m(F_1)).
\]
Otherwise $c(z^m(F_1))=m_0(F_1)-\frac32\ge m_0(F)-\frac32\ge c(z^m(F))$. In both cases $\mu^\BS_{\max}(F)\ge\mu^\BS_{\max}(F_1)$. The second inequality holds because the rank-weighted mean of the Boij--Söderberg slopes of $F_1$ is $\mu_1$.
\end{proof}

The first inequality compares the ends of the chains of $F$ and $F_1$; Theorem~\ref{thm:HNmono} below compares the whole polygons when $F_1$ is a step of the Harder--Narasimhan filtration.

\begin{remark}[no monotonicity for arbitrary subbundles]\label{rem:nomono}
The slope hypothesis cannot be dropped. Let $G=\ker(\OO^3\to\OO(2))$ for a regular sequence of three conics $q_1,q_2,q_3$ (for instance three general conics); then $H^1(G(t))$ is the degree $t+2$ part of $K[x_0,x_1,x_2]/(q_1,q_2,q_3)$, of dimensions $1,3,3,1$ for $t=-2,\dots,1$, and $H^0(G(1))=0$, since there are no linear syzygies. So $G$ is stable of rank $2$ and slope $-1$, with chain $\frac13\gamma^{(2,-2)}+\frac13\gamma^{(1,-2)}+\frac13\gamma^{(1,-3)}$. The subbundle $F_1=\OO\oplus G\subset\OO^4$ has $m_0(F_1)=0$ and $m_1(F_1)=-2$, so by Corollary~\ref{cor:endsP2} its chain ends with $(-1,-3)$; indeed it is
\[
\tfrac13\gamma^{(2,-2)}+\tfrac7{12}\gamma^{(1,-2)}+\tfrac16\gamma^{(1,-3)}+\tfrac1{12}\gamma^{(0,-3)}+\tfrac13\gamma^{(-1,-3)},
\]
with largest slope $\frac12$, while all Boij--Söderberg slopes of $\OO^4$ are $0$. So $\Psi_{F_1}(1)=\frac14>0=\Psi_{\OO^4}(1)$ and $\Lambda_{F_1}(0)=\frac13>0=\Lambda_{\OO^4}(0)$. Here $\mu(F_1)=-\frac23<2=\mu(\OO^4/F_1)$. Moreover $\gamma(\OO)=\frac12\gamma^{(-1,-2)}$, and $(-1,-2)$ is incomparable with $(1,-3)$; so the union of the chains of $\OO$ and $G$ is not a chain, and the re-chaining creates the root sequences $(0,-3)$ and $(-1,-3)$, which occur in neither.
\end{remark}

\subsection{The three-term formula}
For integers $a>b$ and a table $\gamma$ on $\PP^2$ put
\[
\phi_{a,b}(\gamma)=\phi_{(a,b)}(\gamma)=\frac{\gamma_{0,a}-\gamma_{0,b}+\gamma_{1,b}}{a-b}.
\]
With the Hilbert polynomial $\chi_\gamma(t)=\sum_j(-1)^j\gamma_{j,t}$, also $\phi_{a,b}(\gamma)=\big(\gamma_{0,a}+\gamma_{2,b}-\chi_\gamma(b)\big)/(a-b)$.

\begin{theorem}[three-term formula]\label{thm:three}
Let $\gamma\in\CC$ be a table on $\PP^2$ with chain decomposition $\sum_{k=1}^mq_k\gamma^{z^k}$, let $\lambda\in\frac12\ZZ$, and put $c=c_\lambda$.
\begin{enumerate}[label=\textup{(\alph*)}]
\item If $a>b$ are integers with $a+b=2c$ and $y$ is a root sequence comparable with $(a,b)$, then $\phi_{a,b}(\gamma^y)=2\big(c-c(y)\big)_+$; for arbitrary $y$, $\phi_{a,b}(\gamma^y)\ge2\big(c-c(y)\big)_+$.
\item There are integers $a>b$ with $a+b=2c$ such that $(a,b)$ is comparable with every $z^k$, and for each of them $\Lambda_\gamma(\lambda)=\phi_{a,b}(\gamma)$.
\item $\Lambda_\gamma(\lambda)=\min\{\phi_{a,b}(\gamma):\ a>b,\ a+b=2c\}$.
\end{enumerate}
Consequently, for a vector bundle $F$ on $\PP^2$ and $0\le m\le\rk F$,
\[
\Psi_F(m)=\min_{a>b}\Big(\frac{h^0(F(a))-h^0(F(b))+h^1(F(b))}{a-b}-m\Big(\frac{a+b}2+\frac32\Big)\Big).
\]
\end{theorem}

\begin{proof}
(a) For comparable $y$ this is Lemma~\ref{lem:comparable}. Otherwise put $P(u)=(u-y_1)(u-y_2)$, so that $P(a)-P(b)=2(a-b)(c-c(y))$, and recall $\phi_{a,b}(\gamma)=(\gamma_{0,a}-\gamma_{0,b}+\gamma_{1,b})/(a-b)$. If $y_1>a$ and $y_2<b$, both columns $a,b$ lie in row $1$, where the entries are $-P$, so $\phi_{a,b}(\gamma^y)=\gamma^y_{1,b}/(a-b)$, while $\gamma^y_{1,a}-\gamma^y_{1,b}=2(a-b)(c(y)-c)$. If $y_1<a$ and $y_2>b$, then $\phi_{a,b}(\gamma^y)=\gamma^y_{0,a}/(a-b)$, while $\gamma^y_{0,a}-\gamma^y_{2,b}=2(a-b)(c-c(y))$. In both cases
\begin{equation}\label{eq:incomparable}
\phi_{a,b}(\gamma^y)=2\big(c-c(y)\big)_++\frac{\min\big(\gamma^y_{\bullet,a},\gamma^y_{\bullet,b}\big)}{a-b}\ \ge\ 2\big(c-c(y)\big)_+,
\end{equation}
where $\gamma^y_{\bullet,u}=|P(u)|$ is the entry of $\gamma^y$ in the column $u$. (b) is Lemma~\ref{lem:adapted} for $n=2$. (c) By (a), $\phi_{a,b}(\gamma)=\sum_kq_k\phi_{a,b}(\gamma^{z^k})\ge\Lambda_\gamma(\lambda)$ for every pair, with equality in (b). The last formula follows with Lemma~\ref{lem:legendre}(i). This proof uses only Lemmas~\ref{lem:comparable}--\ref{lem:adapted}, not Theorem~\ref{thm:main}.
\end{proof}

We call a pair $(a,b)$ as in (b) \emph{adapted} to the chain of $\gamma$ (or of $F$). The values of $\phi_{a,b}$ on the cancellation vectors are $\phi_{a,b}(e_{0,a}+e_{1,a})=\phi_{a,b}(e_{1,b}+e_{2,b})=\frac1{a-b}$, by Lemma~\ref{lem:localvalues}.

On incomparable root sequences, $\phi_{a,b}$ measures how much the chain gains over other decompositions. For a root sequence $y$ and $u\in\ZZ$ let $\gamma^y_{\bullet,u}=|(u-y_1)(u-y_2)|$ be the entry of $\gamma^y$ in the column $u$, and for a root sequence $z=(s,t)$ put
\[
E_y(z)=\begin{cases}0 & \text{if $z$ and $y$ are comparable},\\[1ex] \dfrac{\min\big(\gamma^y_{\bullet,s},\gamma^y_{\bullet,t}\big)}{s-t} & \text{otherwise}.\end{cases}
\]

\begin{corollary}[the gain formula]\label{cor:gain}
Let $\gamma=\sum_yq_y\gamma^y$, $q_y\ge0$, be any decomposition of a table $\gamma\in\CC$ on $\PP^2$, and let $c=c_\lambda\in\frac12\ZZ$. Then
\[
\Lambda_\gamma(\lambda)=\sum_y2q_y\big(c-c(y)\big)_++\min_{z:\ c(z)=c}\ \sum_yq_y\,E_y(z),
\]
the minimum over all root sequences $z$ with center $c$. It is attained at the root sequences with center $c$ that are comparable with the chain of $\gamma$.
\end{corollary}

\begin{proof}
By Theorem~\ref{thm:three}, $\Lambda_\gamma(\lambda)=\min_z\sum_yq_y\phi_z(\gamma^y)$, with the minimum attained as stated, and $\phi_z(\gamma^y)=2(c-c(y))_++E_y(z)$ for $z$ with $c(z)=c$, by Theorem~\ref{thm:three}(a) for comparable $y$ and $z$ and by \eqref{eq:incomparable} otherwise.
\end{proof}

So the chain improves on any decomposition by the \emph{gain} $\min_z\sum_yq_yE_y(z)$, which is positive exactly when no root sequence with center $c$ is comparable with all the $y$ with $q_y>0$. This quantifies Theorem~\ref{thm:concave}(a) on $\PP^2$.

\subsection{The four defects}
On $\PP^2$ the defect of Proposition~\ref{prop:marginPn} is a sum of nonnegative terms.

\begin{proposition}[four defects]\label{prop:fourdefect}
Let $F$ be a vector bundle on $\PP^2$, $F'\subset F$ a nonzero subsheaf of rank $n'$ and slope $\mu'$ with torsion-free quotient $F''=F/F'$, $\lambda\in\frac12\ZZ$, and $(a,b)$ a pair adapted to the chain of $F$ with $a+b=-2\lambda-3$. Write $h^j_t(G)=h^j(G(t))$, and let $\partial^0_a\colon H^0(F''(a))\to H^1(F'(a))$ and $\partial^1_b\colon H^1(F''(b))\to H^2(F'(b))$ be the connecting homomorphisms. Then
\begin{multline*}
(a-b)\big(\Lambda_F(\lambda)-n'(\mu'-\lambda)\big)\\=\big(h^2_b(F')-h^2_a(F')\big)+\big(h^0_a(F'')-h^0_b(F'')\big)+\big(h^1_a(F')-\rk\partial^0_a\big)+\big(h^1_b(F'')-\rk\partial^1_b\big),
\end{multline*}
and all four terms are nonnegative.
\end{proposition}

\begin{proof}
By Theorem~\ref{thm:three}(b), $(a-b)\Lambda_F(\lambda)=h^0_a(F)-h^0_b(F)+h^1_b(F)$. The long exact cohomology sequences give $h^0_a(F)=h^0_a(F')+h^0_a(F'')-\rk\partial^0_a$ and $h^1_b(F)-h^0_b(F)=h^1_b(F')-h^0_b(F')+h^1_b(F'')-h^0_b(F'')-\rk\partial^1_b$. Writing $h^0_a(F')=\chi(F'(a))+h^1_a(F')-h^2_a(F')$ and $h^1_b(F')-h^0_b(F')=h^2_b(F')-\chi(F'(b))$, we obtain the identity, since $\chi(F'(a))-\chi(F'(b))=\frac{n'}2(a-b)(a+b+2\mu'+3)=(a-b)\,n'(\mu'-\lambda)$ by \eqref{eq:RRP2}. The last two terms are nonnegative trivially. A nonzero linear form is a nonzerodivisor on the torsion-free sheaves $F'$ and $F''$; multiplication by it gives injections $H^0(F''(t))\to H^0(F''(t+1))$, and surjections $H^2(F'(t))\to H^2(F'(t+1))$, since the cokernel of $F'(t)\to F'(t+1)$ is supported on a line. So $h^0_t(F'')$ is nondecreasing and $h^2_t(F')$ is nonincreasing in $t$, and $a>b$.
\end{proof}

For $n=2$, \eqref{eq:marginPn} is this identity before the rearrangement: $B_1=h^1_a(F')-h^2_a(F')+h^0_a(F'')-\rk\partial^0_a$ and $B_2=h^2_b(F')+h^1_b(F'')-h^0_b(F'')-\rk\partial^1_b$; so $B_1$ can be negative, e.g.\ $B_1=-1$, $B_2=3$ for $0\to\OO\to\OO^2\to\OO\to0$ and $(a,b)=(-3,-4)$, $\lambda=2$. Proposition~\ref{prop:fourdefect}, Theorem~\ref{thm:three}(b) and Lemma~\ref{lem:legendre} give an elementary proof of Theorem~\ref{thm:main} on $\PP^2$; it corresponds to the test sheaf $\OO_{(a-b)L}(a)$, with resolution $\OO(b)\to\OO(a)$. They also describe the equality cases: the supporting-line equality $\Lambda_F(\lambda)=\deg F'-\lambda\rk F'$ at a step $F'$ of the Harder--Narasimhan filtration holds iff all four terms vanish for a pair $(a,b)$ adapted at $\lambda$.

\begin{corollary}\label{cor:Q}
For every vector bundle $F$ on $\PP^2$ with Harder--Narasimhan factors $\cQ_i$ of weights $w_i$,
\[
G(F)\ \ge\ 1+4\sum_iw_i\,\delta(\cQ_i),
\]
where $\delta(\cQ_i)\ge0$ by Bogomolov's inequality \cite{Bo78,La04}.
\end{corollary}

\begin{proof}
Corollary~\ref{cor:supernatural}(c) and Lemma~\ref{lem:basic}(v).
\end{proof}

The inequality $P_F\le\Psi_F$ cannot be obtained from the generic splitting type on a line. That type dominates $P_F$, but it is in general more spread out than $\Psi_F$: for $T_{\PP^2}(-1)$ it is $(1,0)$, while the Boij--Söderberg slopes are $(\frac12,\frac12)$.

\subsection{Tables of extensions}
Proposition~\ref{prop:fourdefect} works at the level of tables. This gives necessary conditions on the cohomology tables of extensions of two given bundles, one of the problems of Erman and Sam \cite[Problem~4.3]{ErS}: if $\gamma_1=\gamma(F')$ and $\gamma_2=\gamma(F'')$ are tables of vector bundles, every extension of $F''$ by $F'$ has a table of the following form, with $d^j_t=\rk\partial^j_t$.

\begin{corollary}[tables of extensions]\label{cor:conicQ}
Let $\gamma_1,\gamma_2\in\CC$ be tables on $\PP^2$, $\gamma_1\ne0$, and $D=\sum_td^0_t(e_{0,t}+e_{1,t})+\sum_td^1_t(e_{1,t}+e_{2,t})$ with $0\le d^0_t\le\gamma_{1;1,t}$ and $0\le d^1_t\le\gamma_{2;1,t}$, where $\gamma_{i;j,t}$ are the entries of $\gamma_i$. If $\gamma=\gamma_1+\gamma_2-D\in\CC$, then $\Psi_\gamma(\rk\gamma_1)\ge\deg\gamma_1$.
\end{corollary}

\begin{proof}
Since $D$ is $\chi$-neutral, $\rk\gamma=\rk\gamma_1+\rk\gamma_2$ and $\deg\gamma=\deg\gamma_1+\deg\gamma_2$. As in Lemma~\ref{lem:legendre}(i),(iii), it suffices to show $\Lambda_\gamma(\lambda)\ge\rk\gamma_1\,(\mu(\gamma_1)-\lambda)$ for $\lambda\in\frac12\ZZ$, where $\mu(\gamma_1)=\deg\gamma_1/\rk\gamma_1$. Let $c=c_\lambda$ and choose $(a,b)$ adapted to $\gamma$. Then
\[
\gamma_{0,a}=\gamma_{1;0,a}+\gamma_{2;0,a}-d^0_a\ \ge\ \gamma_{1;0,a}+\gamma_{2;0,a}-\gamma_{1;1,a},
\]
\[
\gamma_{1,b}-\gamma_{0,b}=(\gamma_{1;1,b}-\gamma_{1;0,b})+(\gamma_{2;1,b}-\gamma_{2;0,b})-d^1_b\ \ge\ \gamma_{1;1,b}-\gamma_{1;0,b}-\gamma_{2;0,b}.
\]
Every table in $\CC$ has a nondecreasing row $0$ and a nonincreasing row $2$, because each $\gamma^z$ has. As in the proof of Proposition~\ref{prop:fourdefect},
\[
(a-b)\Lambda_\gamma(\lambda)\ \ge\ \chi_{\gamma_1}(a)-\chi_{\gamma_1}(b)+(\gamma_{1;2,b}-\gamma_{1;2,a})+(\gamma_{2;0,a}-\gamma_{2;0,b})\ \ge\ \chi_{\gamma_1}(a)-\chi_{\gamma_1}(b),
\]
and by the formula for the Hilbert polynomial in Section~\ref{sec:tables}, $\chi_{\gamma_1}(a)-\chi_{\gamma_1}(b)=(a-b)\rk\gamma_1\,(\mu(\gamma_1)-\lambda)$.
\end{proof}

The tables of actual extensions also satisfy $d^0_t\le\gamma_{2;0,t}$ and $d^1_t\le\gamma_{1;2,t}$. Even with these bounds, the conditions of Corollary~\ref{cor:conicQ} are necessary but not sufficient. Let $F'=\Omega^1(1)\oplus\OO(1)$ and $F''=\OO(2)$. Then $\Ext^1(F'',F')=H^1(\Omega^1(-1))\oplus H^1(\OO(-1))=0$, so the only extension is split. The Euler sequence gives $\gamma(F')=3\gamma(\OO)+e_{0,-1}+e_{1,-1}$, and the table $\gamma(F')+\gamma(F'')-(e_{0,-1}+e_{1,-1})=3\gamma(\OO)+\gamma(\OO(2))$ satisfies all the bounds, with $d^0_{-1}=1\le\min\big(h^1(F'(-1)),h^0(F''(-1))\big)$; but it is not the table of an extension of $F''$ by $F'$.

\begin{remark}\label{rem:linear}
The chain enters through the choice of an adapted pair; a linear lower bound valid on all of $\CC$ cannot replace it. A linear functional $L$ with $L(\gamma^z)\le2(\mu(z)-\lambda)_+$ for all $z$ gives a lower bound $L(\gamma(F))\le\Lambda_F(\lambda)$. But then $L(\gamma(F))\le\sum_j2q'_j(\mu(y_j)-\lambda)_+$ for \emph{every} decomposition $\gamma(F)=\sum_jq'_j\gamma^{y_j}$; by Theorem~\ref{thm:concave}(a) the chain gives the largest of these values, and the smallest can be below the vertex target. Let $F_1=\coker(\OO^2\to\OO(1)^2\oplus\OO(2)^2)$ and $F_2=\coker(\OO(-1)\to\OO\oplus\OO(1)^2)$ be general, of slopes $3$ and $\frac32$, and let $F=F_1\oplus F_2$. By Serre duality $H^1(F_1(t))$ and $H^1(F_2(t))$ are dual to the graded pieces of degree $-t-3$ of the cokernels of the transposed maps $S(-1)^2\oplus S(-2)^2\to S^2$ and $S\oplus S(-1)^2\to S(1)$, $S=K[x_0,x_1,x_2]$. For general maps these cokernels have finite length, with Hilbert functions $2,4,4,2$ in degrees $0,\dots,3$ (from the Buchsbaum--Rim complex, a resolution here) and $1,2,1$ in degrees $-1,0,1$ (a complete intersection of type $(1,2,2)$); the $h^0$- and $h^2$-rows then follow from the presentations, which are injective on global sections, and Riemann--Roch. Then $\gamma(F)=\gamma(F_1)+\gamma(F_2)$, with chain
\[
\tfrac12\gamma^{(-1,-4)}+\tfrac12\gamma^{(-2,-5)}+\tfrac25\gamma^{(-2,-6)}+\tfrac15\gamma^{(-2,-7)}+\tfrac25\gamma^{(-3,-7)},
\]
and $\Lambda_F(2)=2=\rk F_1\cdot(\mu_1-2)$. But also
\[
\gamma(F)=\gamma^{(-3,-4)}+\tfrac12\gamma^{(-2,-7)}+\tfrac14\gamma^{(-2,-5)}+\tfrac14\gamma^{(-1,-6)},
\]
which gives the value $1$ at $\lambda=2$. The functionals $\phi_{a,b}$ work the other way: each is an upper bound for $\Lambda$ on all of $\CC$, and $\Lambda$ is their minimum.
\end{remark}

Theorem~\ref{thm:main} says that $\Lambda_F\ge\Lambda_{P_F}$. On $\PP^2$ the table of every step of the Harder--Narasimhan filtration gives a stronger lower bound.

\begin{lemma}\label{lem:lowpair}
Let $\gamma\in\CC$ on $\PP^2$, $\gamma\ne0$, and let $B=M_2+1$ be the second root of the first element of its chain (Corollary~\textup{\ref{cor:endsP2}}). For every $c\in\frac12\ZZ$ there are integers $a>b$ with $a+b=2c$ and $b\le B$ such that $(a,b)$ is comparable with every element of the chain.
\end{lemma}

\begin{proof}
Let $z^1>\dots>z^m$ be the chain, so $z^1_2=B$. If $c\ge c(z^1)$, take $(a,b)=(2c-B,B)$; then $(a,b)\ge z^1$ and $a\ge z^1_1>B$. If $c\le c(z^m)$, take $(a,b)=(z^m_1,2c-z^m_1)$; then $(a,b)\le z^m$, and $b\le z^m_2<a$, $b\le z^1_2=B$. Otherwise $c(z^{k+1})\le c\le c(z^k)$ for some $k$; take $a=\min(z^k_1,\,2c-z^{k+1}_2)$ and $b=2c-a$. Then $z^{k+1}\le(a,b)\le z^k$ and $a>b$, so $(a,b)$ is comparable with the whole chain, and $b\le z^k_2\le B$.
\end{proof}

\begin{proposition}\label{prop:Mcrit}
Let $0\to G\to F\to Q\to0$ be an exact sequence of coherent sheaves on $\PP^2$, with $F$ a vector bundle, $G\ne0$ and $Q\ne0$ torsion-free, and let $M_2(F)=\max\{t:H^2(F(t))\ne0\}$. If $H^0\big(Q(M_2(F)+1)\big)=0$, then $G$ is locally free and $\Lambda_G\le\Lambda_F$; equivalently, $\Psi_G\le\Psi_F$ on $[0,\rk G]$. The hypothesis holds in each of the following cases:
\begin{enumerate}[label=\textup{(\alph*)}]
\item $\mu_{\max}(Q)<\mu_{\min}(F)+2$;
\item $Q$ is $\mu$-semistable and $\mu(Q)<\mu_{\min}(G)+2$.
\end{enumerate}
\end{proposition}

\begin{proof}
$G$ is saturated in $F$, hence reflexive, hence locally free. Let $\lambda\in\frac12\ZZ$, $c=-\lambda-\frac32$, and let $(a,b)$ be a pair for $\gamma(F)$ as in Lemma~\ref{lem:lowpair}. Multiplication by a linear form shows that $h^0(Q(t))$ is nondecreasing in $t$, since $Q$ is torsion-free; so $H^0(Q(b))=0$. The long exact cohomology sequences give $h^0_a(F)=h^0_a(G)+h^0_a(Q)-\rk\partial^0_a$, $h^0_b(F)=h^0_b(G)$ and $h^1_b(F)=h^1_b(G)+h^1_b(Q)-\rk\partial^1_b$, so
\[
(a-b)\big(\phi_{a,b}(\gamma(F))-\phi_{a,b}(\gamma(G))\big)=\big(h^0_a(Q)-\rk\partial^0_a\big)+\big(h^1_b(Q)-\rk\partial^1_b\big)\ \ge\ 0 .
\]
By Theorem~\ref{thm:three}, $\Lambda_F(\lambda)=\phi_{a,b}(\gamma(F))\ge\phi_{a,b}(\gamma(G))\ge\Lambda_G(\lambda)$. Both sides are linear between consecutive points of $\frac12\ZZ$, so $\Lambda_G\le\Lambda_F$ everywhere. Since $\Lambda_\gamma(\lambda)=\max_{0\le m\le\rk\gamma}(\Psi_\gamma(m)-m\lambda)$ and $\Psi_\gamma(m)=\min_\lambda(m\lambda+\Lambda_\gamma(\lambda))$ (Lemma~\ref{lem:legendre}(i)), this is equivalent to $\Psi_G\le\Psi_F$ on $[0,\rk G]$.

If $H^2(F(t))\ne0$, Serre duality gives a nonzero map $F\to\OO(-t-3)$, whose image is a torsion-free quotient of rank one and slope at most $-t-3$; so $M_2(F)+1\le-\mu_{\min}(F)-2$. In case (a), $\mu_{\max}(Q(M_2(F)+1))<0$, so $Q(M_2(F)+1)$ has no sections. Case (b) implies (a): the last Harder--Narasimhan factor $R$ of $F$, of slope $\mu_{\min}(F)$, is semistable, so $\mu(R)\ge\mu_{\min}(G)$ if $G\to R$ is nonzero, and otherwise $R$ is a quotient of $Q$; hence $\mu_{\min}(F)\ge\min(\mu_{\min}(G),\mu(Q))$, and $\mu(Q)<\mu_{\min}(G)+2$ gives $\mu_{\max}(Q)=\mu(Q)<\mu_{\min}(F)+2$.
\end{proof}

\begin{theorem}[monotonicity along the Harder--Narasimhan filtration]\label{thm:HNmono}
Let $0=F_0\subset F_1\subset\dots\subset F_s=F$ be the Harder--Narasimhan filtration of a vector bundle $F$ on $\PP^2$. Then
\[
\Lambda_{F_1}\le\Lambda_{F_2}\le\dots\le\Lambda_{F_s}=\Lambda_F,\qquad\text{i.e.}\qquad \Psi_{F_j}\le\Psi_{F_k}\ \text{on }[0,\rk F_j]\ \ (j\le k).
\]
\end{theorem}

\begin{proof}
Each $F_j$ is saturated in $F$, hence locally free. For $j\ge2$ the quotient $\cQ_j=F_j/F_{j-1}$ is $\mu$-semistable with $\mu(\cQ_j)<\mu(\cQ_{j-1})=\mu_{\min}(F_{j-1})$; apply Proposition~\ref{prop:Mcrit}(b) to $0\to F_{j-1}\to F_j\to\cQ_j\to0$.
\end{proof}

\begin{remark}\label{rem:HNmono}
\begin{enumerate}[label=(\arabic*)]
\item For $\lambda\ge\mu_{j+1}$ (with $\mu_{s+1}=-\infty$), $\Lambda_{P_F}=\Lambda_{P_{F_j}}\le\Lambda_{F_j}\le\Lambda_F$ by Theorem~\ref{thm:main} for $F_j$, so in this range Theorem~\ref{thm:HNmono} strengthens Theorem~\ref{thm:main}: it compares the whole table of each step with that of $F$, not only its rank and degree. Its proof uses Theorem~\ref{thm:three}, Corollary~\ref{cor:endsP2} and Serre duality, but neither Theorem~\ref{thm:main} nor Bogomolov's inequality, and it holds over any algebraically closed field. It uses the two extreme rows of the plane, and does not extend as it stands to $\PP^n$, $n\ge3$.
\item The constant $2$ in Proposition~\ref{prop:Mcrit}(a) is optimal: in Remark~\ref{rem:nomono}, $G=\OO\oplus\ker(\OO^3\to\OO(2))\subset F=\OO^4$ has $Q=\OO(2)$, so $\mu_{\max}(Q)=\mu_{\min}(F)+2$, while $\Lambda_G(0)=\frac13>0=\Lambda_F(0)$. In (b), $\mu_{\min}(G)$ cannot be replaced by $\mu(G)$: adding $\OO(11)$ to $G$ and to $F$ gives $\mu(G\oplus\OO(11))+2=\frac{17}4>2=\mu(Q)$, and the chains, which are the unions of those of $\OO(11)$ and of $G$, resp.\ $F$, still give $\Lambda_{G\oplus\OO(11)}(0)>\Lambda_{F\oplus\OO(11)}(0)$. Proposition~\ref{prop:firstedge}, under the different hypothesis $\mu_{\max}(F/F_1)<\mu(F_1)$, only compares the ends of the chains.
\end{enumerate}
\end{remark}

On $\PP^3$ the analogue of Corollary~\ref{cor:conicQ} is false, even with both bounds $d^j_t\le\min(\gamma_{1;j+1,t},\gamma_{2;j,t})$ that hold for actual extensions. So no argument that uses only the tables of $F'$ and $F''$ and these numerical bounds for the ranks of the connecting homomorphisms can prove Theorem~\ref{thm:main} on $\PP^3$; some geometric input beyond the tables is needed, such as the test sheaves.

\begin{example}[tables do not suffice on $\PP^3$]\label{ex:conicP3}
On $\PP^3$ put
\begin{align*}
\gamma_1&=\gamma^{(5,4,-1)}+4\gamma^{(5,4,-2)}+15\gamma^{(1,0,-3)}+6\gamma^{(1,0,-4)},\\
\gamma_2&=4\gamma^{(4,0,-1)}+11\gamma^{(3,0,-1)}+50\gamma^{(2,0,-1)}+9\gamma^{(1,-4,-5)},\\
\gamma&=100\,\gamma^{(1,0,-1)}=600\,\gamma(\OO(-2)).
\end{align*}
A direct computation from the definition of $\gamma^z$ gives
\[
\gamma_1+\gamma_2-\gamma=\sum_{t=-3}^{3}d_t\,\big(e_{1,t}+e_{2,t}\big),\qquad (d_{-3},\dots,d_3)=(72,162,216,180,168,114,48),
\]
with $d_t=\gamma_{1;2,t}=\gamma_{2;1,t}$; moreover $\gamma_1$ has no $h^1$ row and $\gamma_2$ no $h^2$ row. So $\gamma$ is what the long exact sequences give for an extension of a bundle with table $\gamma_2$ by a bundle with table $\gamma_1$ in which $\partial^0=\partial^2=0$ and every $\partial^1_t\colon H^1(F''(t))\to H^2(F'(t))$ is an isomorphism; and up to a common positive integer multiple, $\gamma_1$ and $\gamma_2$ are tables of direct sums of supernatural bundles. All hypotheses of the analogue of Corollary~\ref{cor:conicQ} hold. But $\rk\gamma_1=156$ and $\deg\gamma_1=-288$, while $\Psi_\gamma$ is linear of slope $-2$, so $\Psi_\gamma(156)=-312<\deg\gamma_1$.

Similarly, with
\begin{align*}
\gamma'_1&=4\gamma^{(6,5,-1)}+20\gamma^{(2,0,-2)}+5\gamma^{(2,0,-3)}+2\gamma^{(2,0,-4)}+\gamma^{(2,0,-5)},\\
\gamma'_2&=2\gamma^{(5,0,-1)}+4\gamma^{(4,0,-1)}+10\gamma^{(3,0,-1)}+2\gamma^{(2,-5,-6)},\\
\gamma'&=50\,\gamma^{(2,0,-1)},
\end{align*}
one finds
\[
\gamma'_1+\gamma'_2-\gamma'=\sum_{t=-4}^{4}d'_t\,\big(e_{1,t}+e_{2,t}\big),\qquad (d'_{-4},\dots,d'_4)=(24,60,96,120,120,160,144,96,40),
\]
with $d'_t=\min(\gamma'_{1;2,t},\gamma'_{2;1,t})$, while $\rk\gamma'_1=192$ and $\deg\gamma'_1=-440>-448=\Psi_{\gamma'}(192)$.
\end{example}

The stated properties follow directly from the definition of $\gamma^z$. By Theorem~\ref{thm:main} neither pair comes from an actual subsheaf, and the test sheaves of Lemma~\ref{lem:testsheaf} show this directly. A vector bundle $F$ with table $\gamma$ has no intermediate cohomology, so it is $\OO(-2)^{\oplus600}$ by Horrocks' criterion \cite{Ho64}. For $z=(1,0,-1)$ the test sheaf is $\Tt=\OO_L(1)$, with resolution $\OO(-1)\to\OO^2\to\OO(1)$, and $h^0(F\otimes \Tt)=0$; but a subbundle $F'\subset F$ with table $\gamma_1$ and locally free quotient, as in an extension, would give, for $L$ general, $0<24=\deg F'+2\rk F'=\chi(F'\otimes\Tt)\le h^0(F'\otimes\Tt)\le h^0(F\otimes\Tt)$. For $\gamma'$, which is a multiple of the supernatural table $\gamma^{(2,0,-1)}$, every bundle with table $\gamma'$ is semistable by Corollary~\ref{cor:supernatural}; the test sheaf is $\Tt=(\OO/\cI_L^2)(2)$, with resolution $\OO(-1)^2\to\OO^3\to\OO(2)$, and $\phi_{(2,0,-1)}(\gamma')=0$ gives $h^0(F\otimes \Tt)=0$, while such a subbundle with table $\gamma'_1$ would give $\chi(F'\otimes\Tt)=3(\deg F'+\frac73\rk F')=24>0$. The cancellation runs through $\partial^1$, between two intermediate rows. On $\PP^2$ every connecting homomorphism has its source in row $0$ or its target in row $2$, and these are the rows whose monotonicity Proposition~\ref{prop:fourdefect} uses; on $\PP^3$ the injection $F'\otimes \Tt\to F\otimes \Tt$ carries information that the tables do not record.

\section{Bridgeland stability and the large volume limit}\label{sec:lvl}

On $\PP^2$ the Harder--Narasimhan polygon is the large volume limit of the Harder--Narasimhan polygons of Bridgeland stability. In this section we show that Theorem~\ref{thm:main} survives the limit: for the stability conditions $\sigma_{\alpha,\beta}$ with $\alpha\gg0$, the Bridgeland Harder--Narasimhan polygon of a vector bundle lies inside a polygon built from its Boij--Söderberg decomposition. The leading order is Theorem~\ref{thm:main}; the second order rests on the fact that its equality cases are determined at the level of Chern characters (Lemma~\ref{lem:Ktight}), and that a Harder--Narasimhan factor along a contact edge is a supernatural bundle (Proposition~\ref{prop:contactedge}, which rests on a Bogomolov-type bound for the gap, Lemma~\ref{lem:gapdisc}). For Theorem~\ref{thm:lvl} and in Remark~\ref{rem:lvl} the ground field is $\mathbb C$, as in the references; Subsection~\ref{sec:equality} holds over any algebraically closed field.

\subsection{Stability conditions and polygons}
We use the stability conditions $\sigma_{\alpha,\beta}$ of \cite{Br08,AB13,ABCH13}, with heart $\operatorname{Coh}^\beta(\PP^2)$ and central charge $Z_{\alpha,\beta}$; see \cite{MS17} for an exposition. For $\beta\in\RR$, the heart $\operatorname{Coh}^\beta(\PP^2)$ is the tilt of $\operatorname{Coh}(\PP^2)$ at the torsion pair $(\mathcal T^\beta,\mathcal F^\beta)$: $\mathcal T^\beta$ consists of the sheaves whose torsion-free part has $\mu_{\min}>\beta$, and $\mathcal F^\beta$ of the torsion-free sheaves with $\mu_{\max}\le\beta$. For $\alpha>0$ the central charge is $Z_{\alpha,\beta}(E)=-\int_{\PP^2}e^{-(\beta+i\alpha)H}\operatorname{ch}(E)$, i.e.
\[
Z_{\alpha,\beta}(E)=-\Big(\operatorname{ch}_2(E)-\beta c_1(E)+\frac{\beta^2-\alpha^2}2\,r(E)\Big)+i\alpha\big(c_1(E)-\beta\,r(E)\big),
\]
which depends only on the class $v(E)=(r,c_1,\operatorname{ch}_2)$ of $E$. Every nonzero object $E$ of the heart has $Z_{\alpha,\beta}(E)\in\RR_{>0}\,e^{i\pi\phi}$ with a phase $\phi(E)\in(0,1]$, and a Harder--Narasimhan filtration with $\sigma_{\alpha,\beta}$-semistable factors $E_1,\dots,E_m$ of strictly decreasing phases \cite{Br07}. The \emph{Harder--Narasimhan polygon} $P^\sigma_E$ is the polygonal path with vertices $0,\ Z(E_1),\ Z(E_1)+Z(E_2),\ \dots,\ Z(E)$; together with the segment $[0,Z(E)]$ it bounds a convex region, which we denote $|P^\sigma_E|$. A vector bundle $F$ with $\mu_{\min}(F)>\beta$ lies in $\operatorname{Coh}^\beta(\PP^2)$.

The Boij--Söderberg decomposition gives a second polygon in the same plane. A supernatural bundle with table $\gamma^z$ has Hilbert polynomial $(t-z_1)(t-z_2)$, so by Riemann--Roch its class is
\[
v(z)=\big(2,\ 2\mu(z),\ z_1z_2-3\mu(z)-2\big),\qquad \mu(z)=-\tfrac{z_1+z_2}2-\tfrac32 .
\]

\begin{definition}\label{def:PsiSigma}
Let $F$ be a vector bundle on $\PP^2$ with chain decomposition $\gamma(F)=\sum_kq_k\gamma^{z^k}$, let $\sigma=\sigma_{\alpha,\beta}$, and suppose that the vectors $q_kZ_{\alpha,\beta}(v(z^k))$ lie in an open half-plane $\mathcal H$ bounded by a line through $0$. The \emph{Boij--Söderberg polygon} $\Psi^\sigma_F$ is the polygonal path from $0$ to $Z_{\alpha,\beta}(F)$ whose edges are these vectors, arranged in order of decreasing argument, measured in an interval of length $\pi$ corresponding to $\mathcal H$; $|\Psi^\sigma_F|$ is the convex region it bounds with the segment $[0,Z_{\alpha,\beta}(F)]$.
\end{definition}

Without the half-plane condition the path need not bound a region: for $F=\OO(2)\oplus\OO\oplus\OO(-2)$ and $(\alpha,\beta)=(1,0)$ the edges are $-\frac32+2i$, $\frac12$ and $-\frac32-2i$, and the first and the last edge meet at $-\frac54+\frac53i$. The condition holds in the two cases we use. For $\alpha\gg0$ all edges have positive real part, by \eqref{eq:Zxy} below; and if $\beta<\mu^\BS_{\min}(F)=\min_k\mu(z^k)$, all edges have positive imaginary part. In the second case the supernatural bundles $E_{z^k}$ lie in $\operatorname{Coh}^\beta(\PP^2)$, and if they are $\sigma$-semistable with distinct phases, then $\Psi^\sigma_F$ is the Harder--Narasimhan polygon of the supernatural model $\bigoplus_kE_{z^k}^{\oplus m q_k}$, up to the factor $m$. Note that $\mu^\BS_{\min}(F)\le\mu_{\min}(F)$ by Corollary~\ref{cor:supernatural}(b), so the hypothesis $\beta<\mu_{\min}(F)$ of Theorem~\ref{thm:lvl} below does not imply $\beta<\mu^\BS_{\min}(F)$.

\begin{theorem}[large volume limit]\label{thm:lvl}
Let $F$ be a vector bundle on $\PP^2$ and $\beta<\mu_{\min}(F)$ a rational number. There is $\alpha_0$ such that for all $\alpha\ge\alpha_0$, with $\sigma=\sigma_{\alpha,\beta}$,
\[
|P^\sigma_F|\ \subseteq\ |\Psi^\sigma_F| .
\]
\end{theorem}

For $r>0$ and $v=(r,c_1,\operatorname{ch}_2)$ write $\mu=c_1/r$ and $\operatorname{ch}_2^\beta=\operatorname{ch}_2-\beta c_1+\beta^2r/2$. Then
\begin{equation}\label{eq:Zxy}
\operatorname{Re}Z_{\alpha,\beta}(v)=\frac{\alpha^2r}2-\operatorname{ch}_2^\beta,\qquad
\operatorname{Im}Z_{\alpha,\beta}(v)=\alpha r(\mu-\beta),\qquad
\frac{\operatorname{Im}Z}{\operatorname{Re}Z}=\frac{2(\mu-\beta)}{\alpha-2\operatorname{ch}_2^\beta/(\alpha r)} .
\end{equation}
So for a fixed finite set of classes of positive rank and $\alpha\gg0$, the real parts are positive, the order of the arguments is the order of the slopes $\mu$, and among classes with the same $\mu>\beta$ the argument increases with $\operatorname{ch}_2/r$: on classes with $\mu>\beta$ the order of phases is the Gieseker order. (For $\mu<\beta$ the second comparison is reversed, but it is not needed: no two elements of a chain have the same slope.) We use the large volume limit in the following form:

\smallskip\noindent\textup{(LV)} \emph{For a class $v$ of positive rank with $\mu(v)>\beta$ there is $\alpha_v$ such that for $\alpha>\alpha_v$ an object of $\operatorname{Coh}^\beta(\PP^2)$ of class $v$ is $\sigma_{\alpha,\beta}$-semistable if and only if it is a $(\beta H-\frac12K_{\PP^2})$-twisted Gieseker semistable sheaf.}

\smallskip\noindent This is the analogue for surfaces of \cite[Prop.~14.2]{Br08}, see \cite[\S6, Exercise~6.27]{MS17}, and for $\PP^2$ \cite[\S6]{ABCH13}. We apply (LV) only for rational $\beta$, only in the direction ``Gieseker semistable $\Rightarrow$ $\sigma_{\alpha,\beta}$-semistable'', and only to the finitely many Gieseker factors of a fixed $F$. In this form it follows from \cite{LQ14}, where $\sigma_{\alpha,\beta}$ is the stability condition $(Z_m,\mathcal P_m)$ with $m=\alpha$ and $\omega=H$. Let $G$ be a twisted Gieseker semistable torsion-free sheaf with $\mu(G)>\beta$, and let $A\subset G$ be a proper subobject in $\operatorname{Coh}^\beta(\PP^2)$. As in the proof of \cite[Lemma~5.1]{LQ14}, $A$ is a torsion-free sheaf, and either $\mu(A)<\mu(G)$, or $A$ is a subsheaf of $G$ with $\mu(A)=\mu(G)$; in the second case the comparison of phases \cite[(2.7)]{LQ14} is the comparison of the twisted reduced Hilbert polynomials of $A$ and $G$. So the phase of $A$ is at most that of $G$ for $\alpha\gg0$: the sheaf $G$ is semistable for the polynomial stability condition of Bayer \cite{Bay09} at the large volume limit, and by \cite[Lemma~4.3]{LQ14} it is $\sigma_{\alpha,\beta}$-semistable for all $\alpha\ge M$, where $M$ depends only on $\beta$ and on the class of $G$. By the following lemma the twist plays no role.

\begin{lemma}\label{lem:twist}
On $\PP^2$, for every $b\in\RR$, a torsion-free sheaf is $bH$-twisted Gieseker semistable if and only if it is Gieseker semistable.
\end{lemma}

\begin{proof}
The $bH$-twisted reduced Hilbert polynomial of a torsion-free sheaf $\cE$ of rank $r>0$ is $\frac1r\int\operatorname{ch}(\cE)e^{-bh}e^{th}\operatorname{td}(\PP^2)$. Its coefficient of $t$ is $\mu(\cE)-b+\frac32$, and for fixed $\mu(\cE)$ its constant term is $\operatorname{ch}_2(\cE)/r$ plus a number depending only on $\mu(\cE)$ and $b$. So twisted and untwisted reduced Hilbert polynomials compare $\mu$ first and then, for equal $\mu$, $\operatorname{ch}_2/r$, in the same way. The argument is unchanged if $\operatorname{td}(\PP^2)$ is replaced by $\sqrt{\operatorname{td}(\PP^2)}$, as in some definitions of twisted stability.
\end{proof}

Consequently, if $F$ is a vector bundle with $\mu_{\min}(F)>\beta$ and $0=G_0\subset G_1\subset\dots\subset G_p=F$ is its Gieseker Harder--Narasimhan filtration, then for $\alpha\gg0$, depending on the finitely many classes of the factors, the factors $G_i/G_{i-1}$ lie in the heart, are $\sigma_{\alpha,\beta}$-semistable, and have strictly decreasing phases; so this filtration is the $\sigma_{\alpha,\beta}$-Harder--Narasimhan filtration of $F$, and $P^\sigma_F$ has vertices $Z(G_i)$. The Gieseker filtration refines the $\mu$-Harder--Narasimhan filtration $0=F_0\subset\dots\subset F_s=F$: on each $\mu$-semistable factor $\cQ_j=F_j/F_{j-1}$ it induces the Gieseker filtration of $\cQ_j$, whose factors are $\mu$-semistable of slope $\mu_j$ with decreasing $\operatorname{ch}_2/r$.

\subsection{The equality cases of Theorem~A on \texorpdfstring{$\PP^2$}{P2}}\label{sec:equality}
We write $n_j=\rk F_j$ and $\mu_1>\dots>\mu_s$ for the slopes of the $\mu$-Harder--Narasimhan factors, and we call a vertex $(n_j,\deg F_j)$ of $P_F$ \emph{tight} if $\Psi_F(n_j)=\deg F_j$.

\begin{lemma}[contact]\label{lem:contact}
If $\Psi_F(m)=P_F(m)$ for some $m$ in the interior of an edge of $P_F$, then $\Psi_F=P_F$ on that edge. So the contact set $\{m:\Psi_F(m)=P_F(m)\}$ is the union of the tight vertices and of the edges of $P_F$ whose both endpoints are tight and on which $\Psi_F$ is linear.
\end{lemma}

\begin{proof}
On the edge, $\Psi_F-P_F$ is concave, nonnegative by Theorem~\ref{thm:main}, and vanishes at an interior point; a nonnegative concave function on an interval that vanishes at an interior point vanishes identically. For the last assertion, if $\Psi_F=P_F$ on an edge then $\Psi_F$ is linear there, and conversely.
\end{proof}

\begin{lemma}[tight vertices are determined by the Chern character]\label{lem:Ktight}
Let $(n_j,\deg F_j)$ be a tight vertex, $0<j<s$, and let $\lambda\in\frac12\ZZ$ be a slope of $\Psi_F$ at $n_j$, i.e.\ $\Psi_F(n_j)=\min_{\lambda'}\big(n_j\lambda'+\Lambda_F(\lambda')\big)$ is attained at $\lambda$. Let $(a,b)$ be a pair adapted to the chain of $F$ with $a+b=-2\lambda-3$, as in Theorem~\textup{\ref{thm:three}(b)}, and let
\[
x=\sum_{\mu(z^k)>\lambda}q_k\gamma^{z^k}+f\,\gamma^{(a,b)},\qquad f\ge0 \text{ such that } \rk x=n_j .
\]
Then $x$ is well defined, $\chi_x(t)=\chi(F_j(t))$ for all $t$, i.e.\ $x$ and $F_j$ have the same Chern character, and
\[
h^0\big((F/F_j)(a)\big)=0,\qquad h^2\big(F_j(b)\big)=0 .
\]
\end{lemma}

\begin{proof}
Put $c=-\lambda-\frac32$, so $c(a,b)=c$. Since $\lambda$ minimizes $n_j\lambda'+\Lambda_F(\lambda')$, whose one-sided derivatives at $\lambda$ are $n_j-\sum_{\mu(z^k)\ge\lambda}2q_k\le0\le n_j-\sum_{\mu(z^k)>\lambda}2q_k$, the pieces of slope $>\lambda$ have total rank $\le n_j$, and those of slope $\ge\lambda$ total rank $\ge n_j$. The pieces of slope $>\lambda$ are those with center $<c$; being comparable with $(a,b)$, they are the $z^k<(a,b)$. Let $x'$ be their sum. If $\rk x'<n_j$, the chain contains a piece of slope exactly $\lambda$, of rank at least $n_j-\rk x'$; it is comparable with $(a,b)$ and has the same center, so it is $(a,b)$. So $x$ is well defined, $f$ is at most the coefficient of $\gamma^{(a,b)}$ in the chain, $\rk x=n_j$, and $\deg x=\deg x'+\lambda(n_j-\rk x')=n_j\lambda+\Lambda_F(\lambda)=\Psi_F(n_j)=\deg F_j$.

For a root sequence $z$ comparable with $(a,b)$ put $P_z(t)=(t-z_1)(t-z_2)$. If $z<(a,b)$ then $z_1\le a$, so the column $a$ of $\gamma^z$ has its entry in row $0$, equal to $P_z(a)$; and $z_2\le b$, so the column $b$ has its entry in row $0$ or $1$, and $\gamma^z_{0,b}-\gamma^z_{1,b}=P_z(b)$ in both cases. If $z\ge(a,b)$ then $\gamma^z_{0,a}=0$ and $\gamma^z_{0,b}=\gamma^z_{1,b}=0$. Since $P_{(a,b)}(a)=P_{(a,b)}(b)=0$, summing over the chain gives
\begin{equation}\label{eq:h0chi}
h^0(F(a))=\chi_x(a),\qquad h^0(F(b))-h^1(F(b))=\chi_x(b).
\end{equation}
The polynomials $\chi_x$ and $\chi(F_j(\,\cdot\,))$ have the same coefficients of $t^2$ and $t$, because $x$ and $F_j$ have the same rank and degree; so $D=\chi_x-\chi(F_j(\,\cdot\,))$ is a constant, equal to $\operatorname{ch}_2(x)-\operatorname{ch}_2(F_j)$.

Put $F'=F_j$, $F''=F/F_j$, which is torsion-free, and write $h^i_t(G)=h^i(G(t))$. By Lemma~\ref{lem:legendre}(i), $\Lambda_F(\lambda)=\Psi_F(n_j)-n_j\lambda=n_j(\mu(F_j)-\lambda)$, so the four terms of Proposition~\ref{prop:fourdefect} vanish; in particular $\rk\partial^0_a=h^1_a(F')$ and $\rk\partial^1_b=h^1_b(F'')$. From the long exact sequences and \eqref{eq:h0chi},
\begin{align*}
D&=h^0_a(F)-\chi(F'(a))=h^0_a(F')+h^0_a(F'')-\rk\partial^0_a-\chi(F'(a))=h^0_a(F'')-h^2_a(F'),\\
D&=h^0_b(F)-h^1_b(F)-\chi(F'(b))=h^0_b(F'')-h^1_b(F'')+\rk\partial^1_b-h^2_b(F')=h^0_b(F'')-h^2_b(F').
\end{align*}
By Theorem~\ref{thm:main}, $\Psi_F\ge P_F$ with equality at $n_j$, so the slopes of $\Psi_F$ at $n_j$ lie between the slopes $\mu_{j+1}$ and $\mu_j$ of $P_F$ there; in particular $\mu_{j+1}\le\lambda\le\mu_j$. Now $h^0_b(F'')=0$: a nonzero section of $F''(b)$ gives a nonzero map $\OO(-b)\to F''$, so $-b\le\mu_{\max}(F'')=\mu_{j+1}\le\lambda$, while $b<c=-\lambda-\frac32$ gives $-b>\lambda$. And $h^2_a(F')=h^0\big(F'^\vee(-a-3)\big)=0$: a nonzero section gives $\OO(a+3)\to F'^\vee$, so $a+3\le\mu_{\max}(F'^\vee)=-\mu_{\min}(F')=-\mu_j$, while $a>c$ gives $-a-3<-c-3=\lambda-\frac32<\mu_j$. Hence $D=h^0_a(F'')\ge0$ and $D=-h^2_b(F')\le0$, so $D=0$, $h^0_a(F'')=0$ and $h^2_b(F')=0$.
\end{proof}

\begin{lemma}[gap and discriminant]\label{lem:gapdisc}
Let $z=(a,b)$ be a root sequence with gap $g=a-b$, and let $\cQ'$ be a $\mu$-semistable torsion-free sheaf on $\PP^2$ of slope $\mu(z)$ with $H^0(\cQ'(a))=0$. Then
\[
\delta(\cQ')\ \ge\ \frac{g^2-1}4=\delta(E_z),
\]
with equality if and only if $H^1(\cQ'(a))=0$; here $E_z$ is a supernatural bundle with table $\gamma^z$.
\end{lemma}

\begin{proof}
By Serre duality $H^2(\cQ'(a))$ is dual to $H^0(\cQ'^\vee(-a-3))$, and $\cQ'^\vee$ is $\mu$-semistable of slope $-\mu(z)$. A nonzero section would give $a+3\le-\mu(z)=\frac{a+b}2+\frac32$, i.e.\ $a+3\le b$, which is false. So $\chi(\cQ'(a))=-h^1(\cQ'(a))\le0$. By \eqref{eq:RRP2}, with $a+\mu(z)+\frac32=a-c(z)=\frac g2$,
\[
\frac{\chi(\cQ'(a))}{\rk\cQ'}=\frac{g^2}8-\frac18-\frac{\delta(\cQ')}2\ \le\ 0 .
\]
For $E_z$, Lemma~\ref{lem:basic}(ii) gives $4\delta(E_z)=G(E_z)-1=g^2-1$.
\end{proof}

\begin{proposition}[contact edges]\label{prop:contactedge}
Let $F$ be a vector bundle on $\PP^2$, and suppose that $\Psi_F=P_F$ on the edge $[n_{j-1},n_j]$ of $P_F$. Then $\cQ_j=F_j/F_{j-1}$ is a vector bundle whose table is a multiple of $\gamma^z$, where $z$ is the unique element of the chain of $F$ of slope $\mu_j$.
\end{proposition}

\begin{proof}
Two comparable root sequences with the same center are equal, so $z=(a,b)$ is unique, and it is the only pair adapted to the chain with $a+b=-2\mu_j-3$. Put $g=a-b$ and $r=\rk\cQ_j$. Since $\Psi_F$ is linear of slope $\mu_j$ on the edge, $\mu_j$ is a slope of $\Psi_F$ at $n_{j-1}$ and at $n_j$. At an interior endpoint, Lemma~\ref{lem:Ktight} with $\lambda=\mu_j$ and the pair $(a,b)$ identifies the Chern character of $F_{j-1}$, resp.\ $F_j$, with that of the prefix of the chain consisting of the pieces of slope $>\mu_j$ and a fraction of $\gamma^z$; at the endpoints $0$ and $\rk F$ the prefixes are $0$ and $\gamma(F)$, and the pieces of slope $\ne\mu_j$ do not occur between them because $\Psi_F$ has slope $\mu_j$ on the whole edge. Subtracting, $\chi(\cQ_j(t))=\frac r2(t-a)(t-b)$, and $\delta(\cQ_j)=\frac{g^2-1}4$.

If $j\ge2$, Lemma~\ref{lem:Ktight} at $F_{j-1}$ gives $H^0((F/F_{j-1})(a))=0$, hence $H^0(\cQ_j(a))=0$; if $j=1$, $z$ is the last element of the chain, every element $z^k$ has $z^k_1\ge a$, so $H^0(F(a))=0$ and again $H^0(\cQ_1(a))=0$. Similarly, if $j<s$, Lemma~\ref{lem:Ktight} at $F_j$ gives $H^2(F_j(b))=0$, and $H^2(F_j(b))\to H^2(\cQ_j(b))$ is surjective; if $j=s$, $z$ is the first element of the chain, every $z^k$ has $z^k_2\le b$, so $H^2(F(b))=0$, which surjects onto $H^2(\cQ_s(b))$. Since $\cQ_j$ is semistable of slope $\mu_j=-\frac{a+b}2-\frac32$, also $H^0(\cQ_j(b))=0$ and, by Serre duality, $H^2(\cQ_j(a))=0$. As $\chi(\cQ_j(a))=\chi(\cQ_j(b))=0$, the sheaf $\cQ_j$ has no cohomology at the twists $a$ and $b$.

Let $E=\cQ_j^{\vee\vee}$ and $\ell=\operatorname{length}(E/\cQ_j)$, so that $\delta(E)=\delta(\cQ_j)-2\ell/r$. The dual $E^\vee=\cQ_j^\vee$ is semistable of slope $-\mu_j$, which is the slope of the root sequence $(-b-3,-a-3)$, of gap $g$, and $H^0(E^\vee(-b-3))$ is dual to $H^2(E(b))\cong H^2(\cQ_j(b))=0$, since $E$ is locally free and $E/\cQ_j$ has finite support. Lemma~\ref{lem:gapdisc} gives $\delta(E)=\delta(E^\vee)\ge\frac{g^2-1}4=\delta(\cQ_j)$, so $\ell=0$ and $\cQ_j$ is locally free. The columns $a$ and $b$ of its table vanish, so every element of its chain is $(a,b)$.
\end{proof}

\begin{corollary}\label{cor:equalityP2}
For a vector bundle $F$ on $\PP^2$, $P_F=\Psi_F$ if and only if the Harder--Narasimhan factors of $F$ are vector bundles whose tables are multiples of supernatural tables $\gamma^{z(1)},\dots,\gamma^{z(s)}$ with pairwise comparable root sequences. In that case $\gamma(F)=\sum_j\gamma(\cQ_j)$ is the chain decomposition.
\end{corollary}

\begin{proof}
If $P_F=\Psi_F$, every edge is a contact edge, and Proposition~\ref{prop:contactedge} applies; the root sequences belong to the chain of $F$, so they are comparable. The converse holds on every $\PP^n$ (Remark~\ref{rem:main}(1)).
\end{proof}

\subsection{Proof of Theorem~\ref{thm:lvl}}
Fix $\beta<\mu_{\min}(F)$ and write $Z=Z_{\alpha,\beta}$, $x(v)=\operatorname{Re}Z(v)$, $y(v)=\operatorname{Im}Z(v)$. Let $\alpha$ be large enough that $P^\sigma_F$ is the Gieseker polygon, with vertices $Z(G_i)$, and that the arguments of the edges of both polygons are ordered as described after \eqref{eq:Zxy}. Then the edges of $\Psi^\sigma_F$ are the pieces in decreasing order of slope. All edges of both polygons have $x>0$ and decreasing $y/x$, so both paths are graphs of concave piecewise linear functions $p$ and $\psi$ on $[0,x(F)]$, and $|P^\sigma_F|\subseteq|\Psi^\sigma_F|$ if and only if $p\le\psi$.

Since $\psi$ is concave and $p$ is linear between consecutive vertices $Z(G_i)$, it suffices to show that each of the finitely many vertices $Z(G_i)$ lies on or below the path $\Psi^\sigma_F$ for $\alpha$ large; then one $\alpha_0$ works for all of them. Put $r_i=\rk G_i$; since the Gieseker filtration refines the Harder--Narasimhan filtration, $\deg G_i=P_F(r_i)$.

The affine map $\Phi(x,y)=\big(2x/\alpha^2,\ y/\alpha+2\beta x/\alpha^2\big)$ sends $Z(v)$ to $(r,c_1)+O(\alpha^{-2})$, uniformly for $v$ in the finite set of classes of the Gieseker factors and of the pieces, and it preserves the vertical order of points with the same first coordinate. So $\Phi(Z(G_i))=(r_i,P_F(r_i))+O(\alpha^{-2})$, and $\Phi(\Psi^\sigma_F)$, as a graph over the first coordinate, lies within $O(\alpha^{-2})$ of $\Psi_F$. Hence if $\Psi_F(r_i)>P_F(r_i)$, the vertex $Z(G_i)$ lies strictly below $\Psi^\sigma_F$ for $\alpha$ large.

Suppose $\Psi_F(r_i)=P_F(r_i)$. If $r_i$ lies in the interior of an edge $[n_{j-1},n_j]$ of $P_F$, then $\Psi_F=P_F$ on that edge by Lemma~\ref{lem:contact}, so $\cQ_j$ is a supernatural bundle by Proposition~\ref{prop:contactedge}, hence Gieseker semistable by Proposition~\ref{prop:sameslope}; then the Gieseker filtration has no step strictly between $F_{j-1}$ and $F_j$, a contradiction. So $G_i=F_j$ and $(n_j,\deg F_j)$ is a tight vertex. If $0<j<s$, Lemma~\ref{lem:Ktight} gives $v(F_j)=v(x)$ for a prefix $x$ of the chain, consisting of the pieces of slope $>\lambda$ and a fraction of the piece of slope $\lambda$, for a slope $\lambda$ of $\Psi_F$ at $n_j$; since the edges of $\Psi^\sigma_F$ are the pieces in order of slope, $Z(F_j)=Z(x)$ lies on $\Psi^\sigma_F$. The endpoints $0$ and $Z(F)$ lie on both paths. This proves $p\le\psi$ on $[0,x(F)]$ for $\alpha\ge\alpha_0$. \qed

\begin{lemma}[support functions]\label{lem:support}
Let $\sigma=\sigma_{\alpha,\beta}$, identify $\mathbb C$ with $\RR^2$, with $\langle u,w\rangle=\operatorname{Re}(\bar uw)$, and for $X$ in $\operatorname{Coh}^\beta(\PP^2)$ let $h_X$ be the support function of $|P^\sigma_X|$. For $\theta\in(0,1)$ put $\ell_\theta=\operatorname{Im}(e^{-i\pi\theta}\,\cdot\,)=\langle e^{i\pi(\theta+1/2)},\cdot\,\rangle$, and let $X_{>\theta}\subset X$ be the step of the Harder--Narasimhan filtration with factors of phase $>\theta$.
\begin{enumerate}[label=\textup{(\alph*)}]
\item $h_X(e^{i\pi(\theta+1/2)})=\ell_\theta(Z(X_{>\theta}))$, and $h_X(u)=\max(0,\langle u,Z(X)\rangle)$ if $\operatorname{Re}u\ge0$. In particular $h_{X\oplus Y}(u)=h_X(u)+h_Y(u)$ for $\operatorname{Re}u\le0$.
\item For every subobject $S\subset X$, $\ell_\theta(Z(S))\le\ell_\theta(Z(X_{>\theta}))$.
\item If $0\to A\to B\to C\to0$ is exact in $\operatorname{Coh}^\beta(\PP^2)$, then $|P^\sigma_{B\oplus D}|\subseteq|P^\sigma_{A\oplus C\oplus D}|$ for every $D$ in $\operatorname{Coh}^\beta(\PP^2)$.
\end{enumerate}
\end{lemma}

\begin{proof}
(a) Along the path $P^\sigma_X$ the increments of $\ell_\theta$ are first positive and then $\le0$, so the maximum of $\ell_\theta$ on $|P^\sigma_X|$ is $\ell_\theta(Z(X_{>\theta}))$; since $(X\oplus Y)_{>\theta}=X_{>\theta}\oplus Y_{>\theta}$, $h_{X\oplus Y}=h_X+h_Y$ in these directions. In the remaining directions $u$, those with $\operatorname{Re}u\ge0$, the increments of $\langle u,\cdot\,\rangle$ along the path are first $\le0$ and then $\ge0$, so $h_X(u)=\max(0,\langle u,Z(X)\rangle)$ depends only on $Z(X)$. Additivity on the boundary $\operatorname{Re}u=0$ follows by continuity from the left half-plane; no additivity is asserted in the right half-plane. (b) We have $S_{>\theta}\subset X_{>\theta}$, since $\operatorname{Hom}(S_{>\theta},X/X_{>\theta})=0$, and $X_{>\theta}/S_{>\theta}$ has factors of phase $>\theta$, while $S/S_{>\theta}$ has factors of phase $\le\theta$; so $\ell_\theta(Z(S))\le\ell_\theta(Z(S_{>\theta}))\le\ell_\theta(Z(X_{>\theta}))$. (c) Applied to $B_{>\theta}\cap A\subset A$ and to its image in $C$, (b) gives $h_B\le h_A+h_C=h_{A\oplus C}$ in the directions $e^{i\pi(\theta+1/2)}$, and $h_B=h_{A\oplus C}$ in the other directions, since $Z(B)=Z(A\oplus C)$. For $\operatorname{Re}u<0$, adding $h_D(u)$ and using (a) gives $h_{B\oplus D}(u)\le h_{A\oplus C\oplus D}(u)$. For $\operatorname{Re}u\ge0$, the two support functions are equal by (a), since $Z(B\oplus D)=Z(A\oplus C\oplus D)$. The support-function criterion for containment of compact convex sets proves the assertion.
\end{proof}

\begin{remark}\label{rem:lvl}
\begin{enumerate}[label=(\arabic*)]
\item Lemma~\ref{lem:Ktight} says that the equality cases of Theorem~\ref{thm:main} on $\PP^2$ are visible in $K$-theory: at a tight vertex the pieces above the vertex have the Chern character of the Harder--Narasimhan step, and in particular the same discriminant. Lemma~\ref{lem:gapdisc} is Bogomolov's inequality improved by the gap: a $\mu$-semistable sheaf of slope $\mu(z)$ without sections at the twist $a$ has discriminant at least that of the supernatural bundles with roots $z$.
\item The statement does not extend to all $(\alpha,\beta)$ at the level of the Gieseker polygon. Let $S=\coker(\OO(-2)^2\to\OO(-1)^2\oplus\OO^2)$ for a general map, a stable bundle of slope $1$ with $h^0(S(t))=2\big(h^0(\OO(t-1))+h^0(\OO(t))-h^0(\OO(t-2))\big)$, $h^2(S(t))=h^0(S(-5-t))$ and $\chi(S(t))=t^2+5t+2$, and let $F=\OO(2)\oplus S\oplus T(-2)$, where here and below $T=T_{\PP^2}$ is the tangent bundle. Its Harder--Narasimhan filtration is $\OO(2)\subset\OO(2)\oplus S\subset F$, and its chain is
\[
\gamma^{(0,-2)}+\tfrac35\gamma^{(0,-4)}+\tfrac1{20}\gamma^{(0,-5)}+\tfrac{13}{30}\gamma^{(-1,-5)}+\tfrac1{12}\gamma^{(-2,-5)}+\tfrac13\gamma^{(-3,-5)} .
\]
The vertex at rank $3$ is tight, and $v(\OO(2)\oplus S)=(3,4,-1)$ is the class of the five pieces of slope $\ge\frac12$, which are the pieces above this vertex in the sense of Lemma~\ref{lem:Ktight} (with $\lambda=\frac12$ and $f=\frac35$, the whole piece $\gamma^{(0,-4)}$), as that lemma predicts. The vertex at rank $1$ is not: $\Psi_F(1)=\frac94>2$. For $(\alpha,\beta)=(\frac15,-\frac{11}{20})$, where all edges of $\Psi^\sigma_F$ have positive imaginary part, neither the Gieseker polygon of $F$ nor its $\sigma$-Harder--Narasimhan polygon is contained in $|\Psi^\sigma_F|$. The sections of $S$ give an exact sequence $0\to\OO^2\to S\to Q\to0$, with $Q=\coker(\OO(-2)^2\to\OO(-1)^2)$ supported on a conic; as $Q$ and $T(-2)$, of slope $-\frac12>\beta$, lie in $\mathcal T^\beta$, $G=\OO(2)\oplus\OO^2$ is a subobject of $F$ in $\operatorname{Coh}^\beta(\PP^2)$. Its central charge $Z(G)=-\frac{559}{160}+\frac{73}{100}\,i$ lies strictly beyond the edge $\frac{13}{30}Z(v(-1,-5))$ of $\Psi^\sigma_F$, of phase $\theta\approx0.648$, and hence so does the vertex $Z(F_{>\theta})$ of $P^\sigma_F$, since $\ell_\theta(Z(G))\le\ell_\theta(Z(F_{>\theta}))$ by Lemma~\ref{lem:support}(b). (So $S$ is $\sigma$-unstable: by the sequence above, $\OO\subset S$ is a subobject in $\operatorname{Coh}^\beta(\PP^2)$, and $\OO$ and $S$ have phases $0.778$ and $0.115$.) In the same example, $v(\OO(2))=(1,2,2)$ is not of the form $\sum_kt_kq_kv(z^k)$ with $0\le t_k\le1$: the combinations with rank $1$ and degree $2$ have $\operatorname{ch}_2\le\frac{31}{16}$. So the containment for all $(\alpha,\beta)$, which would follow if every Harder--Narasimhan step were such a fractional sum of pieces, cannot be proved by an argument on classes alone.
\item The containment cannot hold for all $\sigma$ in the geometric chamber $\{\sigma_{\alpha,\beta}:\alpha>0,\ \beta\in\RR\}$, because the pieces themselves cross walls. For $F=T(-1)$, with $\gamma(F)=\gamma^{(-1,-3)}$, the polygon $\Psi^\sigma_F$ is the segment $[0,Z(T(-1))]$ for every $\sigma$. But for $-1\le\beta<0$ the Euler sequence becomes the exact sequence $0\to\OO^3\to T(-1)\to\OO(-1)[1]\to0$ in $\operatorname{Coh}^\beta(\PP^2)$, and inside the semicircle $(\beta+\frac12)^2+\alpha^2<\frac14$ the subobject $\OO^3$ has larger phase than $T(-1)$ (at $(\alpha,\beta)=(\frac3{10},-\frac12)$ the phases are $0.656$ and $0.583$), so $T(-1)$ is unstable and $|P^\sigma_{T(-1)}|$ has nonempty interior. Below the walls of the pieces, the polygon to compare with is the Harder--Narasimhan polygon of a supernatural model, which is an object of $\operatorname{Coh}^\beta(\PP^2)$ when $\beta<\mu^\BS_{\min}(F)$. We ask: for a vector bundle $F$ on $\PP^2$ with chain $\sum_kq_k\gamma^{z^k}$, are there supernatural bundles $E_{z^k}$ with tables $\gamma^{z^k}$, and $m$ with $mq_k\in\ZZ$, such that
\[
|P^\sigma_F|\ \subseteq\ \tfrac1m\,\big|P^\sigma_{\bigoplus_kE_{z^k}^{\oplus mq_k}}\big|
\]
for every $\sigma=\sigma_{\alpha,\beta}$ with $\alpha>0$ and $\beta<\mu^\BS_{\min}(F)$? For rational $\beta$ and $\alpha\gg0$ this holds for every choice, by Theorem~\ref{thm:lvl}.
\item The answer is yes for $F=\OO\oplus E$, where $E=\coker(\OO(-2)^2\to\OO(-1)^4)$ is a stable bundle with $c_1=0$, $c_2=2$ and table $\gamma^{(0,-3)}$. The chain of $F$ is $\frac13\gamma^{(0,-2)}+\frac56\gamma^{(0,-3)}+\frac13\gamma^{(-1,-3)}$, so with $m=6$ a supernatural model is $T(-2)^{\oplus2}\oplus E^{\oplus5}\oplus T(-1)^{\oplus2}$, and $\mu^\BS_{\min}(F)=-\frac12$. The four rows of the presentation of $E$ are linearly independent in $\operatorname{Hom}(\OO(-2)^2,\OO(-1))\cong H^0(\OO(1))^{\oplus2}$, because $\operatorname{Hom}(E,\OO(-1))=H^0(E(-1))=0$; here and below we use $E^\vee\cong E$, as $E$ has rank $2$ and $c_1=0$. Completing them to a basis gives a presentation $\OO(-2)^2\to\OO(-1)^6$ of $T(-2)^{\oplus2}$, and projecting to the first four coordinates gives an exact sequence $0\to\OO(-1)^2\to T(-2)^{\oplus2}\to E\to0$. Since $\operatorname{Ext}^1(E,\OO)=H^1(E)=0$, its class is the image of a map $E\to T(-1)^{\oplus2}$ under the connecting map of the Euler sequence $0\to\OO(-1)^2\to\OO^6\to T(-1)^{\oplus2}\to0$; so $T(-2)^{\oplus2}$ is the fiber product of $\OO^6$ and $E$ over $T(-1)^{\oplus2}$, and
\[
0\to T(-2)^{\oplus2}\to\OO^{6}\oplus E\to T(-1)^{\oplus2}\to0
\]
is exact. For $\beta<-\frac12$ all its terms lie in $\mathcal T^\beta$, so it is exact in $\operatorname{Coh}^\beta(\PP^2)$. By Lemma~\ref{lem:support}(c), with $B=\OO^6\oplus E$, $A=T(-2)^{\oplus2}$, $C=T(-1)^{\oplus2}$ and $D=E^{\oplus5}$, we get $6\,|P^\sigma_F|=|P^\sigma_{F^{\oplus6}}|\subseteq|P^\sigma_{T(-2)^{\oplus2}\oplus E^{\oplus5}\oplus T(-1)^{\oplus2}}|$ for all $\alpha>0$ and $\beta<-\frac12$.
\end{enumerate}
\end{remark}

\section{Polarized varieties}\label{sec:polarized}

In this section $n\ge2$ is the dimension of a polarized variety, which we project to $\PP^n$.

\subsection{The gap moment on \texorpdfstring{$\PP^n$}{Pn}}
For a root sequence $z=(z_1>\dots>z_n)$ put
\[
\rho(z)=\frac{12}{n(n^2-1)}\sum_{i=1}^n\big(z_i-c(z)\big)^2 ,
\]
so that $\rho(z)=g(z)^2$ for $n=2$. For a vector bundle $F$ of rank $N$ on $\PP^n$ with chain decomposition $\gamma(F)=\sum_\ell q_\ell\gamma^{z^\ell}$ put $p_\ell=n!\,q_\ell/N$, $c_\ell=c(z^\ell)$, and
\[
G(F)=\sum_\ell p_\ell\,\rho(z^\ell),\qquad \bar c=\sum_\ell p_\ell c_\ell,\qquad V_c(F)=\sum_\ell p_\ell(c_\ell-\bar c)^2 ,
\]
which agree with the definitions of Section~\ref{sec:plane} for $n=2$. For a torsion-free sheaf $\cF$ of rank $N$ on $\PP^n$ let $c_1$ and $\operatorname{ch}_2$ be the degrees of $c_1(\cF)\cdot h^{n-1}$ and $\operatorname{ch}_2(\cF)\cdot h^{n-2}$, $\mu(\cF)=c_1/N$, and $\delta(\cF)=\mu^2-2\operatorname{ch}_2/N$, the degree of $\Delta(\cF)\cdot h^{n-2}/N^2$ with $\Delta=2Nc_2-(N-1)c_1^2$.

\begin{lemma}\label{lem:gapPn}
Let $F$ be a vector bundle of rank $N$ on $\PP^n$, $n\ge2$.
\begin{enumerate}[label=\textup{(\roman*)}]
\item $\displaystyle\frac{\chi(F(t))}N=\frac{t^n}{n!}-\bar c\,\frac{t^{n-1}}{(n-1)!}+\Big(\frac{V_c(F)+\bar c^{\,2}}2-\frac{(n+1)\,G(F)}{24}\Big)\frac{t^{n-2}}{(n-2)!}+O(t^{n-3})$.
\item $\bar c=-\mu(F)-\frac{n+1}2$, and $G(F)-1=\dfrac{12}{n+1}\big(V_c(F)+\delta(F)\big)$.
\item $G(F)\ge1$, with equality iff $F$ has no intermediate cohomology, iff $F$ splits.
\item $\displaystyle\sum_{0<i<n}\sum_th^i(F(t))\ \ge\ C_n\,N\big(G(F)-1\big)$, where $C_n=\dfrac{n\,b_n}{12\,(n-2)!}$ and $b_n=(\lfloor n/2\rfloor-1)!\,(\lceil n/2\rceil-1)!$; so $C_2=\frac16$, $C_3=\frac14$, $C_4=\frac16$. The constant is optimal: equality holds for $F=\Omega^u_{\PP^n}(u)$, $u=\lfloor n/2\rfloor$, and in general iff every element of the chain of $F$ has all gaps $z_u-z_{u+1}$ equal to $1$, or exactly one equal to $2$, in a position $u\in\{\lfloor n/2\rfloor,\lceil n/2\rceil\}$.
\item If $0=F_0\subset\dots\subset F_s=F$ is the Harder--Narasimhan filtration, with factors $\cQ_i$ of ranks $N_i$, weights $w_i=N_i/N$ and slopes $\mu_i$, then
\[
G(F)-1=\frac{12}{n+1}\Big(V_c(F)-V_{HN}(F)+\sum_iw_i\,\delta(\cQ_i)\Big),\qquad V_{HN}(F)=\sum_iw_i(\mu_i-\mu(F))^2 .
\]
\end{enumerate}
For $n=2$, (i)--(iii) and (v) are Lemma~\ref{lem:basic}(i)--(iii),(v), and (iv) is the first bound of Lemma~\ref{lem:basic}(iv).
\end{lemma}

\begin{proof}
(i) The Hilbert polynomial is linear in the table and $N=n!\sum_\ell q_\ell$, so $\chi(F(t))/N=\frac1{n!}\sum_\ell p_\ell\prod_i(t-z^\ell_i)$. For a root sequence $z$ with center $c$,
\[
\prod_i(t-z_i)=t^n-nc\,t^{n-1}+e_2(z)\,t^{n-2}-\dots,\qquad e_2(z)=\sum_{i<j}z_iz_j=\frac{n(n-1)}2c^2-\frac12\sum_i(z_i-c)^2 ,
\]
where $\frac12\sum_i(z_i-c)^2=\frac{n(n^2-1)}{24}\rho(z)$; and $\sum_\ell p_\ell c_\ell^2=V_c(F)+\bar c^{\,2}$.

(ii) By Riemann--Roch, $\chi(F(t))=\int\operatorname{ch}(F)e^{th}\operatorname{td}(\PP^n)$ with $\operatorname{td}(\PP^n)=1+\frac{n+1}2h+\frac{(n+1)(3n+2)}{24}h^2+\dots$, so the coefficients of $t^{n-1}/(n-1)!$ and of $t^{n-2}/(n-2)!$ in $\chi(F(t))/N$ are $\mu+\frac{n+1}2$ and $\operatorname{ch}_2/N+\frac{n+1}2\mu+\frac{(n+1)(3n+2)}{24}$. Comparing with (i) gives $\bar c=-\mu-\frac{n+1}2$, and then, since $\bar c^{\,2}=\mu^2+(n+1)\mu+\frac{(n+1)^2}4$,
\[
\frac{V_c+\mu^2}2+\frac{n+1}{24}-\frac{(n+1)G}{24}=\frac{\operatorname{ch}_2}N ,
\]
which is the identity.

(iii) For a root sequence $z$, $n\sum_i(z_i-c)^2=\sum_{i<j}(z_i-z_j)^2\ge\sum_{i<j}(j-i)^2=\frac{n^2(n^2-1)}{12}$, since $z_i-z_j\ge j-i$ for $i<j$; so $\rho(z)\ge1$, with equality iff the roots are consecutive, iff $\gamma^z$ has no intermediate cohomology. The entries of $\gamma(F)=\sum_\ell q_\ell\gamma^{z^\ell}$ are sums of nonnegative terms, so $F$ has no intermediate cohomology iff every $z^\ell$ is consecutive, iff $G(F)=1$; and a bundle on $\PP^n$ without intermediate cohomology splits by Horrocks' criterion \cite{Ho64}.

(iv) Let $z$ be a root sequence, $g_u=z_u-z_{u+1}=1+a_u$ with integers $a_u\ge0$, $1\le u<n$, and let $I(z)=\sum_{0<i<n}\sum_t\gamma^z_{i,t}$ be its total intermediate cohomology. The left side of (iv) is $\sum_\ell q_\ell I(z^\ell)$, and $\sum_\ell q_\ell(\rho(z^\ell)-1)=N(G(F)-1)/n!$, so it suffices to prove $I(z)\ge n!\,C_n(\rho(z)-1)=\frac{b_nn^2(n-1)}{12}(\rho(z)-1)$. Put $D(z)=\sum_{i<j}\big((z_i-z_j)^2-(j-i)^2\big)$; by the proof of (iii), $D(z)=\frac{n^2(n^2-1)}{12}(\rho(z)-1)$. Writing $z_i-z_j=(j-i)+\sum_{u=i}^{j-1}a_u$ and expanding,
\begin{equation}\label{eq:Dexp}
D(z)=\sum_uu(n-u)\big(na_u+a_u^2\big)+2\sum_{u<v}u(n-v)\,a_ua_v .
\end{equation}
An integer $t$ with $z_{u+1}<t<z_u$ contributes $\prod_j|t-z_j|=(z_u-t)(t-z_{u+1})\prod_{j<u}(z_j-t)\prod_{j>u+1}(t-z_j)$ to $I(z)$, and
\[
z_j-t\ge u-j+1+\sum_{v=j}^{u-1}a_v\quad(j<u),\qquad t-z_j\ge j-u+\sum_{v=u+1}^{j-1}a_v\quad(j>u+1).
\]
Multiplying these lower bounds and keeping the constant and the linear terms, $\prod_{j\ne u,u+1}|t-z_j|\ge B_u\big(1+\sum_{v\ne u}c_{uv}a_v\big)$, where $B_u=u!\,(n-u)!$, $c_{uv}=\sum_{j=v+1}^n\frac1{j-u}$ for $v>u$ and $c_{uv}=\sum_{j=1}^v\frac1{u-j+1}$ for $v<u$. Summing over the gap, $\sum_t(z_u-t)(t-z_{u+1})=f(a_u)$ with $f(a)=\frac{a(a+1)(a+2)}6\ge a$, so
\[
I(z)\ \ge\ \sum_uB_uf(a_u)+\sum_{u<v}\big(B_uc_{uv}+B_vc_{vu}\big)a_ua_v .
\]
Now $B_u=u(n-u)\,(u-1)!\,(n-u-1)!\ge u(n-u)\,b_n$, and for $u<v$ we have $c_{uv}\ge\frac{n-v}{n-u}$ and $c_{vu}\ge\frac uv$, so $B_uc_{uv}+B_vc_{vu}\ge2b_nu(n-v)$. Moreover, for integers $a\ge0$,
\[
f(a)-\frac{na+a^2}{n+1}=\frac{a(a-1)\big((n+1)a+4n-2\big)}{6(n+1)}\ \ge\ 0 .
\]
Comparing with \eqref{eq:Dexp} gives $I(z)\ge\frac{b_n}{n+1}D(z)=\frac{b_nn^2(n-1)}{12}(\rho(z)-1)$. For $F=\Omega^u_{\PP^n}(u)$ the roots are $\{0,-1,\dots,-n\}\setminus\{-u\}$ by Bott's formula, so $a_u=1$ and all other $a_v$ vanish; $\rk F=\binom nu$, the only intermediate cohomology is $H^u(F(-u))=K$, and $G(F)-1=\rho-1=\frac{12u(n-u)}{n^2(n-1)}$. Since $\binom nu u(n-u)=n!/b_n$ for $u=\lfloor n/2\rfloor$, this gives equality. For the equality cases, it suffices by additivity to treat one root sequence. Since $B_uc_{uv}+B_vc_{vu}\ge2b_nu(n-v)>2b_nu(n-v)/(n+1)$, equality forces at most one $a_u>0$; the displayed identity excludes $a_u\ge2$; and for $a_u=1$ the remaining condition $B_u=u(n-u)b_n$ means $(u-1)!(n-u-1)!=b_n$, i.e.\ $u\in\{\lfloor n/2\rfloor,\lceil n/2\rceil\}$. Conversely, if all $a_v$ vanish then $I(z)=D(z)=0$, and if only $a_u=1$ then $I(z)=B_u$ and $D(z)=(n+1)u(n-u)$, so equality holds for central $u$.

(v) The numbers $c_1$ and $\operatorname{ch}_2$ are additive, so $\delta(F)=\sum_iw_i\delta(\cQ_i)+\mu(F)^2-\sum_iw_i\mu_i^2=\sum_iw_i\delta(\cQ_i)-V_{HN}(F)$; insert this in (ii).
\end{proof}

\subsection{The defect identity}
Let $X$ be a smooth projective variety of dimension $n\ge2$ over $K$ and $H$ an ample, base-point-free divisor. Then $n+1$ general sections of $\OO_X(H)$ define a finite morphism $\pi\colon X\to\PP^n$ with $\pi^*\OO(1)=\OO_X(H)$, flat because $X$ is Cohen--Macaulay and $\PP^n$ is smooth. Put
\[
d=H^n,\qquad k=K_X\cdot H^{n-1},\qquad \sigma=\frac{(K_X^2-2c_2(X))\cdot H^{n-2}}{n+1}.
\]
For $n=2$, Noether's formula gives $\sigma=K_X^2-8\chi(\OO_X)$, the topological signature of $X$ over $\mathbb C$. We keep the notation $\sigma$, $\alpha$, $\beta$ of \cite{An1,An2}; in this section these letters have nothing to do with the stability conditions $\sigma_{\alpha,\beta}$ of Section~\ref{sec:lvl}. For a vector bundle $\cE$ of rank $r$, $\pi_*\cE$ is a vector bundle of rank $N=rd$, and $H^i(X,\cE(tH))=H^i(\PP^n,\pi_*\cE(t))$. We set $G(\cE)=G(\pi_*\cE)$ and $V_c(\cE)=V_c(\pi_*\cE)$. The bundle $\cE$ is aCM, i.e.\ $H^i(X,\cE(tH))=0$ for $0<i<n$ and all $t$, iff $\pi_*\cE$ splits, iff $G(\cE)=1$ (Lemma~\ref{lem:gapPn}(iii)).

For a torsion-free sheaf $\cQ$ of rank $m$ on $X$ put $v=c_1(\cQ)/m\in\operatorname{NS}(X)_\QQ$, $\mu_H(\cQ)=v\cdot H^{n-1}$ and
\[
\delta(\cQ)=\Big(v^2-\frac{2\operatorname{ch}_2(\cQ)}m\Big)\cdot H^{n-2}=\frac{\Delta(\cQ)\cdot H^{n-2}}{m^2},\qquad
\beta(\cQ)=\frac{\mu_H(\cQ)-k/2}d,
\]
\[
\alpha(\cQ)=v-\frac{K_X}2-\beta(\cQ)H ,
\]
so that $\alpha(\cQ)\cdot H^{n-1}=0$. In the computations below every product of two classes is understood to be intersected with $H^{n-2}$, and a product of a class with $H$ to be intersected with $H^{n-2}$ as well; in particular $\alpha(\cQ)^2$ denotes the number $\alpha(\cQ)^2\cdot H^{n-2}$ and $v\cdot H$ the number $v\cdot H^{n-1}$. If $\cQ$ is $\mu_H$-semistable, then $\delta(\cQ)\ge0$ by Bogomolov's inequality \cite{Bo78}, \cite[Thm.~7.3.1]{HL10}. Always $\alpha(\cQ)^2\le0$: choose $k_0\ge1$ with $k_0H$ very ample; by Bertini, $n-2$ general members of $|k_0H|$ cut out a smooth connected surface $S$, and $k_0^{n-2}\alpha(\cQ)^2=(\alpha(\cQ)|_S)^2\le0$ by the Hodge index theorem on $S$, since $\alpha(\cQ)|_S\cdot H|_S=0$.

Let $0=\cE_0\subset\dots\subset\cE_s=\cE$ be the $\mu_H$-Harder--Narasimhan filtration of $\cE$, with factors $\cQ_i$ of ranks $r_i$ and slopes $\mu_i$. Put $w_i=r_i/r$, $\beta_i=\beta(\cQ_i)=(\mu_i-k/2)/d$, $\bar\beta=\sum_iw_i\beta_i$ and $V_\beta(\cE)=\sum_iw_i(\beta_i-\bar\beta)^2$.

\begin{lemma}\label{lem:push}
Let $\cQ$ be a torsion-free sheaf of rank $m$ on $X$.
\begin{enumerate}[label=\textup{(\alph*)}]
\item $\pi_*\cQ$ is a torsion-free sheaf of rank $md$ and slope $\beta(\cQ)-\frac{n+1}2$ on $\PP^n$.
\item $\displaystyle\frac{\chi(\cQ(tH))}m=d\,\frac{t^n}{n!}+d\,\beta(\cQ)\,\frac{t^{n-1}}{(n-1)!}+\Big(\frac d2\beta(\cQ)^2-\frac{(n+1)\,\sigma}{24}+\frac{\alpha(\cQ)^2-\delta(\cQ)}2\Big)\frac{t^{n-2}}{(n-2)!}+O(t^{n-3})$.
\end{enumerate}
For $n=2$ this is \cite[Lemma~2.1 and (2)]{An2}; in every dimension it is equivalent to \cite[Prop.~10.6]{An3}.
\end{lemma}

\begin{proof}
(b) By Riemann--Roch, $\chi(\cQ(tH))=\int_X\operatorname{ch}(\cQ)e^{tH}\operatorname{td}(X)$ with $\operatorname{td}(X)=1-\frac{K_X}2+\frac{K_X^2+c_2(X)}{12}+\dots$. The coefficient of $t^{n-1}/(n-1)!$ is $(c_1(\cQ)-\frac m2K_X)\cdot H^{n-1}=m\,d\,\beta(\cQ)$. The coefficient of $t^{n-2}/(n-2)!$ is $\big(\operatorname{ch}_2(\cQ)-\frac12c_1(\cQ)K_X+\frac m{12}(K_X^2+c_2(X))\big)\cdot H^{n-2}$; dividing by $m$ and substituting $\operatorname{ch}_2(\cQ)\cdot H^{n-2}=\frac m2(v^2-\delta(\cQ))$ and $v\cdot K_X=v^2+\frac{K_X^2}4-2\beta v\cdot H+\beta k+\beta^2d-\alpha(\cQ)^2$, with $v\cdot H=\beta d+\frac k2$, gives the stated coefficient.

(a) Since $\pi$ is finite and $X$ is integral, $\pi_*\cQ$ is torsion-free, of rank $md$. As $\chi(\pi_*\cQ(t))=\chi(\cQ(tH))$, the coefficient of $t^{n-1}/(n-1)!$ computed on $\PP^n$, $\deg\pi_*\cQ+md\,\frac{n+1}2$, equals $md\,\beta(\cQ)$.
\end{proof}

\begin{proposition}[defect identity]\label{prop:defect}
For every vector bundle $\cE$ on $X$,
\[
\frac{n+1}{12}\big(d\,G(\cE)-\sigma\big)=d\,\big(V_c(\cE)-V_\beta(\cE)\big)+\sum_iw_i\big(\delta(\cQ_i)-\alpha(\cQ_i)^2\big).
\]
\end{proposition}

\begin{proof}
Summing Lemma~\ref{lem:push}(b) over the $\cQ_i$ with weights $w_i$ gives $\chi(\cE(tH))/r$. On the other hand $\chi(\cE(tH))/r=d\,\chi(\pi_*\cE(t))/N$, which is given by Lemma~\ref{lem:gapPn}(i). Comparing the coefficients of $t^{n-1}$ gives $\bar c=-\bar\beta$; comparing the coefficients of $t^{n-2}/(n-2)!$ and subtracting $\frac d2\bar c^{\,2}=\frac d2\bar\beta^2$ gives
\[
\frac d2V_c-\frac{(n+1)\,dG}{24}=\frac d2V_\beta-\frac{(n+1)\,\sigma}{24}+\frac12\sum_iw_i(\alpha_i^2-\delta_i). \qedhere
\]
\end{proof}

In terms of the variance $v(\pi_*\cE)=-\delta(\pi_*\cE)$ of the push-forward, the part of Proposition~\ref{prop:defect} that lives on $X$ is the identity $\frac{n+1}{12}(d-\sigma)=d\big(v(\pi_*\cE)-V_\beta\big)+\sum_iw_i(\delta_i-\alpha_i^2)$ of \cite[(10.3)]{An3} (for $n=2$, \cite[Prop.~9.1]{An3}; for aCM bundles, \cite[Prop.~2.6]{An2}); Lemma~\ref{lem:gapPn}(ii), $\delta(\pi_*\cE)=\frac{n+1}{12}(G-1)-V_c$, converts it into Proposition~\ref{prop:defect}. What the table adds is the term $V_c$, which Theorem~\ref{thm:main} compares with $V_\beta$ for every bundle.

\begin{corollary}\label{cor:cases}
Let $\cE$ be a vector bundle of rank $r$ on $X$.
\begin{enumerate}[label=\textup{(\alph*)}]
\item $G(\cE)\ge\sigma/d$, with equality iff $V_c(\cE)=V_\beta(\cE)$ and $\delta(\cQ_i)=\alpha(\cQ_i)^2=0$ for all $i$. Moreover
\[
\sum_{0<i<n}\sum_th^i(\cE(tH))\ \ge\ C_n\,r\,(\sigma-d)_+ ,
\]
with $C_n$ as in Lemma~\ref{lem:gapPn}(iv); for $n=2$ also $\sum_th^1(\cE(tH))\ge rd\,L(\kappa)$, where $\kappa=\max(1,\sigma/d)$ and $L\colon[1,\infty)\to\RR$ is the function with $L(k^2)=\frac{k^3-k}{12}$ for integers $k\ge1$ that is linear on each interval $[k^2,(k+1)^2]$.
\item If $\cE$ is $\mu_H$-semistable, then $G(\cE)\ge\sigma/d+\frac{12}{n+1}V_c(\cE)$.
\item \textup{(aCM, \cite[Thm.~1.1, Prop.~5.1]{An2}, \cite[\S10.7]{An3})} If $\cE$ is aCM, then $d\ge\sigma$, i.e.\ $(n+1)H^n\ge(K_X^2-2c_2(X))\cdot H^{n-2}$.
\end{enumerate}
\end{corollary}

\begin{proof}
(a) The subsheaves $\pi_*\cE_i\subset\pi_*\cE$ form a filtration with torsion-free quotients $\pi_*\cQ_i$ of slopes $\beta_i-\frac{n+1}2$ (Lemma~\ref{lem:push}(a)), which strictly decrease. So its polygon $P_\beta$, with these slopes and lengths $dr_i$, is concave, has the same endpoints as $\Psi_{\pi_*\cE}$, and its vertices are the points $(\rk\pi_*\cE_i,\deg\pi_*\cE_i)$. By Theorem~\ref{thm:main} they lie below $\Psi_{\pi_*\cE}$, so $P_\beta\le\Psi_{\pi_*\cE}$ by Lemma~\ref{lem:legendre}(ii). As in the proof of Corollary~\ref{cor:supernatural}(c), Karamata's inequality gives $V_\beta(\cE)\le V_c(\cE)$. Both terms on the right of Proposition~\ref{prop:defect} are now nonnegative, and they vanish iff $V_c=V_\beta$ and $\delta(\cQ_i)=\alpha(\cQ_i)^2=0$ for all $i$. The bound on the intermediate cohomology follows from Lemma~\ref{lem:gapPn}(iv) with $N=rd$, since $G(\cE)-1\ge(\sigma-d)/d$; for $n=2$, Lemma~\ref{lem:basic}(iv) gives $\sum_th^1(\cE(tH))=\frac{rd}{12}\sum_\ell p_\ell(g_\ell^3-g_\ell)=rd\sum_\ell p_\ell L(g_\ell^2)$; the slopes of $L$ on the intervals $[k^2,(k+1)^2]$ are $\frac{k(k+1)}{4(2k+1)}$, which increase with $k$, so $L$ is convex and increasing, and Jensen's inequality gives $\sum_th^1(\cE(tH))\ge rd\,L(G(\cE))\ge rd\,L(\kappa)$.

(b) Here $V_\beta=0$, and the last sum in Proposition~\ref{prop:defect} is $\ge0$.

(c) If $\cE$ is aCM, then $G(\cE)=1$, and (a) gives $d\ge\sigma$.
\end{proof}

\begin{remark}
\begin{enumerate}[label=(\arabic*)]
\item Corollary~\ref{cor:cases}(b) is the Boij--Söderberg form of Bogomolov's inequality. The point of (a) is uniformity: all cohomology tables of bundles on $(X,H)$ lie in the region $G\ge\sigma/d$, a uniform lower bound for the gap moment, while the aCM bundles are those with $G=1$. This is not a bound for a distance between tables: on $\PP^2$ the normalized tables $(1-p)\gamma^{(-1,-2)}+p\gamma^{(m,0)}$ in $\CC$, of total rank $2$, with $a\in\QQ_{>0}$ fixed, $p=a/(m^2-1)$ and $m^2>a+1$, have $G=1+a$ for every $m$, while they converge entrywise to $\gamma^{(-1,-2)}=2\gamma(\OO)$ as $m\to\infty$. For the Hesse surfaces $Y_\nu$ with $H_\nu=4A_\nu-E_\nu$ of \cite[Cor.~1.4]{An2} (denoted $Y_n$ there), $\sigma/d=(3\nu^2-11)/7$, which is $\frac{16}7$ for $\nu=3$ and $\frac{37}7$ for $\nu=4$. Since $L$ is the convex piecewise linear interpolation of $s\mapsto\frac{\sqrt s(s-1)}{12}$ at the squares, it dominates this function and $s\mapsto\frac{s-1}6$; for $\nu=4$ it gives $\sum_th^1(\cE(tH))\ge\frac{31}{35}rd$, against $\frac57rd$ from the linear bound.
\item If $S$ is a smooth surface cut out by $n-2$ general members of $|H|$, adjunction gives $K_S^2-2c_2(S)=(n+1)\sigma-(n-2)d$, i.e.\ $3(\sigma_S-d)=(n+1)(\sigma-d)$; so $\sigma>d$ holds for $(X,H)$ iff it holds for its surface sections; this is the case $a=1$ of the formula for the excess $(n+1)(\sigma-d)$ under sections in $|aH|$ \cite[Rem.~5.7]{An2}, \cite[Lemma~3.3]{An3}, compare the proof of \cite[Lem.~3.1]{Lo26}. For $X=Y\times\PP^{n-2}$ with $H=H_Y\boxtimes\OO(1)$ one finds $(n+1)(\sigma-d)=3\big(\sigma_Y-(n-1)^2H_Y^2\big)$ \cite[Cor.~5.3]{An2}, \cite[Cor.~10.9]{An3}; for $Y=Y_\nu$ this is positive iff $(3\nu^2-11)/7>(n-1)^2$, e.g.\ for $\nu\ge4$ when $n=3$. Lopez \cite[Thm.~1]{Lo26} constructs from the Hesse surfaces varieties $Y_\nu\times Z$, $\nu\gg0$, of every dimension $n\ge3$, and their hyperplane sections, with $\sigma>d$ for very ample $H$. On all of them (a) bounds the intermediate cohomology of every vector bundle from below.
\item The rank of $\pi_*\cE$ is $rd\ge d$, and $\pi_*\cQ_i$ need not be semistable, so $P_\beta$ is in general not the Harder--Narasimhan polygon of $\pi_*\cE$; the proof uses only that its vertices come from subsheaves. When the $\cQ_i$ are locally free, Corollary~\ref{cor:monoext}(a) gives the opposite bound $V_c(\cE)\le V_c(\gr\cE)$: cancellation can only decrease the spread of the centers, and Theorem~\ref{thm:main} says that it cannot decrease it below the spread $V_\beta$ of the slopes.
\item For $n=2$, Lemma~\ref{lem:twotwist} transports verbatim: $G(\cE)-1\ge\frac4{rd}\big(h^1(\cE(tH))+h^1(\cE((t+1)H))\big)$ for every $t$.
\item The proof of (c) in \cite[Prop.~5.1]{An2} restricts to a surface section; \cite[\S10.7]{An3} uses the projection to $\PP^n$ directly, with Horrocks' splitting and the majorization for split bundles \cite[Lemma~10.8]{An3}. Here Theorem~\ref{thm:main} replaces that majorization, so the identity of Proposition~\ref{prop:defect} measures the defect of every bundle, not only of the aCM ones.
\item Characteristic zero enters (a)--(c) only through Bogomolov's inequality, and it cannot be dropped. Lemmas~\ref{lem:gapPn} and~\ref{lem:push}, Proposition~\ref{prop:defect} and Theorem~\ref{thm:main} hold over any algebraically closed field. Over a field of characteristic $p>0$, every polarized smooth surface carries Ulrich bundles \cite[Thm.~1.1]{Pi26a} (stated for $H$ very ample; for $H$ ample see \cite[Lemma~2.11]{Pi26a}), and on every polarized smooth variety a Frobenius direct image $F^e_*(\cE(aH))$ of any vector bundle is aCM for suitable $e$ and $a$ (\cite[Prop.~4.11]{Pi26b}, for $H$ very ample; the argument works for every ample $H$: choose $M$ with $H^i(\cE(tH))=0$ for $0<i<n$ and $|t|\ge M$, and take $a=M$, $p^e\ge2M$; then $a+p^et\notin(-M,M)$ for all $t$, and the projection formula applies); these bundles have $G=1$ whatever $\sigma/d$ is. For them Proposition~\ref{prop:defect} gives $\sum_iw_i\delta(\cQ_i)\le\frac{n+1}{12}(d-\sigma)$, so if $\sigma>d$ some Harder--Narasimhan factor violates Bogomolov's inequality and is therefore not strongly semistable \cite{La04}. For the Ulrich bundles of \cite{Pi26a} on a surface this is the computation of \cite[Rem.~2.9]{Pi26a}. Conversely, in characteristic $p$ the bound $G(\cE)\ge\sigma/d$ of (a) holds for every $\cE$ whose Harder--Narasimhan factors are strongly semistable. The contrast also shows up in the cone question of the next subsection: if $\cU$ is Ulrich on $(X,H)$, then $\pi_*(\cU\otimes\pi^*S)\cong S^{\oplus\rk\pi_*\cU}$ for every vector bundle $S$ on $\PP^n$, so the tables $\gamma(\pi_*\cE)$ of vector bundles $\cE$ on $(X,H)$ generate the cone $\CC$ (see \cite[Thm.~5.3]{ESi}); for surfaces in characteristic $p$ this is \cite[Cor.~1.2]{Pi26a}. In every dimension, over an infinite perfect field of characteristic $p$ and for $H$ ample and base-point-free, the cones of cohomology tables of coherent sheaves on $X$ and on $\PP^n$ have the same closure under entrywise limits, by a theorem of Iyengar, Ma and Walker \cite[Thm.~3.4]{IMW23}. In characteristic zero with $\sigma>d$, by (a) the tables $\gamma(\pi_*\cE)$ lie in the proper subcone $\{G\ge\sigma/d\}$; by part~(1) of this remark, this bound does not pass to entrywise limits.
\end{enumerate}
\end{remark}

\subsection{Vanishing at \texorpdfstring{$n$}{n} twists}\label{sec:twists}
After \cite{An1}, D.~Eisenbud and F.-O.~Schreyer asked the first author (personal communication) whether a surface $Y$ carries vector bundles $\cE$ with $h^i(\cE(a_1H))=h^i(\cE(a_2H))=0$ for $i=0,1,2$ at two non-consecutive integers $a_1<a_2$, the first case being $a_1=-1$, $a_2=1$, and what the Boij--Söderberg cone of cohomology tables of aCM bundles on $Y$ is. For consecutive integers such bundles are, up to twist, the Ulrich bundles; the aCM part is answered in \cite{An2}. For $H^2<\sigma$ the question becomes one about all vector bundles, and \cite[\S12, Question~4]{An3} asks whether their cone is cut out from that of $\PP^2$ by conditions like $v(\pi_*\cE)\le\frac14-\frac\sigma{4d}$, which holds for $\mu_H$-semistable $\cE$, once unstable bundles are taken into account. Corollary~\ref{cor:cases}(a) gives such a condition for every bundle, $G\ge\sigma/d$. Supernatural bundles with bounded roots on polarized varieties are studied in \cite{GL26}; what Corollary~\ref{cor:ES} adds is the explicit dependence on the normalized spread $\rho$ of the roots, through the defect identity. The structure of $\Psi$ gives a necessary condition for each configuration of $n$ twists.

\begin{corollary}\label{cor:ES}
Let $\cE\ne0$ be a vector bundle on $X$ and $a_1<\dots<a_n$ integers with $h^i(\cE(a_jH))=0$ for all $i$ and $j$. Then $\gamma(\pi_*\cE)$ is a multiple of $\gamma^{z}$, $z=(a_n,\dots,a_1)$, $\cE$ is $\mu_H$-semistable, and
\[
\frac{n+1}{12}\big(d\,\rho(z)-\sigma\big)\ =\ \delta(\cE)-\alpha(\cE)^2\ \ge\ 0 .
\]
In other words, a supernatural ray $\RR_{>0}\gamma^{z}$ contains the table of $\pi_*\cE$ for some vector bundle $\cE$ on $X$ only if $\rho(z)\ge\sigma/d$. If the $a_j$ form an arithmetic progression of difference $e$, then $\rho(z)=e^2$ and the condition reads $e^2d\ge\sigma$. In particular, if $\sigma>4d$ there is no such bundle on a surface with $a_2-a_1=2$, e.g.\ with $a_1=-1$ and $a_2=1$. On the Hesse surfaces $Y_\nu$ of \cite[Cor.~1.4]{An2} there is none with $(a_2-a_1)^2<(3\nu^2-11)/7$; for instance none with $a_2-a_1=2$ once $\nu\ge4$.
\end{corollary}

\begin{proof}
The columns $a_1,\dots,a_n$ of $\gamma(\pi_*\cE)=\sum_\ell q_\ell\gamma^{z^\ell}$ vanish. Its entries are sums of nonnegative terms, so these columns vanish in every $\gamma^{z^\ell}$: the $a_j$ are roots of every $z^\ell$, i.e.\ $z^\ell=z$. By Corollary~\ref{cor:supernatural}(a) the bundle $\pi_*\cE$ is semistable. If $\cE$ were not $\mu_H$-semistable, its maximal destabilizing subsheaf $\cE_1$ would give $\pi_*\cE_1\subset\pi_*\cE$ of slope $\beta_1-\frac{n+1}2>\bar\beta-\frac{n+1}2=\mu(\pi_*\cE)$ (Lemma~\ref{lem:push}(a)), a contradiction. Now $G(\cE)=\rho(z)$, $V_c(\cE)=0$ and $V_\beta(\cE)=0$, and Proposition~\ref{prop:defect} gives the identity; $\delta(\cE)\ge0$ by Bogomolov's inequality and $\alpha(\cE)^2\le0$ by the Hodge index theorem. A table of $\pi_*\cE$ on the ray of $\gamma^{z}$ is exactly one whose columns $a_1,\dots,a_n$ vanish. For an arithmetic progression, $\sum_j(a_j-c)^2=e^2\sum_j(j-\frac{n+1}2)^2=e^2\,\frac{n(n^2-1)}{12}$.
\end{proof}

Let $a_j=a_1+(j-1)e$. The bundle $\cU=\cE((a_n+e)H)$ has no cohomology at the twists $-eH,\dots,-neH$, so it satisfies the cohomological condition for Ulrich bundles on the polarized variety $(X,eH)$. When $eH$ is very ample, the semistability of $\cE$ also follows from \cite[Thm.~2.9]{CH12}, and the inequality $e^2d\ge\sigma$, i.e.\ $(n+1)(eH)^n\ge(K_X^2-2c_2(X))\cdot(eH)^{n-2}$, from the inequality of Lopez for Ulrich bundles \cite[Lem.~3.1]{Lo26} applied to the polarization $eH$, whose proof is the computation of Proposition~\ref{prop:defect} in this case; see also \cite{An1} for surfaces. What Corollary~\ref{cor:ES} adds is the case of root sequences not in arithmetic progression, which occurs only for $n\ge3$ (two roots always form one) and does not reduce by a twist to the Ulrich condition for a multiple of $H$, and the interpretation of the inequality as the condition for a supernatural ray to be realized. For $e=1$ it gives $d\ge\sigma$, which is the Ulrich case: then $\pi_*\cE\cong\OO(-a_n-1)^{\oplus rd}$, and $\cE((a_n+1)H)$ is an Ulrich bundle.

\section{Sheaves of lower dimension}\label{sec:lowerdim}

In this section $1\le d\le n$, and $d$ denotes a dimension, not the degree $H^n$ of Section~\ref{sec:polarized}. For a coherent sheaf $\cF$ of dimension at most $d$ on $\PP^n$ write
\[
\chi(\cF(t))=\sum_{i=0}^d(-1)^i\,e_i(\cF)\binom{t+d-i}{d-i},
\]
and put $\rk_d\cF=e_0(\cF)$ and $\deg_d\cF=-e_1(\cF)$. Both are additive on exact sequences. The number $\rk_d\cF\ge0$ is the multiplicity of the $d$-dimensional part of $\cF$; if it vanishes, $\deg_d\cF\ge0$ is the multiplicity of the $(d-1)$-dimensional part. Since $\binom{t+n}n=\frac{t^n}{n!}+\frac{n+1}2\,\frac{t^{n-1}}{(n-1)!}+\dots$, Riemann--Roch (Section~\ref{sec:tables}) gives $\rk_n\cF=\rk\cF$ and $\deg_n\cF=\deg\cF$. In general we use the normalized slope $\mu_d(\cF)=\deg_d\cF/\rk_d\cF=\hat\mu(\cF)-\frac{d+1}2$, where $\hat\mu$ is the slope of \cite[\S1.6]{HL10}. Let $P_\cF$ be the upper envelope of the points $(\rk_d\cF',\deg_d\cF')$ for coherent subsheaves $\cF'\subset\cF$; for $\cF$ pure of dimension $d$ it is the Harder--Narasimhan polygon of $\cF$ with respect to $\mu_d$. Equivalently, if $P^{\hat\mu}_{\cF}$ denotes the polygon with slopes $\hat\mu$ and horizontal lengths $e_0$, then $P_\cF(m)=P^{\hat\mu}_\cF(m)-\frac{d+1}2m$ \cite[\S1.6]{HL10}.

The table $\gamma(\cF)$ has rows $0,\dots,d$. We use root sequences $z=(z_1>\dots>z_d)$ of length $d$, with $c(z)=\frac1d\sum_jz_j$ and $\mu(z)=-c(z)-\frac{d+1}2$, the functionals $\phi_z$ given by \eqref{eq:phi} with $n$ replaced by $d$, and, as in Section~\ref{sec:coherent},
\begin{align*}
\Lambda_\cF(\lambda)&=\inf\big\{\phi_z(\gamma(\cF)):\ \mu(z)=\lambda\big\},\qquad \lambda\in\tfrac1d\ZZ,\\
\Psi_\cF(m)&=\inf_{\lambda\in\frac1d\ZZ}\big(m\lambda+\Lambda_\cF(\lambda)\big),\qquad 0\le m\le\rk_d\cF .
\end{align*}
For $d=n$ these are the functions of Section~\ref{sec:coherent}.

\begin{lemma}[projection]\label{lem:projection}
Let $\cF$ be a coherent sheaf of dimension at most $d$ on $\PP^n$, let $\Pi\subset\PP^n$ be a linear subspace of dimension $n-d-1$ disjoint from $\operatorname{Supp}\cF$ (empty if $d=n$), and let $p\colon\PP^n\setminus\Pi\to\PP^d$ be the projection from $\Pi$. On coherent sheaves $\mathcal G$ with $\operatorname{Supp}\mathcal G\subset\operatorname{Supp}\cF$ the functor $p_*$ is exact, and:
\begin{enumerate}[label=\textup{(\roman*)}]
\item $H^j(\PP^d,(p_*\mathcal G)(t))\cong H^j(\PP^n,\mathcal G(t))$ for all $j$ and $t$, functorially in $\mathcal G$; in particular $\gamma(p_*\mathcal G)=\gamma(\mathcal G)$, $\rk p_*\mathcal G=\rk_d\mathcal G$ and $\deg p_*\mathcal G=\deg_d\mathcal G$;
\item if $\mathcal G$ is pure of dimension $d$, then $p_*\mathcal G$ is torsion-free;
\item if $\mathcal G_x$ is a Cohen--Macaulay module of dimension $d$ for every closed point $x\in\operatorname{Supp}\mathcal G$, then $p_*\mathcal G$ is locally free.
\end{enumerate}
\end{lemma}

\begin{proof}
Such a $\mathcal G$ is a module over a closed subscheme $Z\subset\PP^n$ with support $\operatorname{Supp}\cF$, for instance the one defined by a power of the ideal of $\operatorname{Supp}\cF$, and $Z\subset\PP^n\setminus\Pi$. For $y\in\PP^d$ the linear space $L_y$ spanned by $\Pi$ and $y$ contains $\Pi$ as a hyperplane, and $p^{-1}(y)\cap Z=Z\cap L_y$ is projective and contained in the affine space $L_y\setminus\Pi$, hence finite. So $p|_Z$ is proper and quasi-finite, hence finite; $p_*\mathcal G=(p|_Z)_*\mathcal G$ is coherent, and $p_*$ is exact, without higher direct images. Since $p^*\OO(1)=\OO(1)$, the projection formula gives $p_*(\mathcal G(t))=(p_*\mathcal G)(t)$, and (i) follows; the rank and the degree are read off from $\chi((p_*\mathcal G)(t))=\chi(\mathcal G(t))$ by Riemann--Roch on $\PP^d$. (ii) Since $p|_Z$ is finite, the associated points of $p_*\mathcal G$ are the images of those of $\mathcal G$, which are generic points of $d$-dimensional subvarieties; so the generic point of $\PP^d$ is the only associated point of $p_*\mathcal G$. (iii) For a closed point $y\in\PP^d$, the stalk $(p_*\mathcal G)_y$ is a finite module over the regular local ring $\OO_{\PP^d,y}$ of dimension $d$, and its depth is the minimum of the depths of the $\mathcal G_x$, $x\in p^{-1}(y)\cap Z$, that is $d$ (or the stalk is zero). By the Auslander--Buchsbaum formula it is free.
\end{proof}

\Needspace{6\baselineskip}
\begin{theorem}\label{thm:lowerdim}
Let $\cF$ be a coherent sheaf of dimension at most $d$ on $\PP^n$.
\begin{enumerate}[label=\textup{(\alph*)}]
\item For every coherent subsheaf $\cF'\subset\cF$ and every root sequence $z$ of length $d$,
\[
\deg_d\cF'-\mu(z)\rk_d\cF'\le\phi_z(\gamma(\cF)).
\]
Hence $\deg_d\cF'\le\Psi_\cF(\rk_d\cF')$, i.e.\ $P_\cF\le\Psi_\cF$. For $\cF$ pure of dimension $d$, the Harder--Narasimhan polygon of $\cF$ with respect to $\mu_d=-e_1/e_0$ lies below $\Psi_\cF$.
\item If $\cF_x$ is a Cohen--Macaulay module of dimension $d$ at every closed point $x$ of its support, for instance if $\cF$ is a vector bundle on a Cohen--Macaulay subvariety of dimension $d$, then $\gamma(\cF)$ is the table of a vector bundle on $\PP^d$, and $\Psi_\cF$ is the Boij--Söderberg polygon of its chain decomposition, with finitely many edges and endpoints $(0,0)$ and $(\rk_d\cF,\deg_d\cF)$.
\item $\Psi_\cF\le\Psi_{\cF'\oplus\cF''}$ for every exact sequence $0\to\cF'\to\cF\to\cF''\to0$. For $d=2$, Theorem~D holds for $\cF$ in the form of Proposition~\ref{prop:cohD}, with $\rk_2$ and $\deg_2$ in place of the rank and the degree.
\end{enumerate}
\end{theorem}

\begin{proof}
Choose $\Pi$ and $p$ as in Lemma~\ref{lem:projection}; a general $\Pi$ of dimension $n-d-1$ is disjoint from $\operatorname{Supp}\cF$. Every subsheaf and every quotient of $\cF$ is supported in $\operatorname{Supp}\cF$, so $p_*$ carries $0\to\cF'\to\cF\to\cF''\to0$ to an exact sequence on $\PP^d$, and by Lemma~\ref{lem:projection}(i) it preserves the tables, the long exact cohomology sequences, and the numbers $\rk_d$ and $\deg_d$, which become the rank and the degree on $\PP^d$. In particular $\Lambda_\cF$ and $\Psi_\cF$ are the functions of Section~\ref{sec:coherent} for the coherent sheaf $p_*\cF$ on $\PP^d$. (a) is Proposition~\ref{prop:cohA}(a) for $p_*\cF'\subset p_*\cF$. (b) By Lemma~\ref{lem:projection}(iii), $p_*\cF$ is a vector bundle on $\PP^d$, and by Proposition~\ref{prop:minformula} and Lemma~\ref{lem:legendre}(i) its function $\Psi$ is its Boij--Söderberg polygon; a vector bundle on a Cohen--Macaulay subvariety $X$ of dimension $d$ has stalks free over the Cohen--Macaulay local rings $\OO_{X,x}$ of dimension $d$. (c) is Proposition~\ref{prop:cohC}(a) and Proposition~\ref{prop:cohD} for $p_*$ of the sequence.
\end{proof}

\begin{remark}\label{rem:lowerdim}
\begin{enumerate}[label=(\arabic*)]
\item \emph{Vector bundles on subvarieties.} Let $X\subset\PP^n$ be a smooth subvariety of dimension $d$, $H=\OO_X(1)$, and $\cE$ a vector bundle on $X$, viewed as a sheaf on $\PP^n$. The subsheaves of $\cE$ on $\PP^n$ are the subsheaves on $X$, and by Riemann--Roch on $X$ a torsion-free subsheaf $\cE'$ of rank $r'$ has $\rk_d\cE'=r'H^d$ and $\deg_d\cE'=c_1(\cE')\cdot H^{d-1}-\frac{r'}2(K_X+(d+1)H)\cdot H^{d-1}$. So $P_\cE$ is the $\mu_H$-Harder--Narasimhan polygon of $\cE$ after a linear change of coordinates, and Theorem~\ref{thm:lowerdim} bounds it by the Boij--Söderberg polygon of the vector bundle $p_*\cE$ on $\PP^d$, whose table is $(h^j(\cE(tH)))_{j,t}$. If $H$ is ample and base-point-free, the proof applies to a finite morphism $\pi\colon X\to\PP^d$ with $\pi^*\OO(1)=\OO_X(H)$ in place of $p$; this is the inequality $P_\beta\le\Psi_{\pi_*\cE}$ in the proof of Corollary~\ref{cor:cases}(a), where the dimension is denoted $n$.
\item \emph{Hilbert coefficients.} Let $M$ be a finitely generated graded module of Krull dimension $d+1\ge2$ over $S=K[x_0,\dots,x_n]$, with Hilbert polynomial $\sum_{i=0}^d(-1)^ie_i(M)\binom{t+d-i}{d-i}$, and let $N\subset M$ be a graded submodule of the same dimension. Then $\widetilde N\subset\widetilde M$, the Hilbert polynomials of $N$ and $\widetilde N$ agree, and Theorem~\ref{thm:lowerdim}(a) gives
\[
e_1(N)\ \ge\ -\Psi_{\widetilde M}\big(e_0(N)\big).
\]
The table of $\widetilde M$, hence the right-hand side, is determined by the Hilbert functions of $M$ and of its local cohomology modules $H^j_{\mathfrak m}(M)$, since $h^0(\widetilde M(t))=\dim M_t-\dim H^0_{\mathfrak m}(M)_t+\dim H^1_{\mathfrak m}(M)_t$ and $h^j(\widetilde M(t))=\dim H^{j+1}_{\mathfrak m}(M)_t$ for $j\ge1$. If $M$ is Cohen--Macaulay, then $M$ is free over a polynomial ring $K[y_0,\dots,y_d]$ generated by linear forms, with $h_k$ basis elements in degree $k$, where $\sum_kh_kt^k/(1-t)^{d+1}$ is the Hilbert series of $M$; so $p_*\widetilde M\cong\bigoplus_k\OO(-k)^{h_k}$ for the projection $p$ from $V(y_0,\dots,y_d)$, and the bound says that $e_1(N)$ is at least the sum of the $e_0(N)$ smallest elements of the multiset in which $k$ occurs $h_k$ times. In this case it reduces to the elementary bound for subsheaves of a split bundle.
\item \emph{The Betti side.} The multiplicity conjecture of Herzog, Huneke and Srinivasan \cite{HS98}, proved by Eisenbud and Schreyer \cite{ES09} and, for the upper bound in the non-Cohen--Macaulay case, by Boij and Söderberg \cite{BS12}, bounds $e_0(M)$ through the pure decomposition of the Betti table of $M$. Theorem~\ref{thm:lowerdim} is of a different kind: it bounds $e_0$ and $e_1$ of all submodules through the cohomology table. The naive analogue on the Betti side, with the number of generators in place of the rank, fails. For $n\ge2$ and $M=S/(x_0)\oplus S/(x_1)$, the Betti table is twice the pure diagram of type $(0,1)$, whose pieces have multiplicity $1$ per generator, while the cyclic submodule $N=S\cdot(1,1)\cong S/(x_0x_1)$ has $e_0(N)=2$. Theorem~\ref{thm:lowerdim} holds in this example, with $d=n-1$: $e_0(N)=e_0(M)=2$ and $e_1(N)=1>0=e_1(M)$.
\end{enumerate}
\end{remark}

\appendix
\section{Supernatural bundles on \texorpdfstring{$\PP^3$}{P3} with small root windows}\label{app:P3}

By Corollary~\ref{cor:supernatural}, every vector bundle on $\PP^3$ whose table is a multiple of a supernatural table $\gamma^z$, $z=(z_1>z_2>z_3)$, is semistable. Two classes of such bundles are familiar. If $z_1-z_3\le3$, they are direct sums of copies of a twist of $\OO$, $\Omega^1$ or $\Omega^2$ (see below). If the roots form an arithmetic progression, a twist is an Ulrich bundle for a multiple of $\OO(1)$ (Remark~\ref{rem:main}(3)). The smallest root sequences outside both classes have $z_1-z_3=4$ and gaps $(z_1-z_2,z_2-z_3)=(3,1)$ or $(1,3)$. In this appendix we describe the bundles with $z_1-z_3=4$, up to a twist. For the gaps $(3,1)$ they are the duals of certain Steiner bundles (Lemma~\ref{lem:beil31}). For the gaps $(2,2)$ they are instanton bundles (Lemma~\ref{lem:beil22} and Remark~\ref{rem:beil22}), and a comparison with the semistability theorem of Jardim \cite{Ja06} shows that the vanishing at all three roots is needed. The consequences for stability are collected in Corollary~\ref{cor:appP3}.

The tool is Beilinson's theorem \cite{Be78}: a coherent sheaf $\cF$ on $\PP^n$ is the abutment of a spectral sequence with $E_1^{p,j}=H^j(\cF(p))\otimes\Omega^{-p}(-p)$, $-n\le p\le0$. If $z_1-z_3\le3$, the three roots lie in a window $\{k,k-1,k-2,k-3\}$, and after twisting by $\OO(k)$, which moves the roots into $\{0,-1,-2,-3\}$, only one column of the spectral sequence survives; so a vector bundle with table $q\gamma^z$ is a direct sum of copies of one of $\OO$, $\Omega^1(1)$, $\Omega^2(2)\cong T(-2)$, $\OO(-1)$, suitably twisted. If $z_1-z_3=4$, the gaps are $(2,2)$, $(3,1)$ or $(1,3)$, and the last two cases are dual to each other. We write $U=H^0(\OO_{\PP^3}(1))$, so that $\PP^3=\PP(U^\vee)$, $H^0(T(-1))=U^\vee$ and $\operatorname{Hom}(T(-2),\OO)=H^0(\Omega^1(2))=\Lambda^2U$, and we use that $T(-3)$ and $T(-5)$ are acyclic; here $T=T_{\PP^3}$.

\begin{lemma}[gaps $(2,2)$]\label{lem:beil22}
Let $z=(1,-1,-3)$ and $q\in\frac13\ZZ_{>0}$. A vector bundle $F$ on $\PP^3$ has cohomology table $q\gamma^z$ if and only if $F\cong\ker\varphi$ for a surjection $\varphi\colon V\otimes T(-2)\to W\otimes\OO$ with $\dim V=\dim W=3q$ such that
\[
B=H^0(\varphi(1))\colon V\otimes U^\vee\longrightarrow W\otimes U
\]
is an isomorphism. The surjectivity of $\varphi$ follows from that of $B$. For $q=\frac13$ these are the twists $N(-1)$ of the null correlation bundles.
\end{lemma}

\begin{proof}
If $\gamma(F)=q\gamma^z$, then $F(-1)$ and $F(-3)$ are acyclic, so the columns $p=-1,-3$ of the Beilinson spectral sequence vanish; the column $p=0$ is $H^1(F)\otimes\OO$ in the row $1$, and the column $p=-2$ is $H^2(F(-2))\otimes\Omega^2(2)$ in the row $2$, with $\Omega^2(2)\cong T(-2)$ on $\PP^3$. The only differential is $d_2\colon H^2(F(-2))\otimes T(-2)\to H^1(F)\otimes\OO$, of total degrees $0\to1$, and convergence to $F$ forces it to be surjective with kernel $F$. Here $h^2(F(-2))=h^1(F)=3q$. Twisting $0\to F(1)\to V\otimes T(-1)\to W\otimes\OO(1)\to0$ and using $H^i(T(-1))=H^i(\OO(1))=0$ for $i>0$, the acyclicity of $F(1)$ says exactly that $B$ is bijective. If $B$ is surjective, the image $I$ of $\varphi(1)$ has $H^0(I)=H^0(W\otimes\OO(1))$, and $W\otimes\OO(1)$ is globally generated, so $\varphi$ is surjective. Conversely, let $G=\ker\varphi$. It is locally free, and $G(t)$ is acyclic for $t=1$ by the hypothesis, and for $t=-1,-3$ because $T(-3)$, $T(-5)$, $\OO(-1)$, $\OO(-3)$ are. Every piece of the chain of $G$ has zero columns at $1,-1,-3$, hence root sequence $z$, so $\gamma(G)=q'\gamma^z$ with $6q'=\rk G=6q$. For $\dim V=1$, $\varphi$ is a nonzero $2$-form $\omega\in\Lambda^2U$, and $B$ is $x\mapsto\omega(\cdot,x)$, bijective iff $\omega$ is nondegenerate; then $\ker\varphi=N(-1)$ \cite[I, \S4.2]{OSS}.
\end{proof}

\begin{remark}\label{rem:beil22}
The bundle $E=F(1)$ is the cohomology of the linear monad $V\otimes\OO(-1)\to V\otimes U^\vee\otimes\OO\to W\otimes\OO(1)$, so it is an instanton bundle of rank $6q$ and charge $3q$ in the sense of \cite{Ja06}, with $E$, $E(-2)$ and $E(-4)$ acyclic; and $F(3)$ is an Ulrich bundle on $(\PP^3,\OO(2))$. For the correspondence between Ulrich bundles on $(\PP^3,\OO(d))$ and instanton bundles see \cite[Thm.~3.6]{CMR21}. Jardim proves that locally free instanton sheaves of rank at most $2n-1$ on $\PP^n$ are semistable, and that this bound is sharp \cite{Ja06}. A simple example on $\PP^3$ is $T(-1)\oplus\Omega^1(1)$, the cohomology of $\OO(-1)\to U^\vee\otimes\OO\oplus U\otimes\OO\to\OO(1)$: an instanton bundle of rank $6$ and charge $1$, destabilized by $T(-1)$, with $h^0=4$ in the column where our tables have a root. So the supernatural hypothesis in Corollary~\ref{cor:supernatural}(a) cannot be weakened to the vanishing conditions that define instanton bundles.
\end{remark}

\begin{lemma}[gaps $(3,1)$]\label{lem:beil31}
Let $q\in\frac12\ZZ_{>0}$. A vector bundle $F$ on $\PP^3$ has cohomology table $q\gamma^{(1,-2,-3)}$ if and only if $F\cong\ker\theta$ for a surjection $\theta\colon V\otimes\Omega^1(1)\to W\otimes\OO$ with $\dim V=4q$, $\dim W=6q$, such that $H^0(\theta(1))\colon V\otimes\Lambda^2U\to W\otimes U$ is bijective. Every map $\Omega^1(1)\to\OO$ is the restriction of a linear form on $U$, so $\theta$ is the restriction of a linear map $\tilde\theta\colon V\otimes U\to W$; put $S=\ker\tilde\theta$. Then $\dim S=10q$,
\[
0\longrightarrow F\longrightarrow S\otimes\OO\longrightarrow V\otimes\OO(1)\longrightarrow0
\]
is exact, the second map being the evaluation, and the condition on $\theta$ says that $S\otimes U\to V\otimes S^2U$ is bijective. So $F^\vee$ is the Steiner bundle $\coker\big(V^\vee\otimes\OO(-1)\to S^\vee\otimes\OO\big)$, of rank $6q$; Steiner bundles of rank $n$ on $\PP^n$ were introduced in \cite{DK93}. For $q=\frac12$ an example is the Hankel space $S=\langle e_1\otimes u_{j}+e_2\otimes u_{j-1}:0\le j\le4\rangle$, for a basis $e_1,e_2$ of $V$ and a basis $u_0,\dots,u_3$ of $U$, with $u_{-1}=u_4=0$.
\end{lemma}

\begin{proof}
If $\gamma(F)=q\gamma^z$, $z=(1,-2,-3)$, then $F(-2)$ and $F(-3)$ are acyclic, and the cohomology of $F$ and $F(-1)$ is concentrated in degree one, with $h^1(F)=6q$ and $h^1(F(-1))=4q$. So the only terms of the Beilinson spectral sequence are $E_1^{-1,1}=H^1(F(-1))\otimes\Omega^1(1)$ and $E_1^{0,1}=H^1(F)\otimes\OO$, and as in Lemma~\ref{lem:beil22} the differential is a surjection $\theta$ with kernel $F$. The sequence $0\to F(1)\to V\otimes\Omega^1(2)\to W\otimes\OO(1)\to0$ and the vanishing of $H^i(\Omega^1(2))$ and $H^i(\OO(1))$ for $i>0$ show that $F(1)$ is acyclic if and only if $H^0(\theta(1))$ is bijective. The assertion on maps $\Omega^1(1)\to\OO$ follows from the Euler sequence $0\to\Omega^1(1)\to U\otimes\OO\to\OO(1)\to0$, since $\operatorname{Hom}(\OO(1),\OO)=\Ext^1(\OO(1),\OO)=0$. Since $\theta$ is surjective, so is $\tilde\theta$, hence $\dim S=10q$, and the map $V\otimes U\otimes\OO\to W\otimes\OO\oplus V\otimes\OO(1)$ given by $\tilde\theta$ and the evaluation is surjective, by the snake lemma applied to the Euler sequence tensored with $V$; restricting it to $S\otimes\OO$ gives the displayed sequence. Its cohomology at the twist $1$ is $0\to H^0(F(1))\to S\otimes U\to V\otimes S^2U\to H^1(F(1))\to0$, which gives the second form of the condition. Conversely, if $\theta$ is a surjection with $H^0(\theta(1))$ bijective, then $F(1)$ is acyclic, and $F(-2)$, $F(-3)$ are acyclic because $\OO(-1),\OO(-2),\OO(-3)$ are; as in Lemma~\ref{lem:beil22}, every piece of the chain of $F$ has the roots $1,-2,-3$, so $\gamma(F)=q'\gamma^z$ with $6q'=\rk F=6q$. Finally, the condition on $S$ alone suffices: if $S\otimes U\to V\otimes S^2U$ is surjective, then, as in Lemma~\ref{lem:beil22}, the image of $S\otimes\OO(1)\to V\otimes\OO(2)$ contains all sections of the globally generated sheaf $V\otimes\OO(2)$, so the evaluation $S\otimes\OO\to V\otimes\OO(1)$ is surjective; with $W=(V\otimes U)/S$, the Euler sequence identifies the cokernel and the kernel of the induced map $\theta\colon V\otimes\Omega^1(1)\to W\otimes\OO$ with those of the evaluation on $S\otimes\OO$, so $\theta$ is surjective with kernel $F$. For the Hankel space, the $2\times5$ matrix of linear forms of the evaluation map has rows $(u_0,u_1,u_2,u_3,0)$ and $(0,u_0,u_1,u_2,u_3)$, which are proportional at no point of $\PP^3$, so the evaluation is surjective; and $S\otimes U\to V\otimes S^2U$, a map between spaces of dimension $20$, is injective over every field. Indeed, with $s_j=e_1\otimes u_j+e_2\otimes u_{j-1}$, an element $\sum_{j=0}^4s_j\otimes\ell_j$, $\ell_j=\sum_ka_{jk}u_k$, maps to $e_1\otimes\sum_{j=0}^3u_j\ell_j+e_2\otimes\sum_{j=0}^3u_j\ell_{j+1}$. The kernel of the multiplication $U\otimes U\to S^2U$ consists of the alternating tensors, so if the image is zero, the matrices $(a_{jk})_{0\le j,k\le3}$ and $(a_{j+1,k})_{0\le j,k\le3}$ are alternating. Then $a_{j+1,k}=-a_{k+1,j}=a_{j,k+1}$ for $0\le j\le3$, $0\le k\le2$, so $a_{jk}=h(j+k)$ depends only on $j+k\in[0,7]$, and the vanishing of the diagonals gives $h(2j)=a_{jj}=0$ and $h(2j+1)=a_{j+1,j}=0$ for $0\le j\le3$. So all $\ell_j$ vanish.
\end{proof}

\begin{corollary}\label{cor:appP3}
\begin{enumerate}[label=\textup{(\alph*)}]
\item For every $q\in\frac13\ZZ_{>0}$, the bundles $F$ of Lemma~\ref{lem:beil22} are semistable. So the instanton bundles $E=F(1)$ of Remark~\ref{rem:beil22}, of rank $6q$ and charge $3q$, are semistable also when their rank exceeds the bound $2n-1=5$ of \cite{Ja06}.
\item For every $q\in\frac12\ZZ_{>0}$, the bundles $F$ of Lemma~\ref{lem:beil31} and their duals, the Steiner bundles $\coker\big(V^\vee\otimes\OO(-1)\to S^\vee\otimes\OO\big)$ of rank $6q$ for which $S\otimes U\to V\otimes S^2U$ is bijective, are semistable. For $q=\frac12$ they are stable, and in characteristic zero the general one is stable for every $q$.
\end{enumerate}
\end{corollary}

\begin{proof}
The semistability is Corollary~\ref{cor:supernatural}(a), in any characteristic; in (a) it also follows from \cite[Thm.~2.9]{CH12}, since $F(3)$ is an Ulrich bundle on $(\PP^3,\OO(2))$. In (b), for $q=\frac12$, $\rk F=3$ and $c_1(F)=-2$ are coprime, so every such $F$ is stable. The subspaces $S\subset V\otimes U$ of dimension $10q$ for which $S\otimes U\to V\otimes S^2U$ is bijective form a nonempty open subset of the Grassmannian; it contains the direct sums of $2q$ Hankel spaces. A choice of basis of $S$ identifies the frame bundle over this open set with a nonempty open subset of the space of $10q\times4q$ matrices of linear forms, so a general point of it is a general matrix of linear forms in the sense of Coskun, Huizenga and Smith, whose theorem \cite[Thm.~5.4]{CHS} (in characteristic zero) applies to the Steiner bundle $F^\vee$ of rank $r=6q\ge3$ with $t=4q$ relations, since $r<(1+\sqrt3)\,t$. So the general vector bundle with table $q\gamma^{(1,-2,-3)}$ is stable, and dually for the gaps $(1,3)$. For $q\ge1$, direct sums give strictly semistable ones.
\end{proof}

\section*{Acknowledgments}
This paper grew out of an attempt to answer the questions on the Boij--Söderberg cone of cohomology tables of vector bundles on the surfaces of \cite{An1} that David Eisenbud and Frank-Olaf Schreyer asked the first author after that paper appeared (see Section~\ref{sec:twists}). We are grateful for their interest and encouragement.

\section*{AI Declaration}
ChatGPT and Claude were used for exploratory computations, literature searches, and interactive feedback on earlier drafts of the manuscript. The authors take full responsibility for all content.

\begin{thebibliography}{EFW11}
\bibitem[An1]{An1} C.~Anghel, \emph{Smooth counterexamples to the Eisenbud--Schreyer--Weyman Ulrich existence problem}, arXiv:2609.02718.
\bibitem[An2]{An2} C.~Anghel, \emph{Polarized varieties without arithmetically Cohen--Macaulay bundles and section rings without graded maximal Cohen--Macaulay modules in characteristic zero}, arXiv:2609.15589.
\bibitem[An3]{An3} C.~Anghel, \emph{The recent counterexamples to the Eisenbud--Schreyer--Weyman Ulrich problem (2003) and to Hochster's small Cohen--Macaulay conjecture (1973): a survey, with a generalization to arbitrary dimension}, arXiv:2610.30027.
\bibitem[AB13]{AB13} D.~Arcara, A.~Bertram, \emph{Bridgeland-stable moduli spaces for $K$-trivial surfaces}, J.\ Eur.\ Math.\ Soc.\ \textbf{15} (2013), 1--38.
\bibitem[ABCH13]{ABCH13} D.~Arcara, A.~Bertram, I.~Coskun, J.~Huizenga, \emph{The minimal model program for the Hilbert scheme of points on $\PP^2$ and Bridgeland stability}, Adv.\ Math.\ \textbf{235} (2013), 580--626.
\bibitem[Bay09]{Bay09} A.~Bayer, \emph{Polynomial Bridgeland stability conditions and the large volume limit}, Geom.\ Topol.\ \textbf{13} (2009), 2389--2425.
\bibitem[Be78]{Be78} A.~A.~Beilinson, \emph{Coherent sheaves on $\mathbb P^n$ and problems of linear algebra}, Funct.\ Anal.\ Appl.\ \textbf{12} (1978), 214--216.
\bibitem[Bo78]{Bo78} F.~A.~Bogomolov, \emph{Holomorphic tensors and vector bundles on projective varieties}, Izv.\ Akad.\ Nauk SSSR Ser.\ Mat.\ \textbf{42} (1978), 1227--1287.
\bibitem[BS12]{BS12} M.~Boij, J.~Söderberg, \emph{Betti numbers of graded modules and the multiplicity conjecture in the non-Cohen--Macaulay case}, Algebra Number Theory \textbf{6} (2012), 437--454.
\bibitem[Br07]{Br07} T.~Bridgeland, \emph{Stability conditions on triangulated categories}, Ann.\ of Math.\ (2) \textbf{166} (2007), 317--345.
\bibitem[Br08]{Br08} T.~Bridgeland, \emph{Stability conditions on K3 surfaces}, Duke Math.\ J.\ \textbf{141} (2008), 241--291.
\bibitem[CH12]{CH12} M.~Casanellas, R.~Hartshorne, \emph{Stable Ulrich bundles}, with an appendix by F.~Geiss and F.-O.~Schreyer, Internat.\ J.\ Math.\ \textbf{23} (2012), 1250083.
\bibitem[CHS25]{CHS} I.~Coskun, J.~Huizenga, G.~Smith, \emph{Stability and cohomology of kernel bundles on $\mathbb P^n$}, Michigan Math.\ J.\ \textbf{75} (2025), 173--198.
\bibitem[CMR21]{CMR21} L.~Costa, R.~M.~Miró-Roig, \emph{Instanton bundles vs Ulrich bundles on projective spaces}, Beitr.\ Algebra Geom.\ \textbf{62} (2021), 429--439.
\bibitem[DK93]{DK93} I.~Dolgachev, M.~Kapranov, \emph{Arrangements of hyperplanes and vector bundles on $\mathbf P^n$}, Duke Math.\ J.\ \textbf{71} (1993), 633--664.
\bibitem[EFW11]{EFW11} D.~Eisenbud, G.~Fløystad, J.~Weyman, \emph{The existence of equivariant pure free resolutions}, Ann.\ Inst.\ Fourier \textbf{61} (2011), 905--926.
\bibitem[ES09]{ES09} D.~Eisenbud, F.-O.~Schreyer, \emph{Betti numbers of graded modules and cohomology of vector bundles}, J.\ Amer.\ Math.\ Soc.\ \textbf{22} (2009), 859--888.
\bibitem[ES10]{ES10} D.~Eisenbud, F.-O.~Schreyer, \emph{Cohomology of coherent sheaves and series of supernatural bundles}, J.\ Eur.\ Math.\ Soc.\ \textbf{12} (2010), 703--722.
\bibitem[ESW03]{ESW03} D.~Eisenbud, F.-O.~Schreyer, \emph{Resultants and Chow forms via exterior syzygies}, with an appendix by J.~Weyman, J.\ Amer.\ Math.\ Soc.\ \textbf{16} (2003), 537--579.
\bibitem[ErSa16]{ErSa16} D.~Erman, S.~V.~Sam, \emph{Supernatural analogues of Beilinson monads}, Compositio Math.\ \textbf{152} (2016), 2545--2562.
\bibitem[ErSa17]{ErS} D.~Erman, S.~V.~Sam, \emph{Questions about Boij--Söderberg theory}, Proc.\ Sympos.\ Pure Math.~\textbf{95}, Amer.\ Math.\ Soc., 2017, 285--304; arXiv:1606.01867.
\bibitem[Fl12]{Fl12} G.~Fløystad, \emph{Boij--Söderberg theory: introduction and survey}, in: Progress in Commutative Algebra~1, de Gruyter, Berlin, 2012, 1--54.
\bibitem[GL26]{GL26} Ö.~Genç, A.~F.~Lopez, \emph{On the existence of polarized varieties without weakly Ulrich bundles, ordinary instantons with bounded quantum number and supernatural bundles with bounded roots}, arXiv:2610.36231v2.
\bibitem[HLP34]{HLP} G.~H.~Hardy, J.~E.~Littlewood, G.~P\'olya, \emph{Inequalities}, Cambridge University Press, Cambridge, 1934.
\bibitem[HK84]{HK84} J.~Herzog, M.~Kühl, \emph{On the Betti numbers of finite pure and linear resolutions}, Comm.\ Algebra \textbf{12} (1984), 1627--1646.
\bibitem[HS98]{HS98} J.~Herzog, H.~Srinivasan, \emph{Bounds for multiplicities}, Trans.\ Amer.\ Math.\ Soc.\ \textbf{350} (1998), 2879--2902.
\bibitem[Ho64]{Ho64} G.~Horrocks, \emph{Vector bundles on the punctured spectrum of a local ring}, Proc.\ London Math.\ Soc.\ (3) \textbf{14} (1964), 689--713.
\bibitem[HL10]{HL10} D.~Huybrechts, M.~Lehn, \emph{The geometry of moduli spaces of sheaves}, 2nd ed., Cambridge Math.\ Library, Cambridge Univ.\ Press, 2010.
\bibitem[IMW23]{IMW23} S.~B.~Iyengar, L.~Ma, M.~E.~Walker, \emph{Lim Ulrich sequences and Boij--Söderberg cones}, Forum Math.\ Sigma \textbf{11} (2023), e115.
\bibitem[Ja06]{Ja06} M.~Jardim, \emph{Instanton sheaves on complex projective spaces}, Collect.\ Math.\ \textbf{57} (2006), 69--91.
\bibitem[Ka32]{Ka32} J.~Karamata, \emph{Sur une inégalité relative aux fonctions convexes}, Publ.\ Math.\ Univ.\ Belgrade \textbf{1} (1932), 145--148.
\bibitem[La04]{La04} A.~Langer, \emph{Semistable sheaves in positive characteristic}, Ann.\ of Math.\ (2) \textbf{159} (2004), 251--276.
\bibitem[LQ14]{LQ14} J.~Lo, Z.~Qin, \emph{Mini-walls for Bridgeland stability conditions on the derived category of sheaves over surfaces}, Asian J.\ Math.\ \textbf{18} (2014), no.~2, 321--344.
\bibitem[Lo26]{Lo26} A.~F.~Lopez, \emph{On the existence of smooth varieties of any dimension carrying no Ulrich bundles for some very ample line bundle}, arXiv:2609.06678v3.
\bibitem[MS17]{MS17} E.~Macrì, B.~Schmidt, \emph{Lectures on Bridgeland stability}, in: Moduli of Curves, Lect.\ Notes Unione Mat.\ Ital.~21, Springer, Cham, 2017, 139--211; arXiv:1607.01262.
\bibitem[MiSp08]{MiSp08} E.~Miller, D.~E.~Speyer, \emph{A Kleiman--Bertini theorem for sheaf tensor products}, J.\ Algebraic Geom.\ \textbf{17} (2008), 335--340.
\bibitem[Mu70]{Mu70} D.~Mumford, \emph{Abelian varieties}, Tata Inst.\ Fund.\ Res.\ Stud.\ Math.~5, Oxford University Press, London, 1970.
\bibitem[OSS80]{OSS} C.~Okonek, M.~Schneider, H.~Spindler, \emph{Vector bundles on complex projective spaces}, Progress in Mathematics~3, Birkhäuser, Boston, 1980.
\bibitem[Pi26a]{Pi26a} J.~Pine, \emph{Existence of Ulrich bundles on polarized surfaces and weakly Ulrich bundles on polarized varieties in positive characteristic}, arXiv:2610.28456v3.
\bibitem[Pi26b]{Pi26b} J.~Pine, \emph{ACM, almost Ulrich, and Ulrich vector bundles arising from Frobenius}, arXiv:2610.28464v1.
\bibitem[Ra66]{Ra66} S.~Ramanan, \emph{Holomorphic vector bundles on homogeneous spaces}, Topology \textbf{5} (1966), 159--177.
\bibitem[SE10]{ESi} F.-O.~Schreyer, D.~Eisenbud, \emph{Betti numbers of syzygies and cohomology of coherent sheaves}, Proc.\ ICM Hyderabad 2010, Vol.~II, 586--602.
\bibitem[Sh77]{Sh77} S.~S.~Shatz, \emph{The decomposition and specialization of algebraic families of vector bundles}, Compositio Math.\ \textbf{35} (1977), 163--187.
\bibitem[Um78]{Um78} H.~Umemura, \emph{On a theorem of Ramanan}, Nagoya Math.\ J.\ \textbf{69} (1978), 131--138.
\end{thebibliography}

\end{document}
